% concatenated from mhd-constant-field.tex
\documentclass[11pt,reqno]{amsart}
\usepackage[margin=2.6cm]{geometry}
\usepackage{amsmath,amssymb,amsthm,mathrsfs}
\usepackage[T1]{fontenc}
\usepackage[utf8]{inputenc}
\usepackage{lmodern}
\usepackage{microtype}
\usepackage{xcolor}
\usepackage{longtable,array,booktabs}
\providecommand{\tightlist}{\setlength{\itemsep}{0pt}\setlength{\parskip}{0pt}}
\usepackage[unicode,colorlinks=true,linkcolor=blue!50!black,citecolor=blue!45!black,urlcolor=blue!60!black]{hyperref}

% One counter for every statement and one for every display, both
% restarting in each section.
\theoremstyle{plain}
\newtheorem{theorem}{Theorem}[section]
\newtheorem{proposition}[theorem]{Proposition}
\newtheorem{lemma}[theorem]{Lemma}
\newtheorem{corollary}[theorem]{Corollary}
\theoremstyle{definition}
\newtheorem{definition}[theorem]{Definition}
\theoremstyle{remark}
\newtheorem{remark}[theorem]{Remark}

\numberwithin{equation}{section}

% Named differential and tensor operators.
\DeclareMathOperator{\Div}{div}
\DeclareMathOperator{\grad}{grad}
\DeclareMathOperator{\curl}{curl}
\DeclareMathOperator{\Curl}{Curl}
\DeclareMathOperator{\tr}{tr}
\DeclareMathOperator{\supp}{supp}

% Recurring regularity indices.
\newcommand{\breg}{\mathfrak b}
\newcommand{\sreg}{\mathfrak s}
\newcommand{\preg}{\mathfrak p}
\newcommand{\rpar}{\mathfrak r}

% Recurring fields along X, label-space tensors, and geometric coefficients.
\newcommand{\Ufield}{\mathsf U}
\newcommand{\Kmat}{\mathsf K}
\newcommand{\Mmat}{\mathsf M}
\newcommand{\Ccoef}{\mathcal G}   % product of two inverse deformation factors
\newcommand{\Hess}{\mathsf H}     % label Jacobian dF of a vector field along X
\newcommand{\eframe}{\mathsf e}   % Cartesian frame along X
\newcommand{\Vfield}{\mathsf V}
\newcommand{\vLag}{\mathsf v}
\newcommand{\BLag}{\mathsf B}
\newcommand{\bLag}{\mathsf b}
\newcommand{\gmetric}{\mathbf g}
\newcommand{\Bzero}{\mathbf B_0}
\newcommand{\Binf}{\mathbf B_\infty}

% Free propagators and standard domains.
\newcommand{\Cprop}{\mathsf C}
\newcommand{\Sprop}{\mathsf S}
\newcommand{\Real}{\mathbb R}
\newcommand{\Rn}{\Real^n}
\newcommand{\Rst}{\Real^{1+n}}
\newcommand{\Riesz}{\mathrm R}


%table of content
\makeatletter
\def\l@section{\@tocline{1}{6pt plus1pt}{1.2pc}{2.2pc}{}}
\def\l@subsection{\@tocline{2}{0pt}{3.4pc}{2pc}{}}
\makeatother

\emergencystretch=2em
\clubpenalty=10000 \widowpenalty=10000 \displaywidowpenalty=10000
\setcounter{tocdepth}{1}

\title[Low-regularity solutions of Ideal MHD]{Ideal MHD below the classical well-posedness threshold}
\author{Matteo Giardi}
\thanks{M.G. was supported by the International Max Planck Research School Mathematics
in the Sciences (IMPRS MiS) and thanks his advisor L\'aszl\'o
Sz\'ekelyhidi Jr. for helpful comments on an earlier version of this
manuscript.}

\date{}

\begin{document}
\begin{abstract}
We establish local existence and uniqueness of solutions for the ideal incompressible magnetohydrodynamics system posed on
$[0,T]\times\Rn$, $n\ge2$, with a nonzero constant initial magnetic field $\Bzero$ and arbitrary divergence-free velocity data $v_0\in H^s$, in the range $(n+1)/2<s\le n/2+1$. The proof uses a Lagrangian wave--Hodge reformulation and exploits an Alfvén null--structure hidden in the pressure forcing. In particular, the constructed Eulerian solutions are induced by a bi-Lipschitz measure-preserving flow map. Zhang first identified this null-structure in \cite{Zhang2024}; the present work
provides a self-contained bridge from that Lagrangian theory to the Eulerian Cauchy
problem.
\end{abstract}

\maketitle

\tableofcontents

\section{Introduction}
\label{sec:introduction}

We study the ideal incompressible magnetohydrodynamics (MHD) system set on $[0,T]\times\Rn_x$.  In Eulerian variables
$(t,x)$ this reads 
\begingroup
\renewcommand{\theHequation}{MHD}
\begin{equation}\label{eq:mhd}\tag{MHD}
\begin{cases}
\partial_tv+\nabla_v v-\nabla_B B+\nabla p=0,\\
\partial_tB+\nabla_v B-\nabla_B v=0,\\
\Div_x v=\Div_x B=0,
\end{cases}
\end{equation}
\endgroup
with data $(v,B)|_{t=0}=(v_0,\Bzero)$. Throughout this paper,
$\Bzero\equiv\Binf\in\Rn\setminus\{0\}$ is a constant non-zero vector, and
we write $v_{\rm A}:=|\Binf|$ for the corresponding Alfv\'en speed.

Let $(u\otimes z)^{ij}=u^iz^j$. For a tensor $T$ we take the
divergence in its second index,
\[
 (\Div_xT)^j:=\partial_iT^{ji}.
\]
Put $b:=B-\Binf$. A pair $(v,B)$ with
$v,b\in C([0,T],L^2_{\rm loc})$ is a distributional solution of \eqref{eq:mhd} with data
$(v_0,\Binf)$ if $\Div_{x}v=\Div_{x}B=0$ in $\mathcal D'$ and
\[
 \int_0^T  \int_{\Rn}
 \Big(v\cdot\partial_t\varphi+(v\otimes v-B\otimes B):\nabla\varphi\Big) dx dt
 +\int_{\Rn}v_0\cdot\varphi(0) dx=0
\]
for every divergence-free
$\varphi\in C_c^\infty([0,T)\times\Rn;\Rn)$, while
\[
 \int_0^T  \int_{\Rn}
 \Big(B\cdot\partial_t\psi+(B\otimes v-v\otimes B):\nabla\psi\Big) dx dt
 +\int_{\Rn}\Binf\cdot\psi(0) dx=0
\]
for every $\psi\in C_c^\infty([0,T)\times\Rn;\Rn)$.\\

\textbf{Brief notation overview.} We denote by $x$ and $y$ the two variables given by the change of coordinates $x=X(t,y)$, with $X$ the Lagrangian flow of the velocity $v$, and call $w:=X-\mathrm{id}$ the \emph{displacement} associated to $X$. The
Eulerian vector fields $v,B$ determine the following vector fields along $X$, written in the fixed Cartesian frame:
\begin{equation}
    \vLag:=v\circ X=\partial_tX, \quad \BLag:=B\circ X=\mathcal BX, \quad \bLag:=b\circ X=\BLag-\Binf.
\end{equation}
We sometimes call $y$ the label coordinate and $x$ the Cartesian or Eulerian coordinate. We reserve $d=d_{y}$ for exterior differentiation in the labels and for the componentwise
differential of label maps; thus $dX$, $dw$, $dF$ and
$d(\Box X)$ are componentwise label differentials.
We reserve $\nabla$ for the flat Eulerian derivative, so that $\nabla_v B=(v\cdot\nabla)B$ has its classical meaning, write $\nabla^{y}$ for the flat
label derivative, and use $\eframe_i=\left.\partial_{x^i}\right|_{X}$ for the Cartesian vector fields along $X$ and $E_i=dX^{-1}\eframe_i=X^*\partial_{x^i}$ for their genuine pullbacks to label space, defined in \eqref{eq:component-derivative}. 
Fourier transforms are taken with the unitary normalisation. We use a hat for the spatial
transform in $y$ and a tilde for the spacetime transform in $(t,y)$; a transform in a single
scalar variable is identified explicitly where it occurs. Harmless dimensional constants are
denoted by $c_n$ or $C$. 
Matrix and tensor norms are Frobenius norms unless stated otherwise; $\mathrm{id}$ denotes
the identity map and $\mathrm{Id}$ the identity matrix. For reference, we collect the most recurrent objects in \S\ref{ssec:table}.\\

\textbf{Comparison with the classical local theory and literature.}
Writing $b=B-\Binf$, the standard Kato--Ponce commutator energy estimate for smooth
solutions gives
\begin{equation}\label{eq:standardHsintro}
    \frac{d}{dt}
 \Bigl(\|v(t)\|_{H^s}^2+\|b(t)\|_{H^s}^2\Bigr)
 \lesssim
 \Bigl(\|\nabla v(t)\|_{L^\infty}
       +\|\nabla b(t)\|_{L^\infty}\Bigr)
 \Bigl(\|v(t)\|_{H^s}^2+\|b(t)\|_{H^s}^2\Bigr).
\end{equation}
Here the principal magnetic terms cancel between the two equations, the
constant-field contribution is skew-adjoint, and the pressure term vanishes
by incompressibility. When $s>\frac n2+1$, the Sobolev embedding
$H^s\hookrightarrow W^{1,\infty}$ closes this estimate and yields the
classical local well-posedness theory. In particular, the velocity is
Lipschitz and its flow is well defined.

This work is concerned with the low-regularity range
\[
 \frac{n+1}{2}<s\leq\frac n2+1,
\]
where the $H^s$ norm no longer controls the Lipschitz norm of the velocity and
the classical transport argument does not close. The decisive advance in this
direction is due to Zhang~\cite{Zhang2024}. For a nonzero constant initial
magnetic field, Zhang identified in Lagrangian coordinates a degenerate
one-dimensional wave system with an Alfv\'en null structure, introduced function spaces adapted to that structure, and proved the
corresponding bilinear null-form estimate. This led to the first construction below the classical threshold, namely local well-posedness for an auxiliary Lagrangian system when $2\leq n\leq4$ and
$\frac{n+1}{2}<s\leq\frac n2+1$. This is the
analytic starting point for the half-derivative Eulerian improvement considered here.

The wave structure mentioned above is the mathematical manifestation of the
\emph{Alfv\'en waves} discovered in \cite{Alfven1942}. Its role in the long-time dynamics of ideal MHD was already exploited by Bardos--Sulem--Sulem~\cite{BardosSulemSulem1988}, who proved global existence and scattering for small perturbations of a strong constant magnetic field. The global dynamics near constant magnetic backgrounds have since been studied extensively; see, for instance, He--Xu--Yu~\cite{HeXuYu2018}, Cai--Lei~\cite{CaiLei2018}, and the references therein. The Lagrangian
approach of Andersson and Kapitanski~\cite{AnderssonKapitanski2023} and the
wave-map formulation for neo-Hookean elasticity developed by
Zhang~\cite{Zhang2023elastic} are then precursors of \cite{Zhang2024} mentioned above.

The self-contained present paper rewrites that Lagrangian mechanism of \cite{Zhang2024} into an invariant, wave--Hodge formulation \eqref{eq:WaveHodgeliterature} and bridges the gap with an Eulerian Cauchy local existence and uniqueness theory. Let
\begin{equation}\label{eq:Bmult}
 \mathcal B:=\Bzero\cdot\nabla^{y}
            =\Binf\cdot\nabla^{y},
 \qquad
 \widehat{\mathcal B f}(\xi)
   =\mathrm i\lambda(\xi)\widehat f(\xi),
 \qquad
 \lambda(\xi):=\Binf\cdot\xi .
\end{equation}
Regularity along this preferred direction plays a crucial role in the
analysis and is measured by the fixed-time scale $Y^{\sreg,\preg}$
defined below in \eqref{eq:Y}; $\sreg$ is the isotropic regularity index and $\preg$ the parallel regularity index.

For the displacement $w=X-\mathrm{id}$, the ideal MHD system will be
rewritten in terms of the Alfv\'en and wave operators
\[
 \mathcal A^\pm:=\partial_t\pm\mathcal B,
 \qquad
 \Box:=\mathcal A^+\mathcal A^-
      =\partial_t^2-\mathcal B^2
\]
as the wave--Hodge system
\begin{equation}\label{eq:WaveHodgeliterature}
 \Box w=-F,\qquad
    \begin{cases}
        \curl_{\gmetric}F=0,\\
        \Div_{\gmetric}F=\mathcal Q[dw],
    \end{cases}
\end{equation}
for the metric $\gmetric=(dX)^{\top}dX$ and with $\mathcal{Q}$ carrying the
aforementioned Alfv\'en null--structure. The free evolution generated by $\Box$ is
non-dispersive: the two carried
half-waves are simply translated in opposite directions along $\Binf$.
Nevertheless, away from the characteristic set the symbol of $\Box$ is
elliptic in the joint Alfv\'en variables, so that the corresponding Duhamel
operator gains a fraction of a derivative in a modulation weight measuring
the distance to that set. Near the characteristic set no such gain is
available, and the nonlinear forcing must instead be handled through the
matching Alfv\'en null structure hidden in $\mathcal{Q}$, which cancels the worst
co-propagating interactions.  In this 1+1 hyperbolic setting, this mechanism is particularly apparent and can be made explicit with the following well-known heuristics. Let us informally write 
\[
 \mathcal Q[dw]\sim Q_0(dw,dw)\sim V^+\otimes V^-,
 \qquad
 V^\pm:=\mathcal A^\pm(dw)=d(\mathcal A^\pm w).
\]
For a free wave, after a fixed rotation for which $\Binf=v_{\rm A}e_1$ and splitting $y=y_1e_1+y'$ one sees that
\begin{equation}\label{eq:motivationintro}
 \mathcal A^\mp V^\pm=0
 \qquad\Longrightarrow\qquad
 V^\pm(t,y)=F^\pm(y_1\pm v_{\rm A}t,y').
\end{equation}
Now consider the null-coordinates change of variables
\[
 (t,y_1)\longmapsto(u,v):=(y_1+v_{\rm A}t,y_1-v_{\rm A}t),
 \qquad
 dtdy_1=\frac1{2v_{\rm A}}dudv.
\]
This is particularly helpful when trying to estimate the non-linearity in $L^2_{t,y_1}$ at fixed $y'$. Indeed, the transport condition \eqref{eq:motivationintro} gives

\begin{equation}\label{eq:keycomputation}
\begin{split}
    \int_\Real\int_\Real |V^+(t,y)|^2|V^-(t,y)|^2d td y_1 &=\int_{\Real}\int_{\Real}
 |F^+(y_1+v_{\rm A}t,y')|^2\,|F^-(y_1-v_{\rm A}t,y')|^2dtdy_1\\
 &=
 \frac1{2v_{\rm A}}
 \int_{\Real}|F^+(u,y')|^2du
 \int_{\Real}|F^-(v,y')|^2dv .
\end{split}
\end{equation}
Thus, the $L^2_{t,y_1}$-norm of the product is controlled by the product of
the two longitudinal $L^2$ norms without requiring $H^s\hookrightarrow L^\infty$. This is exactly the missing ingredient in \eqref{eq:standardHsintro}. The precise lower bound for $s$ then comes from the transverse count that follows: the differentiated unknown $dw$ in $Q_0$ carries $s-1$ isotropic derivatives
and requiring the summation of the transverse $\Real^{n-1}$ sector gives
\[
 s-1>\frac{n-1}{2},
 \qquad\text{equivalently}\qquad
 s>\frac{n-1}{2}+1=\frac{n+1}{2}.
\]
This is just the classical Sobolev well-posedness threshold with one
spatial dimension removed.  Theorem~\ref{thm:null-form} replaces the exact
free-transport relations \eqref{eq:motivationintro} by the corresponding
modulation estimate, while the wave--Hodge argument reinstates the
coefficients and contractions in $\mathcal Q$. We remark that the key computation \eqref{eq:keycomputation} also underlies the uniqueness argument of Proposition \ref{prop:unconditional-uniqueness} and thus forms the backbone of this work and \cite{Zhang2024}. See \S\ref{ssec:overview} for more details.

Loosely speaking, the \eqref{eq:mhd} system behaves like
incompressible Euler in the transverse sector, while magnetic tension and
the null-form cancellations supply the missing longitudinal control. 
No additional longitudinal regularity is required for existence or uniqueness
in Theorem~\ref{thm:main}, which we refer to as regime~\textup{(R1)}. If
\[
 \nabla v_0\in Y^{s-1,\rpar},
 \qquad 0<\rpar\leq1,
\]
then the lifespan can be bounded in terms of the base $H^s$ norm and the
corresponding anisotropic norm. We call this regime~\textup{(R2)}.

A related anisotropic condition already appears in the literature.
Definition~1.2(iv) of Ifrim, Pineau, Tataru, and Taylor~\cite{ifrim}, in the
setting of free-boundary MHD, requires
\[
 \nabla_Bv,\nabla_BB\in H^{s-1/2}
\]
for an $H^s$ state, with $s>\frac n2+1$ in their result.  Although the settings and techniques differ, both
conditions reflect the `smoothing' role of derivatives along the
characteristic magnetic field.

The threshold in \cite{ifrim} is sharp for the class treated there, which
contains the hydrodynamic case $B\equiv0$; the authors are explicit in observing that since the system contains Euler, ``the best we
could hope for in the case of MHD is also $s>\frac{n+2}{2}$'' \cite[p.~5]{ifrim}. The present result shows that once the background field is
bounded away from zero $\Bzero\equiv\Binf\ne0$, for the Cauchy problem
on $\Rn$ the threshold drops to $s>\frac{n+1}2$. In particular, Theorem~\ref{thm:main} produces solutions at an
index where, for $n=2,3$, the corresponding Euler problem on $\Rn$ is strongly
ill-posed by Bourgain--Li~\cite{BourgainLi2015Invent,BourgainLi2021IMRN}. This is the familiar situation in which a
structural hypothesis beats a threshold that is sharp for the unstructured
class. Consistently, the relevant constants below degenerate as $v_{\rm A}\downarrow0$, see \eqref{eq:keycomputation} above, so the two
regimes are separated rather than competing, and no Euler solution is obtained
in the limit. In this sense, \cite{Zhang2024} and the present work isolate the regime in which ideal
magnetohydrodynamics genuinely separates from hydrodynamics.

Beyond the invariant reformulation \eqref{eq:WaveHodgeliterature}, compared to \cite{Zhang2024}, we run the fixed point producing the Lagrangian solution, on its differentiated form but for trial \emph{exact} deformation
matrices $\Ufield$, \emph{without} imposing the volume constraint. We then prove that exactness is preserved, that the resulting displacement
satisfies $dw=\Ufield$, and that $\det d X=\det(\mathrm{Id}+\Ufield)=1$ is propagated dynamically by the coupled
wave--Hodge equations, as a theorem at the regularity of the
fixed-point class (Proposition~\ref{prop:constraint}) in every dimension $n\ge2$. The return
to Eulerian variables is then justified through the weak reconstruction of \S\ref{ssec:eulerian}. Eulerian uniqueness is finally proved in \S\ref{ssec:unconditional-uniqueness}. This is done directly
in the physical difference equations after pulling them back by the
already constructed map $X$; no flow for the comparison solution is needed. The use of an interaction-flux adapted from Ifrim Tataru~\cite{IfrimTataru2025}, substituting the role of the energy, and surfacing the null-structure again is crucial to the argument. Detailed descriptions of the above mechanisms and tools are given in \S\ref{ssec:overview}.\\

\textbf{The wave-adapted spaces and parameters.} The various operators and observations above motivate the scale of spaces we introduce next. These are an anisotropic version of the classical Bourgain $X^{s,b}$ spaces widely used in non-linear dispersive equations, see for example \cite{Tao2006Dispersive}.  

Throughout $n\ge2$, $\frac{n+1}2<s\le\frac n2+1$, $\kappa:=n/2+1-s$, $\breg\in(\frac12,1)$, $\sreg\ge0$, $\preg\in\Real$, with
$\varepsilon:=1-\breg$, and $T\le1$; we refer to $(n,s,\breg,v_{\rm A})$ as the
standing parameters.  The bilinear estimates of
\S\ref{ssec:null-form} are proved for the wider range $\breg\in(\frac12,1]$,
which costs nothing and is recorded there; the upper endpoint is excluded in the
main argument only because the linear estimate \eqref{eq:linear} gains
$T^{1-\breg}$. Moreover, we let $\Lambda:=(1-\Delta_{y})^{1/2}$ and associate to the Alfvén and magnetic derivations the self-adjoint operators
\begin{equation}\label{eq:self-adjoint-derivatives}
    D_t:=-\mathrm i\partial_t,\qquad
    D_{\parallel}:=-\mathrm i\mathcal B, \qquad
    D_{\pm}:=-\mathrm i\mathcal A^{\pm}=D_t\pm D_{\parallel},
\end{equation}
where $\parallel$ denotes the full direction $\Binf$, without normalising its
magnitude, so that the parallel and Alfv\'en/null directional symbols are given by
\begin{equation}\label{eq:taupm}
    \lambda=\Binf\cdot\xi, \qquad\tau_+:=\tau+\lambda,\qquad \tau_-:=\tau-\lambda.
\end{equation}
Define the weights
\begin{equation}\label{eq:weightshm}
     h:=\max\big(\langle\tau_+\rangle,\langle\tau_-\rangle\big),\qquad
 m:=\min\big(\langle\tau_+\rangle,\langle\tau_-\rangle\big).
\end{equation}
and the corresponding modulation-weighted spacetime and fixed-time scales
\begin{equation}\label{eq:Y}
\|f\|_{Y^{\sreg,\preg}_{\breg}}:=\big\|\langle D_{\parallel}\rangle^{\preg}m^{\breg}\Lambda^{\sreg}f\big\|_{L^{2}(\mathbb{R}^{1+n})},
\qquad
\|f\|_{Y^{\sreg,\preg}}:=\big\|\langle D_{\parallel}\rangle^{\preg}\Lambda^{\sreg}f\big\|_{L^{2}(\mathbb{R}^{n})}.
\end{equation}
We finally construct the wave-adapted class directly by its two Cauchy components:
\begin{equation}\label{eq:Hpair}
\|f\|_{H^{\sreg,\preg}_{\breg}}:=
\Big(\|f\|_{Y^{\sreg,\preg}_{\breg}}^{2}
+\|\partial_tf\|_{Y^{\sreg,\preg-1}_{\breg}}^{2}\Big)^{1/2}.
\end{equation}
This definition may be interpreted as follows. We measure
$\sreg$ isotropic derivatives on both Cauchy components, together with
$\preg$ parallel derivatives on $f$ and $\preg-1$ on $\partial_tf$
($\preg=1$ is the principal case), together with $\breg$ derivatives with respect to the regularised distance $m$ to the characteristic set $\{\tau_+=0\}\cup\{\tau_-=0\}$ of $\Box$, which for this degenerate operator is a pair of hyperplanes rather than a cone. A detailed explanation of these spaces is given in \S\ref{ssec:overview}.

To connect the wave-adapted norm with the null-form multiplier, put
$\Omega:=hm^{\breg}$. Unfolding $\Omega$ over the two Alfv\'en directions gives
\begin{equation}\label{eq:pmax}
\Omega
 =h m^{\breg}
 =\max\big(\langle\tau_+\rangle\langle\tau_-\rangle^{\breg},
             \langle\tau_+\rangle^{\breg}\langle\tau_-\rangle\big)
 =\frac{\langle\tau_+\rangle\langle\tau_-\rangle}{m^{\varepsilon}}
\end{equation}
that is, we measure one derivative along each Alfv\'en
characteristic, the smaller of the two discounted by $\varepsilon$. Secondly, from the relations
\[
\max\big(|\tau_+|,|\tau_-|\big)=|\tau|+|\lambda| ,\qquad
\min\big(|\tau_+|,|\tau_-|\big)=\big| |\tau|-|\lambda| \big| ,\qquad
h\ \le\ \langle\lambda\rangle+|\tau|\ \le\ 2h ,
\]
we obtain the equivalent single-multiplier expression
\begin{equation}\label{eq:Hsymbol}
\|f\|_{H^{\sreg,\preg}_{\breg}}\ \simeq\
\big\|\langle D_{\parallel}\rangle^{\preg-1}\Omega(D_+,D_-)\Lambda^{\sreg}f\big\|_{L^2(\mathbb{R}^{1+n})},
\end{equation}
for every $\preg\in\Real$ (again $\preg=1$ is the principal case).  Indeed, Plancherel's identity rewrites the
square of the norm in \eqref{eq:Hpair} as
\[
 \int \langle\xi\rangle^{2\sreg}m^{2\breg}
 \left(\langle\lambda\rangle^{2\preg}
       +\tau^2\langle\lambda\rangle^{2\preg-2}\right)
 |\widetilde f(\tau,\xi)|^2 d\tau d\xi .
\]
Factoring out $\langle\lambda\rangle^{2\preg-2}$, the remaining
weight is $\langle\lambda\rangle^2+\tau^2$.  The preceding comparison
$h\simeq\langle\lambda\rangle+|\tau|$ gives
\[
 (\langle\lambda\rangle^2+\tau^2)^{1/2}\simeq h,
\]
with constants independent of $(\tau,\xi)$.  Thus the integrand is
equivalent to
$\langle\xi\rangle^{2\sreg}\langle\lambda\rangle^{2\preg-2}
\Omega^2|\widetilde f|^2$, which proves \eqref{eq:Hsymbol}.  This
calculation also explains why the second Cauchy component carries one fewer
longitudinal derivative: together, $\langle D_\parallel\rangle f$ and $\partial_t f$ supply the
large characteristic weight $h$.

Finally, we define the longitudinal $H^s$ modulus of the datum by
\begin{equation}\label{eq:datum-modulus}
 \omega_{v_0}(T)
 :=\big\|\min(1,T\langle D_{\parallel}\rangle)
                 \Lambda^{s-1}\nabla v_0\big\|_{L^2},
 \qquad 0<T\le1.
\end{equation}
It satisfies $\omega_{v_0}(T)\to0$ as $T\downarrow0$, and if
$\nabla v_0\in Y^{s-1,\rpar}$, $0<\rpar\le1$, then
\begin{equation}\label{eq:omegav0bound}
    \omega_{v_0}(T)\le T^{\rpar}\|\nabla v_0\|_{Y^{s-1,\rpar}}
\end{equation}
which follows immediately from $\min(1,r)\leq r^\rpar$ for $r\geq 0$ and $0<\rpar\leq 1$.\\
 
\textbf{Conserved quantities.} A field tending to $\Binf\ne0$ carries infinite energy, so the constant part is subtracted. The finite conserved quantities are
\begin{equation}\label{eq:energy}
E:=\tfrac12\int\big(|v|^2+|b|^2\big) dx,\qquad
H_{\rm c}:=\int v \cdot b dx,
\end{equation}
and both integrals are absolutely convergent for a solution in the class of
Theorem~\ref{thm:main}.

\subsection{The main results}
With the preliminary definitions in place, we state our main result. 

\begin{theorem}[local existence and uniqueness below the classical threshold]\label{thm:main}
Fix $\breg\in(\frac12,1)$. Let $n\ge2$, let $\frac{n+1}2<s\le\frac n2+1$, let
$\Bzero\equiv\Binf\in\Rn\setminus\{0\}$, and let $v_0\in H^{s}(\Rn)$ be divergence-free. There are
constants $C\ge1$ and $\delta\in(0,1)$, depending only on
$n,s,\breg$ and $v_{\rm A}$, and there exist a time $T_*>0$, a distributional
solution $(v,B, p)$ of \eqref{eq:mhd} on $[0,T_*]\times\Rn$ with data
$(v_0,\Binf)$, with the following properties:
\begin{enumerate}
\item writing $b:=B-\Binf$, the fields $v$, $b$ and $p$ lie in
$C([0,T_*],H^{s}(\Rn))$. The pressure is normalised by
\[
 p=\Riesz_i\Riesz_j
 \big(v^iv^j-\Binf^ib^j-b^i\Binf^j-b^ib^j\big),
\]
where $\Riesz_i$ are the Euclidean Riesz transforms;
\item writing $X=\mathrm{id}+w$, the solution is carried by a Lagrangian flow
$X:[0,T_*]\times\Rn\to\Rn$ with $X(0,\cdot)=\mathrm{id}$, uniformly bi-Lipschitz
and volume-preserving, $\det dX\equiv1$, such that $v=(\partial_tX)\circ X^{-1}$ and
$B=(\mathcal BX)\circ X^{-1}$;
moreover there is a spatially exact spacetime field
$\Ufield\in H^{s-1,1}_{\breg}(\Rst)$ such that
$\Ufield=dw$ on $[0,T_*]\times\Rn$ and, for constants fixed by
the standing parameters (one may take $C=2C_{\rm lin}$, with $C_{\rm lin}$
the constant of the linear estimate \eqref{eq:linear}),
\begin{equation}\label{eq:main-adapted-bound}
 \|\Ufield\|_{H^{s-1,1}_{\breg}}
 \le C\|v_0\|_{H^s},
 \qquad
 \sup_{t\in[0,T_*]}\|dw(t)\|_{Y^{s-1,1}}\le\delta;
\end{equation}
here $\delta$ is chosen so that the second bound implies
$\|dw\|_{L^\infty_{t,y}}\le\frac12$;
\item the renormalised energy $E$ and cross helicity $H_{\rm c}$ of \eqref{eq:energy} are
conserved on $[0,T_*]$;
\item the solution is unique among all distributional solutions with the same
initial data for which $v$ and $B-\Binf$ belong to
$C([0,T_*],H^s)$. 
\end{enumerate}
The existence time is governed by the longitudinal regularity of the datum. We
distinguish the following two regimes.
\begin{enumerate}
\item[\textbf{\emph{(R1)}}]  $T_*$ depends on $n,s,\breg,
v_{\rm A}$, on $\|v_0\|_{H^{s}}$ and on the function
$T\mapsto\omega_{v_0}(T)$ defined in \eqref{eq:datum-modulus}.
\item[\textbf{\emph{(R2)}}] If $\|\nabla v_0\|_{Y^{s-1,\rpar}}<\infty$
for some $\rpar\in(0,1]$, then $T_*$ depends on $n,s,\breg,\rpar,v_{\rm A}$
and on $\|v_0\|_{H^{s}}+\|\nabla v_0\|_{Y^{s-1,\rpar}}$ alone. Moreover,
the spacetime field in~(2) satisfies
\[
 \|\Ufield\|_{H^{s-1,1+\rpar}_{\breg}}
 \le C\|\nabla v_0\|_{Y^{s-1,\rpar}}.
\]
\end{enumerate}
\end{theorem}

\subsection{Overview of the proof}\label{ssec:overview}

The starting points are the Lagrangian flow equation for the velocity~$v$, namely
\[
 v\circ X=\partial_t X
\]
and the fact that reading the Maxwell--Faraday--Ohm law in \eqref{eq:mhd} as Lie transport of the magnetic field $B$ along $v$, one has
\[
 B\circ X=\mathcal BX,
 \qquad \mathcal B=\Binf\cdot\nabla^y.
\]
The momentum equation along $X$ can then be read as a \emph{constant-coefficients} wave equation:  
\[
 (\partial_tv+\nabla_vv)\circ X=\partial_t\big(v\circ X\big)=\partial_t^2X, \qquad (\nabla_BB)\circ X= \mathcal B(B\circ X)=\mathcal B^2X
\]
that is
\[
    \partial_t^2X-\mathcal B^2X=-(\nabla p)\circ X.
\]
Setting
\[
 w:=X-\mathrm{id},\qquad
 F:=(\nabla p)\circ X,\qquad
 \Box:=\partial_t^2-\mathcal B^2,
\]
and using $\mathcal B^2\mathrm{id}=0$, we obtain the invariant
wave--Hodge system:
\begin{equation}\label{eq:whintro}
 \Box w=-F,\qquad 
    \begin{cases}
        \curl_{\gmetric}F=0,\\
        \Div_{\gmetric}F=\mathcal Q[dw],
    \end{cases}
\end{equation}
for the metric $\gmetric=(dX)^{\top}dX$. The Hodge datum contains the Alfv\'en null-form
\[
 \mathcal Q[dw]
 =-\Kmat^a_i\Kmat^b_jQ_0(\partial_aw^j,\partial_bw^i),
 \qquad
 Q_0(f,g)=\partial_tf \partial_tg-\mathcal Bf \mathcal Bg, \qquad \Kmat=(d X)^{-1}.
\]
Lemma~\ref{lem:weak-pullback} gives the weak geometric calculus in which the
two Hodge equations of \eqref{eq:whintro} are posed, while
Lemma~\ref{lem:wave-hodge-reduction} supplies both the algebraic identity for the null--form and the derivation of \eqref{eq:whintro} from
smooth ideal MHD. 

Existence of a solution for this auxiliary problem is proved by a fixed-point argument run at the level of the differentiated equation for the displacement
$\Ufield$. $w$ is reconstructed only afterwards, and the constraint $dw=\Ufield$ will hold on the interval of definition. At the very end, the return to the Eulerian side from the constructed weak auxiliary solution $X=\mathrm{id}+w$ and Eulerian uniqueness are addressed.\\

\textbf{Wave-adapted spaces, linear and nonlinear wave theory.} Let us first explain what the function spaces \eqref{eq:Y},\eqref{eq:Hpair} measure.  The characteristic
derivatives and Fourier variables of $\Box$ are
\[
 \mathcal A^\pm=\partial_t\pm\mathcal B,
 \qquad \tau_\pm=\tau\pm\Binf\cdot\xi.
\]
If
\[
 h=\max(\langle\tau_+\rangle,\langle\tau_-\rangle),
 \qquad
 m=\min(\langle\tau_+\rangle,\langle\tau_-\rangle),
\]
then the equivalent description \eqref{eq:Hsymbol} says that
$H_{\breg}^{\sreg,\preg}$ measures
\[
 \langle D_\parallel\rangle^{\preg-1}
 h m^{\breg}\Lambda^{\sreg}f.
\]
The factor $\Lambda^{\sreg}$ quantifies isotropic spatial regularity,
$\langle D_\parallel\rangle^{\preg-1}$ adjusts the regularity along
$\Binf$, think of $\preg=1$, and $h m^{\breg}$ measures one full derivative in
the larger Alfv\'en direction and $\breg$ derivatives in the smaller
one.  The smaller modulation $m$ is the regularised distance to the nearer
characteristic branch $\tau_+=0$ or $\tau_-=0$, where $\Box$ is not elliptic; this modulation weight compensates for that degeneracy. This is the usual modulation mechanism for a wave equation, adapted in \cite{Zhang2024} to a
strongly anisotropic operator; the no-endpoint conditions $\frac12<\breg <1$ are spent as follows. 

\textit{Linear wave estimates.} The upper endpoint $\breg <1$ is needed to gain a positive power of $T$ in the
Duhamel/forcing term coming from the solution operator of the linear equation
\[
    \Box u=G,
\]
leading to the linear estimate \eqref{eq:linear}.  We now recall how this classical $H^{\sreg,\preg}_{\breg}$ estimate is obtained. After time
localisation at scale $T$, one splits according to the modulation $m$.  In the
high-modulation region $m>2/T$ the two Japanese brackets are comparable with
the corresponding absolute values, so $|\tau_+\tau_-|\gtrsim hm$ and
division by the (elliptic) symbol of $\Box$ against the $hm^{\breg}$ solution weight
leaves $m^{\breg-1}\lesssim T^{1-\breg}$.  In the low-modulation region
$m\le2/T$, that is, near the characteristic set, no elliptic gain is
possible. There one simply spends the modulation weight $hm^{\breg}\lesssim T^{-\breg}h$, and cures the loss exploiting time localisation which ensures the `energy bound' $\|h\widetilde u\|_{L^2}\lesssim T\|G\|_{L^2}$, on an interval of length $T$; see Lemma~\ref{lem:oscillator}.  The two regions therefore give the same gain $T^{1-\breg}$.

\textit{Wave non-linearity and null-structure.} The condition $\breg>\frac12$ is used in the non-linear part of the argument to gain from the null structure using the `stored regularity' in the modulation weight $m$. The rhs of $\Box$ in \eqref{eq:whintro} contains the bilinear form $Q_0$ which shares with $\Box$ the Alfv\'en factorization:
\[
 Q_0(f,g)=\tfrac12\big(\mathcal A^+f \mathcal A^-g
                         +\mathcal A^-f \mathcal A^+g\big).
\]
Each term pairs a $+$ derivative
of one input with a $-$ derivative of the \emph{other}, while both carry a wave class weight
\[
    hm^{\breg}=\frac{\langle\tau_+\rangle\langle\tau_-\rangle}{m^{\varepsilon}}
\]
ensuring one derivative per type is controlled, in particular the product in $Q_0$, although each factor loses one Alfv\'en derivative, still has both available and stays in the class. Quantitatively, when performing estimates in the wave class $H_{\breg}^{s-1,\preg}$ in which the linear estimate \eqref{eq:linear} from above is read,  dividing $|\tau_\pm|$ coming from the differentiation $\mathcal{A}^\pm$ by
the available weight $hm^{\breg}=\frac{\langle\tau_+\rangle\langle\tau_-\rangle}{m^{\varepsilon}}$ leaves the decay $\langle\tau_\mp\rangle^{-\breg}$
in the \emph{other} null variable. The square of this kernel is integrable in
the two null variables precisely for
$\breg>\frac12$, equivalently $\varepsilon=1-\breg<\frac12$. This is the content of Theorem~\ref{thm:null-form} and results in 
\begin{equation}\label{eq:NFintro}
    \|Q_0(\Ufield,\Ufield)\|_{L^2_tH^{s-1}}
 \lesssim \|\Ufield\|_{H_{\breg}^{s-1,1}}^2.
\end{equation}
The only multiplier in $H_{\breg}^{\sreg,\preg}$ that we haven't justified yet is the weight $\Lambda^\sreg$. Unlike the usual wave operator, $\Box$ contains no transverse spatial
derivatives, leaving the transverse frequencies as parameters.  In the differentiated equation, the null form first produces the Hodge datum for $dF$, which is estimated in $L^2_tH^{s-1}$; transverse frequencies are therefore summed with the isotropic
$\Lambda^{s-1}$ weight.  This is why the core null-form estimate contains
\[
 \int_{\Real^{n-1}}\langle\eta'\rangle^{-2(s-1)} d\eta',
\]
which is finite exactly when $s>(n+1)/2$.  At fixed time, the additional
longitudinal derivative in $Y^{s-1,1}$, we control in \eqref{eq:NFintro}, completes the missing spatial bound
and gives the $L^\infty$, algebra and multiplier estimates needed for the
elliptic/Hodge step coming afterwards.\\

\textbf{Wave estimates meet Hodge estimates.}  The above explained the wave part of \eqref{eq:whintro} and the adapted spaces $H_{\breg}^{\sreg,\preg}$. What is left to be addressed in the auxiliary system is how the Hodge part interacts with the null-form to give the forcing $F$.  From a trial field
$\Ufield$ (recall that we work at the level $\Ufield=d w$ not on \eqref{eq:whintro} directly), the null form first produces the scalar datum
$\mathcal Q[\Ufield]$,  the metric Hodge system recovers $F[\Ufield]$, and
$dF[\Ufield]$ is then the source in the differentiated \eqref{eq:whintro}, namely
\[
 \Box\Ufield=-dF[\Ufield].
\]
This concatenation in fact produces a map $\Ufield\mapsto dF[\Ufield]$ from
$H^{s-1,1}_{\breg}$ to $L^2_tH^{s-1}$ which costs \textit{no derivatives} beyond those
measured by the wave-adapted norm. In addition to the null-form gain explained above, the technical content of this statement is the $-1$ order nature of the metric Hodge system, which compensates for $d$. Expressing it in terms of the flat one, the inverse deformation matrix $\Kmat[\Ufield]=(\mathrm{Id}+\Ufield)^{-1}$ appears \emph{only} as a coefficient of $d$, and gives
\[
 (\Div_{\gmetric},\curl_{\gmetric})
 = (\Div_y,\curl_y)
   +\Mmat[\Ufield]d, \qquad
 \Mmat[\Ufield]:=\Kmat[\Ufield]-\mathrm{Id}.
\]
The second term is a principal-order perturbation, but the fixed-point class we run the argument in, namely $\mathcal{U}_{R,\delta}$ from \eqref{eq:fixedpoint-class}, ensures $\Mmat[\Ufield]$ is $\delta$ small in the right norm, and thus the operator has small coefficients. Standard perturbation theory, namely expansion in a Neumann series, can then be applied. 

Concretely, we first rewrite the Hodge system (see \eqref{eq:pressure-Hessian}) as the resolvent identity
\begin{equation}\label{eq:resolventintro}
    dF[\Ufield]
 =
 \mathscr T_0\mathcal Q[\Ufield]
 +\mathscr T_{\Mmat[\Ufield]}dF[\Ufield] \quad\leadsto \quad dF[\Ufield]=(\mathrm{Id}-\mathscr T_{\Mmat[\Ufield]})^{-1}\mathscr T_0\mathcal Q[\Ufield]
\end{equation}
where $\mathscr T_0$ is an order-zero Calderón-Zygmund type operator and
$\mathscr T_{\Mmat[\Ufield]}$ consists of multiplication by
$\Mmat[\Ufield]$ followed by an order-zero CZ type operator. This is then a small perturbation of the identity ensuring the displayed invertibility.  This rewriting together with the null--form estimate \eqref{eq:NFintro} then yields
\begin{equation}\label{eq:hessianintro}
    \|dF[\Ufield]\|_{L^2_tH^{s-1}}
 \lesssim_\delta
 \|\mathcal Q[\Ufield]\|_{L^2_tH^{s-1}}
 \lesssim_\delta R^2
\end{equation}
for $\Ufield\in \mathcal{U}_{R,\delta}$. No Christoffel symbols or derivatives of the metric coefficients (which would be of second order in $X$) enter this argument.

The identity \eqref{eq:resolventintro} also makes the dependence of the pressure Hessian $d F[\Ufield]$ quantitative. Indeed, bilinearity of $Q_0$ and of \eqref{eq:resolventintro} in both the argument and the coefficients/geometry leads to
\[
 \|dF[\Ufield]-dF[\Vfield]\|_{L^2_tH^{s-1}}
 \lesssim_\delta
 (R+R^2)\|\Ufield-\Vfield\|_{H_{\breg}^{s-1,1}}.
\]
The term proportional to $R$ comes from replacing one factor of the null form,
whereas the term proportional to $R^2$ comes from varying the geometry. This Hessian pressure-difference estimate is the additional observation that turns the
fixed-point map into a contraction.

This explains the small parameter $\delta$ absent in \cite{Zhang2024}.  It is not a smallness assumption
on the initial datum; indeed, the space--time radius $R\gtrsim\|v_0\|_{H^s}$ may be
arbitrarily large.  Instead, the fixed-point class requires
\[
 \sup_{t\in[0,T]}\|\Ufield(t)\|_{Y^{s-1,1}}\le\delta.
\]
The anisotropic embedding then makes $dX=\mathrm{Id}+\Ufield$ uniformly
invertible, so every trial map $X(t,\cdot)$ in the fixed-point class is globally bi-Lipschitz and the
pullback geometry is meaningful.  This is a small-deformation
condition and short time buys it even
for large data: the homogeneous part of the solution for $\Box$ is controlled by the longitudinal datum modulus
$\omega_{v_0}(T)\to0$, and the Duhamel/forced part gains a positive power of
$T$: $T^{1/2}$ in the fixed-time norm $Y^{s-1,1}$, by the second line of
\eqref{eq:cutoff-duhamel}, and $T^{1-\breg}$ in the spacetime norm, as explained
before.\\

\textbf{The fixed-point argument.} Starting from an exact matrix-valued $1$-form $\Ufield$, define $\mathcal M\Ufield$ on
$[0,T]$ by
\[
 \Box\mathcal M\Ufield=-dF[\Ufield],
 \qquad
 \begin{cases}
  \curl_{\gmetric[\Ufield]}F[\Ufield]=0,\\
  \Div_{\gmetric[\Ufield]}F[\Ufield]=\mathcal Q[\Ufield],
 \end{cases}
\]
with Cauchy data $(0,\nabla v_0)$. After time localisation, the two estimates
defining the class have the schematic form
\[
 \|\mathcal M\Ufield\|_{H_{\breg}^{s-1,1}}
 \le \frac R2+C_\delta T^{1-\breg}R^2, \qquad
 \sup_{t\in[0,T]}
 \|\mathcal M\Ufield(t)\|_{Y^{s-1,1}}
 \le
 C\big(
   \omega_{v_0}(T)+T^{1/2}C_\delta R^2\big).
\]
Small $T$ therefore makes $\mathcal M$ a self-map of the closed exact class
$\mathcal U_{R,\delta}$ with the prescribed Cauchy data.

For two elements of the class, their free-wave evolution parts cancel. Combining the
pressure-difference estimate with the linear Duhamel estimate gives
\[
 \|\mathcal M\Ufield-\mathcal M\Vfield\|_{H_{\breg}^{s-1,1}}
 \lesssim_\delta
 T^{1-\breg}(R+R^2)
 \|\Ufield-\Vfield\|_{H_{\breg}^{s-1,1}}.
\]
Eventually decreasing $T$, Banach's fixed-point
theorem yields a unique fixed point  $\Ufield$ in $\mathcal U_{R,\delta}$.  Uniqueness of $\Ufield$ does not enter the final Eulerian uniqueness argument.

\textit{Off-constraint class.} The fixed-point map is deliberately defined without imposing the volume
constraint on its trial fields. The point is that
the nonlinear condition
\[
 \det(\mathrm{Id}+\Ufield)=1
\]
is not known to be preserved by one application of $\mathcal M$. For
$\Vfield=\mathcal M\Ufield$, the Jacobi identity involves the metric and
Hodge datum determined by $\Vfield$, whereas the wave source used to define
$\Vfield$ was determined by $\Ufield$, see Remark \ref{rem:generalcaseJacobi} and compare to \eqref{eq:jacobiintro} below. These agree only at a fixed point. Another argument for this being the good `relaxed space' to set the contraction argument in is discussed in Remark \ref{rem:weak-limit-heuristic}.\\

\textbf{Back to $w$, the constraint $J=1$ and the Eulerian side.}
Once the fixed point is obtained, the explicit Duhamel formula reconstructs $w$,
and differentiation of that gives $dw=\Ufield$ on $[0,T]$. It remains
to prove that the auxiliary system is genuinely ideal MHD. For $X=\mathrm{id}+w$, set
\[
 J:=\det dX=\det(\mathrm{Id}+\Ufield),
 \qquad
 \varrho:=\log J.
\]
The weak Jacobi identity gives
\begin{equation}\label{eq:jacobiintro}
    \Box\varrho
 =\mathcal Q[dw]
  -\Div_{\gmetric[dw]}F[dw]
 =0.
\end{equation}
Since $\varrho$ has zero Cauchy data, it follows that
$\varrho\equiv0$, and therefore $J\equiv1$.
Together with the bi-Lipschitz bound supplied by $\delta$, this allows us to
return to Eulerian variables. At the available Sobolev regularity, the divergence conditions, the induction equation and the momentum equation are verified directly against test functions. Finally, a Liouville argument identifies the candidate pressure gradient
$F\circ X^{-1}$ with the normalised
Riesz pressure gradient $\nabla p_{\mathrm R}$. This also supplies the missing $L^2$ control on $F$ which is needed
because a priori Lemma~\ref{lem:pressure} gives only $F\in L^\infty_tL^4$ and
$dF\in L^2_tH^{s-1}$, with no $L^2$ bound on $F$ itself. A final step then puts $v,b,p\in C_tH^s$ and closes the Eulerian existence proof.\\

\textbf{The two regimes.}
The two regimes in Theorem~\ref{thm:main} use longitudinal regularity only to
determine the lifespan and the propagated norm. In \textup{(R1)} no additional assumption is made:
$\omega_{v_0}(T)\to0$ for each fixed $H^s$ datum, but its rate is not uniform
on bounded $H^s$ sets. This is because the finite $H^s$ mass may sit at arbitrarily large longitudinal frequencies by becoming more and more concentrated; see the discussion after Lemma \ref{lem:contraction}.
In \textup{(R2)}, any positive longitudinal fraction gives
\[
 \omega_{v_0}(T)
 \le
 T^{\rpar}\|\nabla v_0\|_{Y^{s-1,\rpar}},
\]
so the lifespan can be chosen in terms of norms alone. The same fraction is in fact propagated within the Picard iteration and thus
\[
 \Ufield\in H_{\breg}^{s-1,1+\rpar}
\]
and the fixed point class $\mathcal{U}_{R,\delta}$ need not be modified.\\

\textbf{Uniqueness.} This is the content of
Proposition~\ref{prop:unconditional-uniqueness} and can be seen as a sort of
weak--strong uniqueness statement: the decaying parts of both solutions
belong to the same Eulerian $C_tH^s$ class, while the constructed one is
additionally induced by a bi-Lipschitz measure-preserving map $X$.  Let
$(v,B,p)$ denote this solution and let
$(\widetilde v,\widetilde B,\widetilde p)$ be any other distributional
solution with the same initial data.  Set
\[
 Z^\pm:=v\pm B,
 \qquad
 \widetilde Z^\pm:=\widetilde v\pm\widetilde B,
 \qquad
 \delta Z^\pm:=\widetilde Z^\pm-Z^\pm,
 \qquad
 z^\pm:=\frac12\bigl(Z^\pm+\widetilde Z^\pm\bigr)\mp\Binf,
\]
and, along the reference map,
\[
 \mathsf Z^\pm:=Z^\pm\circ X=\mathcal A^\pm X,
 \qquad
 \delta\mathsf Z^\pm:=\delta Z^\pm\circ X,
 \qquad
 \delta p:=\widetilde p-p,
 \qquad
 \delta\mathsf p:=\delta p\circ X.
\]
Reading the difference equation in Lagrangian coordinates gives
\begin{equation}\label{eq:overview-conservative-difference}
 \mathcal A^\mp\delta\mathsf Z^\pm
 +\Div_{\gmetric}\bigl(
   (\mathsf Z^\pm+\delta\mathsf Z^\pm)
   \otimes\delta\mathsf Z^\mp\bigr)
 +\grad_{\gmetric}\delta\mathsf p=0.
\end{equation}
The usual difference energy does not close.  Indeed, for the full
deformation tensor we have
\[
 \bigl(E_i\mathsf Z^{\pm,j}\bigr)_{i,j}
 \in C\bigl(I;H^{s-1}_y\bigr)
 \hookrightarrow
 L^2\bigl(I;L^2_{y_1}H^{s-1}_{y'}\bigr)
 \hookrightarrow
 L^2\bigl(I;L^2_{y_1}L^\infty_{y'}\bigr),
 \qquad
 s-1>\frac{n-1}{2},
\]
but
\[
 L^2_{y_1}L^\infty_{y'}\not\hookrightarrow L^\infty_y.
\]
Thus the transverse $L^\infty$ control is available, but the longitudinal
$L^\infty$ control is missing.  Here $E_i=\Kmat_i^a\partial_a$,
$\Kmat=(dX)^{-1}$ and, after a fixed rotation,
$\Binf=v_{\rm A}e_1$.  Now writing
\[
 \mathcal E(t):=\frac12\sum_\pm
 \int_{\Real^n}|\delta\mathsf Z^\pm(t,y)|^2dy,
\]
the ordinary energy identity gives
\[
\begin{aligned}
 \partial_t\mathcal E(t)
 &=-\sum_\pm\int_{\Rn}
 E_i\mathsf Z^{\pm,j}
 \delta\mathsf Z^{\mp,i}
 \delta\mathsf Z^{\pm,j}dy\lesssim
 \mathcal E(t)\sum_\pm
 \left\|
 \bigl(E_i\mathsf Z^{\pm,j}\bigr)_{i,j}(t)
 \right\|_{L^\infty_y}.
\end{aligned}
\]
This is the same obstruction as in
\eqref{eq:standardHsintro}, which motivated the whole construction.

The replacement is an interaction functional adapted from Ifrim and
Tataru~\cite{IfrimTataru2025}, namely a transport version of the wave
null-form pairing.  First define the \textit{energy densities} and
\textit{line energies}
\[
 e_\pm:=\frac12|\delta\mathsf Z^\pm|^2,
 \qquad
 \rho_\pm(t,r):=\int_{\Real^{n-1}}e_\pm(t,r,y')dy'.
\]
For $I=[t_0,t_1]$, introduce the \textit{time-localised energy} and the
\textit{interaction flux} norm
\[
 \mathcal E_I:=\sup_{t\in I}\int_{\Rn}(e_++e_-)(t,y)dy,
 \qquad
 \Phi_I:=\left(
 \int_I\int_{\Real}\rho_+(t,r)\rho_-(t,r)dr dt
 \right)^{1/2}.
\]
Letting $t_0\in[0,T^*]$ be a time where the solutions agree, they satisfy
\begin{equation}\label{eq:fluxenergydominationintro}
 \mathcal E_I\lesssim c_I\Phi_I,
\end{equation}
where
\[
 c_I:=
 \sup_{t\in I}\bigl(
   \|\delta Z^+(t)\|_{H^s}+\|\delta Z^-(t)\|_{H^s}\bigr)+\left(\int_I
   \bigl(\|z^+(t)\|_{H^s}^2+\|z^-(t)\|_{H^s}^2\bigr)dt
   \right)^{1/2}.
\]
The usefulness of
\eqref{eq:fluxenergydominationintro} is that it brings back in the game quantities that are propagating along opposite Alfv\'en directions, and in a fashion similar to \eqref{eq:keycomputation} this will compensate for the missing longitudinal $L^\infty$ bound.

Indeed, integrating the local energy identities in $y'$ and \eqref{eq:overview-conservative-difference} give
\begin{equation}\label{eq:densitytransportintro}
 (\partial_t\mp v_{\rm A}\partial_r)\rho_\pm
 +\partial_rf_\pm=h_\pm,
\end{equation}
where $f_\pm$ contain the nonlinear transport and pressure fluxes and
$h_\pm$ are the deformation sources.  Thus the principal parts of
$\rho_\pm$ travel in opposite directions with speed $v_{\rm A}$.

To estimate $\Phi_I$ one reasons geometrically, and this is the key insight
we borrow from
\cite[Sections~6.3.2 and~6.6.3]{IfrimTataru2025}.  Define the
\textit{interaction functional}
\[
 \mathscr I(t):=\int_{r_+>r_-}
 \rho_+(t,r_+)\rho_-(t,r_-)dr_+dr_-.
\]
This counts pairs still awaiting an encounter: the $+$ energy lies to the
right and moves left, while the $-$ energy lies to the left and moves
right.  For the \emph{free transports} these pairs cross the encounter diagonal
$r_+=r_-$ with relative speed $2v_{\rm A}$, and hence
\[
 \mathscr I'(t)
 =-2v_{\rm A}\int_{\Real}\rho_+(t,r)\rho_-(t,r)dr.
\]
Thus $2v_{\rm A}\Phi_I^2$ is the total flux of the pair density through
the encounter diagonal.  Since $\rho_\pm\ge0$, one has
$\mathscr I(t)\ge0$, while agreement at $t_0$ gives
$\mathscr I(t_0)=0$.  After integrating over $I$ and including the \emph{error
terms} in \eqref{eq:densitytransportintro}, one obtains
\begin{equation}\label{eq:fluxbalanceintro}
    \mathscr I(t_1)+2v_{\rm A}\Phi_I^2
 =\text{nonlinear transport, pressure and source contributions}.
\end{equation}
For the free
transports, computing as in \eqref{eq:keycomputation} would produce only the product of the two total line energies.  In this sense the interaction functional is exactly the
transport version of the wave null-form pairing.

The nonlinear transport contribution is bounded by
$Cc_I\Phi_I^2$, while the pressure and source contributions are bounded by
$Cc_I^2\Phi_I^2$. Consequently, after discarding the non-negative endpoint contribution $\mathscr I(t_1)$ in \eqref{eq:fluxbalanceintro}, we reach 
\[
 2v_{\rm A}\Phi_I^2
 \le C(c_I+c_I^2)\Phi_I^2.
\]
Since $c_I\to0$ as $t_1\downarrow t_0$, on a sufficiently short interval
$C(c_I+c_I^2)<2v_{\rm A}$.  Hence $\Phi_I=0$, and
\eqref{eq:fluxenergydominationintro} gives $\mathcal E_I=0$.  A
continuation argument proves equality on the whole lifespan and uniqueness follows.


\subsection{What is left open} Three issues remain. 

\emph{Continuity of the data-to-solution map}: the present argument does not
prove continuity of the Eulerian data-to-solution map. Comparing solutions
arising from two different data would require controlling both the change in
the fixed-point map and composition with two different inverse flows at the
top $H^s$ regularity. We do not pursue this additional step here. Moreover, in
regime~\textup{(R1)} the available lifespan is not uniform on bounded subsets
of $H^s$, because the modulus $\omega_{v_0}$ is not uniformly controlled
there. Accordingly, no Hadamard well-posedness statement is asserted.

\emph{A variable initial field}: the Lagrangian reduction itself does not use
constancy.  For any divergence-free $B_0$, setting
$\mathcal B:=B_0\cdot\nabla^{y}$, one still has $B\circ X=dX B_0=\mathcal BX$
and $\partial_t^2X=\mathcal B^2X-(\nabla p)\circ X$.  What constancy buys is
that $\mathcal B$ is a constant-coefficient operator: $\Box$ is then a Fourier
multiplier, $\Box \mathrm{id}=0$, and $[\Lambda^{\sreg},\mathcal B]=0$ by
\eqref{eq:comm-zero}.  Each of these fails for nonconstant $B_0$.  Since
$\Box \mathrm{id}=-\nabla_{B_0}B_0$, the displacement equation acquires the
source term $\Box w=-(\nabla p)\circ X+\nabla_{B_0}B_0$, and the wave-adapted
norms acquire commutator errors; neither is estimated here.  A joint functional
calculus for the Alfv\'en operators does survive: if $B_0$ is regular enough (and still non-vanishing)
to generate a measure-preserving flow, then $D_t$ and
$D_{\parallel}=-\mathrm i\mathcal B$ are commuting self-adjoint operators, and
the spectral theorem for the associated $\Real^{2}$ unitary action supplies objects of the form $\Phi(D_+,D_-)$.  What is lost is that this action is no longer simultaneously
diagonalised with $\Lambda^{\sreg}$, so the weights
$\langle D_{\parallel}\rangle^{\preg}m^{\breg}\Lambda^{\sreg}$ of \eqref{eq:Y}
no longer combine into a single Fourier multiplier.  Determining the minimal
regularity of $B_0$ for which the commutator estimates close and thus the present construction can be carried out in this generalised setup is
left to future work.

\emph{Restarting the construction}: at a positive time, the evolved magnetic
field is generally nonconstant, so the theorem cannot simply be restarted with
the new Cauchy data. Consequently, no maximal-lifespan or continuation
criterion follows from the present result.\\

\textbf{Degeneration of the parameters and Besov regularity.} The number $\breg$ is an auxiliary proof parameter.
The trace and null-form constants diverge as $\breg\downarrow\frac12$, while the short-time gain
$T^{1-\breg}$ degenerates as $\breg\uparrow1$.  The constants also deteriorate as
$v_{\rm A}\downarrow0$: \eqref{eq:core-const} contains the factor
$v_{\rm A}^{-1/2}$, and, more robustly, the smallness parameter $\delta$ of
\eqref{eq:delta-choice} itself degenerates, since the constant in
\eqref{eq:embed-frac} at $(\sreg,\preg)=(s-1,1)$ is
$\big(\int\langle\xi\rangle^{-2(s-1)}\langle\Binf\cdot\xi\rangle^{-2}d\xi\big)^{1/2}
\simeq v_{\rm A}^{-\kappa}$ for $\kappa>0$, and
$\simeq(\log v_{\rm A}^{-1})^{1/2}$ at
$\kappa=0$. Through the third line of \eqref{eq:Tstar}, the lifespan furnished
by this construction tends to zero for nontrivial data as
$v_{\rm A}\downarrow0$: at $\Binf=0$ the system is
incompressible Euler, for which the range $s\le\frac n2+1$ is not covered by any
local existence theory of this kind.
We believe that an anisotropic Besov $\ell^1$-type adaptation may reach the
remaining modulation and isotropic regularity endpoints.



\section{Legenda of Symbols}\label{ssec:table}

The most common symbols are collected below. Fonts are used as consistently as possible: plain italic
letters denote physical fields, scalar variables and the flow; sans-serif letters denote
vector fields along $X$, label-space tensors, null-frequency variables and explicit wave propagators; calligraphic letters
denote operators and function spaces; fraktur letters denote regularity indices, auxiliary
parameters or kernels.
The letter $D$ carries the self-adjoint Fourier derivatives $D_t$, $D_{\parallel}$, $D_{\pm}$, boldface distinguishes the fixed
background field and the pullback metric, and blackboard boldface is used for the Leray
projection $\mathbb P$ and the Eulerian stress tensor $\mathbb T$ in Lemma~\ref{lem:pressure-recovery}.

\begingroup\small
\setlength{\LTleft}{0pt}\setlength{\LTright}{0pt}
\begin{table}[htbp]
  \centering
  \small
  \renewcommand{\arraystretch}{0.85}
  
  \begin{tabular}{@{}>{\raggedright\arraybackslash}p{0.205\textwidth}>{\raggedright\arraybackslash}p{0.565\textwidth}>{\raggedright\arraybackslash}p{0.155\textwidth}@{}}
    \hline
\emph{symbol} & \emph{meaning} & \emph{introduced} \\ \hline
$\Bzero\equiv\Binf$ & fixed initial field; Alfv\'en speed $v_{\rm A}:=|\Binf|$ & \S\ref{sec:introduction} \\
$v,\ B,\ b;\ \vLag,\ \BLag,\ \bLag$ & velocity, total magnetic field, perturbation $b=B-\Binf$, and the corresponding vector fields along $X$ & \S\ref{sec:introduction} \\
$s,\ \sigma,\ \kappa,\ \varepsilon$ & $\frac{n+1}2<s\le\frac n2+1$; $\sigma=s-1$; $\kappa=\frac n2+1-s$; $\varepsilon=1-\breg$ & \S\ref{sec:introduction} \\
$\sreg,\ \preg,\ \breg$ & isotropic, parallel and modulation regularity indices; $\breg\in(\frac12,1)$ in the main argument & \S\ref{sec:introduction} \\
$x,\ y;\ X,\ w$ & Eulerian coordinate, material label, flow and displacement $X-\mathrm{id}$ & \S\ref{sec:introduction} \\
$\mathcal C_Xf=f\circ X$ & precomposition with the flow & \S\ref{sec:wave-hodge} \\
$d=d_{y}$ & exterior differential and componentwise label differential of $X,w,F$ and $\Box X$ & \S\ref{sec:introduction} \\
$\nabla,\ \nabla^{y}$ & flat Eulerian and flat label connections & \S\ref{sec:introduction} \\
$\eframe_i,\ E_i$ & Cartesian vector field $\partial_{x^i}\circ X$ along $X$ and its genuine pullback $\Kmat^a_i\partial_a$ to label space & \S\ref{sec:wave-hodge}, \eqref{eq:component-derivative} \\
$\Hess$ & label Jacobian $dF$ of a vector field along $X$, $\Hess_a^i=\partial_aF^i$ & \S\ref{sec:wave-hodge}, \S\ref{ssec:pressure-hessian} \\
$\Ccoef_{\Ufield},\ N_{\Ufield}$ & null-form coefficient $\Kmat[\Ufield]^a_i\Kmat[\Ufield]^b_j$ and numerator $Q_0(\Ufield_a^j,\Ufield_b^i)$, contracted in $\mathcal Q[\Ufield]$ & Lemma~\ref{lem:pressure} \\
$\Ufield$ & spatially exact spacetime fixed-point field; at a fixed point, $\Ufield=dw$ on $[0,T]$ & \eqref{eq:fixedpoint-class}, \eqref{eq:potential-reconstruction} \\
$\gmetric[\Ufield],\ \Kmat[\Ufield],\ J,\ \mu$ & pullback metric, inverse deformation matrix, Jacobian determinant and metric measure & \S\ref{sec:wave-hodge}, \eqref{eq:datum} \\
$A,\ \Mmat$ & deformation matrix $A_a^i=\partial_aX^i$ and the deviation $\Mmat=\Kmat-\mathrm{Id}$ & \S\ref{sec:wave-hodge}, \eqref{eq:KMdef} \\
$\mathcal B,\ D_{\parallel}$ & $\Binf \cdot \nabla^{y}$, of symbol $\mathrm i\lambda$; $D_{\parallel}=-\mathrm i\mathcal B$ self-adjoint & \eqref{eq:Bmult}, \eqref{eq:self-adjoint-derivatives} \\
$\mathcal A^{\pm},\ \Box$ & $\partial_t\pm\mathcal B$, and $\mathcal A^{+}\mathcal A^{-}=\partial_t^2-\mathcal B^2$ & \S\ref{sec:introduction} \\
$D_{\pm},\ D_t$ & $-\mathrm i\mathcal A^{\pm}=D_t\pm D_{\parallel}$ and $-\mathrm i\partial_t$ & \eqref{eq:self-adjoint-derivatives} \\
$\lambda,\ \tau,\ \tau_\pm$ & symbols of $D_{\parallel}$ and $D_t$, with $\tau_\pm=\tau\pm\lambda$ & \eqref{eq:Bmult}, \eqref{eq:joint-res} \\
$\mathsf z_1,\ \mathsf z_2,\ \mathsf z$ & the two factor null frequencies and their sum in $(\tau_+,\tau_-)$ coordinates & \eqref{eq:null-coordinate-plancherel} \\
$\Lambda,\ \langle D_{\parallel}\rangle$ & $(1-\Delta_{y})^{1/2}$; the multiplier $\langle\lambda\rangle$ & \S\ref{sec:introduction} \\
$m,\ h,\ \Omega$ & the smaller and larger null weights, and $\Omega=h m^{\breg}$ & \eqref{eq:pmax} \\
$Y^{\sreg,\preg},\ Y^{\sreg,\preg}_{\breg}$ & fixed-time and modulation-weighted spacetime scales & \eqref{eq:Y} \\
$H^{\sreg,\preg}_{\breg}$ & wave-adapted spacetime Hilbert norm & \eqref{eq:Hpair} \\
$Q_0,\ \mathcal Q[\Ufield]$ & bilinear Alfv\'en null form and the nonlinear scalar pressure datum & \eqref{eq:Q0}, \eqref{eq:datum} \\
$F[\Ufield]$ & canonical $L^4$ pressure force along $X$, characterised by the metric div--curl system & \eqref{eq:FU} \\
$\mathcal P_{\gmetric},\ d\mathcal P_{\gmetric}$ & pressure-gradient operator and its label Jacobian, obtained from the flat Hodge resolvent & \eqref{eq:pressure-gradient}, \eqref{eq:pressure-Hessian} \\
$\varrho$ & $\log\det dX$ & \S\ref{sec:wave-hodge} \\
$\Phi(D_{+},D_{-})$ & joint Fourier multiplier with symbol $\Phi(\tau_+,\tau_-)$ & \eqref{eq:calculus} \\
$\nu$ & $d\tau_+ d\tau_-/(2v_{\rm A})$, the null-coordinate measure & \eqref{eq:null-coordinate-plancherel} \\
$\Cprop,\ \Sprop$ & free cosine and sine propagators for $\Box$ & \eqref{eq:propagator} \\
$c_\perp,\ c_\parallel$ & transverse Sobolev integral and one-dimensional null integral & \eqref{eq:embed-frac}, \eqref{eq:core-const} \\
$C_{\rm lin}$ & the constant of the linear estimate & \eqref{eq:linear} \\
$\mathcal U_{R,\delta},\ \mathcal Y,\ R$ & fixed-point class, ambient Hilbert space and radius & \eqref{eq:fixedpoint-class} \\
$\rpar,\ Y^{s-1,\rpar}$ & additional longitudinal regularity exponent of the datum and its fixed-time scale & \eqref{eq:Y}, Thm.~\ref{thm:main} \\
$\omega_{v_0}$ & the longitudinal high-frequency modulus of the datum & \eqref{eq:datum-modulus} \\
$E,\ H_{\rm c}$ & renormalised energy and cross helicity & \eqref{eq:energy} \\ \hline
\end{tabular}
\end{table}
\endgroup

\section{Ideal MHD as a Wave--Hodge System}
\label{sec:wave-hodge}

This section introduces the first-order weak Lagrangian geometry used below
and derives the wave--Hodge system from smooth ideal MHD.  The basic Hodge
solvability statement is proved here; the first-order derivative estimates needed for
the fixed-point argument are deferred to Section~\ref{sec:elliptic}.\\

\textbf{Standing geometric definitions.} Let $X=\mathrm{id}+w:\Real_y^n\to\Real_x^n$ be a possibly time-dependent,
orientation-preserving global bi-Lipschitz map, and suppress the time
dependence at each fixed time.  Label indices are denoted by $a,b$ and
physical Cartesian indices along $X$ by $i,j$; repeated indices are summed.  We write
\[
 \eframe_i(y):=\left.\partial_{x^i}\right|_{X(y)}
\]
for the Cartesian vector fields along $X$. The Cartesian indices $i,j$ are raised and lowered with the Euclidean fiber metric,
which is the identity in the frame $\{\eframe_i\}$, so their position is (mostly) cosmetic; raising or lowering a label index $a,b$ applies $\gmetric$ and does not
leave the components unchanged. Here and below $d=d_y$ is the exterior differential in the labels and acts
componentwise on maps and on vector fields along $X$ written in the frame
$\{\eframe_i\}$; in particular,
\[
 dF=\Hess_a^i dy^a\otimes\eframe_i
 \qquad\text{for }F=F^i\eframe_i, \qquad \Hess_a^i:=\partial_aF^i,
\]
and thus
\[
 dX=A_a^i dy^a\otimes\eframe_i, \qquad A_a^i:=\partial_aX^i
\]
The map $X$ induces the metric and volume measure
\begin{align*}
 \gmetric&:=X^*\langle\cdot,\cdot\rangle=(dX)^\top dX,
 &\gmetric_{ab}&=A_a^iA_b^i,\\
 J&:=\det dX>0,
 &d\mu&:=J dy,
 &\varrho&:=\log J.
\end{align*}
Put
\[
 \Kmat^a_i:=(A^{-1})^a_i,
 \qquad
 A_a^j\Kmat^a_i=\delta^j_i,
 \qquad
 \Kmat^b_iA_a^i=\delta^b_a.
\]
Let $\mathcal C_Xf:=f\circ X$ be the precomposition operator, acting componentwise in the Cartesian frame on Eulerian vector fields and producing vector fields along $X$. The geometry in the $\gmetric$ metric is characterised by the pullback of the Cartesian frame,
\begin{equation}\label{eq:component-derivative}
 E_i:=dX^{-1}\eframe_i
 =\mathcal C_X\partial_{x^i}\mathcal C_X^{-1}
 =\Kmat^a_i\partial_a
\end{equation}
For a vector field $F=F^i\eframe_i$ along $X$, sufficiently regular (we will address this precisely later), set
\begin{align}\label{eq:divcurlg}
 \Div_{\gmetric}F
 &:=E_iF^i
   =\Kmat^a_i\partial_aF^i,\\
 (\curl_{\gmetric}F)^{ij}
 &:=E_iF^j-E_jF^i
 =\Kmat^a_i\partial_aF^j-\Kmat^a_j\partial_aF^i.
\end{align}
For a scalar function $q$ along $X$ also write
\begin{equation}\label{eq:metricgradient}
    (\mathrm{grad}_\gmetric q)^i=E_i q=\Kmat^a_i\partial_aq.
\end{equation}
For a tensor $T=T^{ji}\eframe_j\otimes\eframe_i$ along $X$, we use the
same second-index convention
\begin{equation}\label{eq:metric-tensor-divergence}
 (\Div_{\gmetric}T)^j:=E_iT^{ji}.
\end{equation}
In particular,
\begin{equation}\label{eq:metric-transport-divergence}
 \Div_{\gmetric}(G\otimes F)
 =F^iE_iG+G\Div_{\gmetric}F.
\end{equation}

\begin{remark}[Geometric interpretation]
The almost-everywhere differential of $X$ is intrinsically a
pullback-bundle-valued $1$-form,
\[
 dX\in L^\infty_{\mathrm{loc}}\big(
 T^*\Real_y^n\otimes X^*T\Real_x^n\big),
 \qquad
 dX_y:T_y\Real_y^n\longrightarrow T_{X(y)}\Real_x^n.
\]
When $X$ is smooth, the pullback bundle $X^*T\Real_x^n$ carries the
Euclidean fiber metric and the pullback connection
$\nabla^{X^*}:=X^*\nabla$. The Cartesian frame $\{\eframe_i\}$ along $X$, is parallel, so
\[
 \nabla^{X^*}_V(F^i\eframe_i)=V(F^i)\eframe_i.
\]
Moreover, $dX$ identifies the genuinely pulled-back $\gmetric$-orthonormal frame $\{E_i\}$ of
label tangent space with the orthonormal Cartesian frame
$\{\eframe_i\}$ along $X$.  Consequently for smooth $X$,
\[
 \Div_{\gmetric}F
 =\tr_{\gmetric} \left(
   \nabla^{\mathrm{LC},\gmetric}(dX^{-1}F)\right)
 =\sum_i\left\langle
    \nabla^{X^*}_{E_i}F,\eframe_i
   \right\rangle,
\]
where the divergence in the middle is the Riemannian divergence of the label
vector field $dX^{-1}F$.  Similarly,
$(\curl_{\gmetric}F)^{ij}$ are the components of
$d((dX^{-1}F)^\flat)$ in the orthonormal frame $\{E_i\}$, with the two
Cartesian indices displayed up as fixed above. We remark that this metric curl is distinct from the rowwise label curl
$\Curl_y$ used for $\Rn$-valued $1$-forms later on, see \eqref{eq:Curl}.  Smoothness of $X$ is used in the calculation above only to justify the Levi--Civita terms involving
$d\gmetric$; the definitions
\eqref{eq:component-derivative}--\eqref{eq:divcurlg} ask no more than
$\Kmat\in L^\infty$, so no weak Levi--Civita connection or Christoffel symbols
are introduced anywhere.  Piola's identity is likewise not needed for these definitions, but will
provide their conservative form and the weak integration-by-parts
formula in Lemma~\ref{lem:weak-pullback}.
\end{remark}

\textbf{The metric Hodge system.} Given a metric $\gmetric$ as above, we now study the Hodge system for a vector field $F$ along $X$,
\begin{equation}\label{eq:invdiv}
    \begin{cases}
        \curl_{\gmetric}F=0,\\
     \Div_{\gmetric}F=q.
\end{cases}
\end{equation}
The following shows the existence and uniqueness of its solutions under sufficient integrability assumptions on the data.
\begin{lemma}[Canonical solvability of the metric Hodge system]
\label{lem:hodge-solvability}
Let $X:\Real_y^n\to\Real_x^n$ be an orientation-preserving global
bi-Lipschitz map with $d\mu=J dy$.  Fix $1<p<n$ and put
\[
 p^*:=\frac{np}{n-p}.
\]
For every $q\in L^p(\mu)$ there is a unique vector field $F=F^i\eframe_i$ along $X$, whose Cartesian components belong to
$L^{p^*}(\mu;\Rn)$, satisfying \eqref{eq:invdiv} in the conjugated
distributional sense: equivalently, $u:=F\circ X^{-1}$ satisfies
\[
 \partial_{x^i}u^j-\partial_{x^j}u^i=0,
 \qquad
 \partial_{x^i}u^i=q\circ X^{-1}
 \quad\text{in }\mathcal D'(\Real_x^n).
\]
The solution is
\begin{equation}\label{eq:pressure-gradient}
 \mathcal P_{\gmetric}q
 :=\mathcal C_X\nabla\Delta_x^{-1}\mathcal C_X^{-1}q, \qquad \mathcal{C}_Xf=f\circ X,
\end{equation}
and it satisfies
\begin{equation}\label{eq:canonical-hls}
 \|\mathcal P_{\gmetric}q\|_{L^{p^*}(\mu)}
 \le C_{n,p}\|q\|_{L^p(\mu)}.
\end{equation}
In particular, for
\begin{equation}\label{eq:rHLS}
    r_{\mathrm{HLS}}:=\frac{4n}{n+4},
\end{equation}
one has $1<r_{\mathrm{HLS}}<n$ for every $n\ge2$ and $r_{\mathrm{HLS}}^*=4$; hence
every $q\in L^{r_{\mathrm{HLS}}}(\mu)$ has a unique solution
$\mathcal P_{\gmetric}q\in L^4(\mu)$ which we refer to as the canonical Hodge vector field along $X$.
\end{lemma}

\begin{proof}
The change of variables $x=X(y)$ makes
$\mathcal C_X:L^r(dx)\to L^r(\mu)$ an isometry.  Put
$f:=\mathcal C_X^{-1}q$ and $u:=\nabla\Delta_x^{-1}f$.  The
Hardy--Littlewood--Sobolev inequality \cite{stein} gives
\begin{equation}\label{eq:flat-hls}
 \|u\|_{L^{p^*}(dx)}\le C_{n,p}\|f\|_{L^p(dx)},
\end{equation}
and the flat identities
\[
 \partial_{x^i}u^j-\partial_{x^j}u^i=0,
 \qquad
 \partial_{x^i}u^i=f
\]
hold in distributions.  Thus $F:=\mathcal C_Xu$ solves
\eqref{eq:invdiv} in the stated sense, and the isometry gives
\eqref{eq:canonical-hls}.

For uniqueness, let $H$ be the difference of two $L^{p^*}(\mu)$ solutions
and put $\tilde u:=H\circ X^{-1}$.  Then
$\tilde u\in L^{p^*}(\Rn)$, $\Div_x\tilde u=0$, and $\curl_x\tilde u=0$ in distributions, hence
\[
 \Delta_x\tilde u^j
 =\partial_{x^i}
   (\partial_{x^i}\tilde u^j-\partial_{x^j}\tilde u^i)
  +\partial_{x^j}\partial_{x^i}\tilde u^i
 =0.
\]
By Weyl's lemma $\tilde u$ is smooth, and the mean-value property together
with $p^*<\infty$ gives
$|\tilde u(x)|\le|B_R|^{-1}\|\tilde u\|_{L^1(B_R(x))}
\le c_nR^{-n/p^*}\|\tilde u\|_{L^{p^*}}\to0$ as $R\to\infty$; hence
$\tilde u=0$. Integrability kills the Hodge kernel.
Finally,
$p=r_{\mathrm{HLS}}$ gives $p^*=4$.
\end{proof}

\textbf{Definition of the off-constraint wave--Hodge auxiliary system.} We set some more notation and define the auxiliary wave--Hodge reformulation of \eqref{eq:mhd} without imposing $\det[d X]=1$.  This is the generalised form used later in the fixed-point argument. For scalar functions, define
\begin{equation}\label{eq:Q0}
 Q_0(f,g)
 :=\tfrac12\big(\mathcal A^+f \mathcal A^-g
                 +\mathcal A^-f \mathcal A^+g\big)
 =\partial_tf \partial_tg-\mathcal Bf \mathcal Bg.
\end{equation}
For an $\Rn$-valued $1$-form $\Ufield=(\Ufield_a^j)$ such that
$\mathrm{Id}+\Ufield$ is invertible, put
\begin{equation}\label{eq:datum}
 \Kmat[\Ufield]:=(\mathrm{Id}+\Ufield)^{-1},
 \qquad
 \mathcal Q[\Ufield]
 :=-\Kmat[\Ufield]^a_i\Kmat[\Ufield]^b_j
 Q_0(\Ufield_a^j,\Ufield_b^i).
\end{equation}
If $\Ufield$ is exact and a primitive $w$ with $dw=\Ufield$ exists such that $X=\mathrm{id}+w$ is an
orientation-preserving global bi-Lipschitz map, write
\begin{equation*}
 \gmetric[\Ufield]:= (\mathrm{Id}+\Ufield)^\top(\mathrm{Id}+\Ufield), \qquad
 J[\Ufield]:=\det(\mathrm{Id}+\Ufield), \qquad d\mu[\Ufield]:=J[\Ufield] dy.
\end{equation*}
Translations in the choice of $w$ do not affect these quantities.  Whenever
$\mathcal Q[\Ufield]\in L^{r_{\mathrm{HLS}}}(\mu[\Ufield])$, let
\begin{equation}\label{eq:FU}
    F[\Ufield]
 :=\mathcal P_{\gmetric[\Ufield]}\mathcal Q[\Ufield]
\end{equation}
denote the unique $L^4(\mu[\Ufield])$ solution of the Hodge system
\eqref{eq:invdiv} with $q=\mathcal Q[\Ufield]$, where the operators in that
system, namely \eqref{eq:divcurlg}, are formed from
$E_i[\Ufield]=\Kmat[\Ufield]^a_i\partial_a$.  Lemma~\ref{lem:hodge-solvability}
gives its existence, uniqueness and Hardy--Littlewood--Sobolev bound; its
first-order derivative estimates are proved in Subsection~\ref{ssec:pressure-hessian}.

If $\Ufield=dw$, the off-constraint auxiliary equation considered below, without imposing $J[\Ufield]=1$, is
the wave--Hodge system
\begin{subequations}\label{eq:auxiliary-system}
\begin{align}
&
\begin{cases}
    \Box w=-F[dw],\\
    (w,\partial_tw)|_{t=0}=(0,v_0),
\end{cases}\\
 &
 \begin{cases}
  \curl_{\gmetric[dw]}F[dw]&=0,\\
  \Div_{\gmetric[dw]}F[dw]&=\mathcal Q[dw]
 \end{cases}
 &\label{eq:force-hodge}
\end{align}
\end{subequations}
The last two equations, together with the normalisation
$F[dw]\in L^4(\mu[dw])$, are precisely the defining characterization of
$F[dw]$ furnished by Lemma~\ref{lem:hodge-solvability}; without the
integrability killing the harmonic components, they determine $F[dw]$ only up to $(\nabla H)\circ X$ with $H$
harmonic.\\

\textbf{From Ideal MHD to the wave--Hodge system.} In the following two lemmas, we first record the properties of the weak geometric calculus induced by a bi-Lipschitz orientation-preserving map $X$. We then show that a smooth solution $(v,B)$ of \eqref{eq:mhd} solves the auxiliary wave--Hodge system \eqref{eq:auxiliary-system} with $w=X-\textrm{id}$ and $X$ the flow map of the velocity $v$. Note that the information $\det d X=1$ seems to be missing from \eqref{eq:auxiliary-system}, but it will be recovered dynamically from divergence-free initial data; see \S \ref{ssec:constraint}. Subsection~\ref{ssec:eulerian} later proves that an auxiliary solution for
which the volume constraint has been recovered determines a solution of
ideal MHD, namely the opposite direction.

\begin{lemma}[weak pullback calculus]\label{lem:weak-pullback}
Let $X$ be an orientation-preserving global bi-Lipschitz map as above.
\begin{enumerate}
\item[(i)] $E_i$ is skew-adjoint on $L^2(\mu)$. In particular
  \[
   \int_{\Real_y^n}(E_if)\phi d\mu=-\int_{\Real_y^n}f(E_i\phi)d\mu
  \]
  whenever $f,\phi,E_if,E_i\phi\in L^2(\mu)$.
\item[(ii)] Let $F=F^i\eframe_i$ be a vector field along $X$ whose Cartesian components belong to $L^1_{\rm loc}(\Real^n_y;\Rn)$. If
  $q,c^{ij}\in L^1_{\rm loc}(\mu)$, then
  \[
   \Div_{\gmetric}F=q,
   \qquad
   (\curl_{\gmetric}F)^{ij}=c^{ij}
  \]
  in the conjugated distributional sense if and only if
  \[
   \partial_a(J\Kmat_i^aF^i)=Jq,
   \qquad
   \partial_a\!\left(J\Kmat_i^aF^j-J\Kmat_j^aF^i\right)=Jc^{ij}
   \quad\text{in }\mathcal D'(\Real^n_y).
  \]
  Equivalently, for every $\varphi\in C_c^\infty(\Real^n_y)$,
  \begin{equation}\label{eq:conservativeform}
   \int q\varphi d\mu=-\int F^i\Kmat_i^a\partial_a\varphi d\mu,\qquad
   \int c^{ij}\varphi d\mu=-\int\bigl(F^j\Kmat_i^a-F^i\Kmat_j^a\bigr)
          \partial_a\varphi d\mu.
  \end{equation}
  With the convention in \eqref{eq:metric-tensor-divergence}, the same
  calculation gives, whenever
  $H:=\Div_{\gmetric}T\in L^1_{\rm loc}(\mu)$,
  \begin{equation}\label{eq:conservative-tensor-divergence}
   JH^j
   =\partial_a(J\Kmat_i^aT^{ji})
   \quad\text{in }\mathcal D'(\Real_y^n)
  \end{equation}
  for every locally integrable tensor $T$.
\item[(iii)] If $\mathcal C_X^{-1}F\in W^{1,p}_{\rm loc}$ for some
  $1\le p\le\infty$, then 
  \begin{equation}\label{eq:biLipchainrule}
      \partial_aF^i=A_a^jE_jF^i,
  \end{equation}
  and the conjugated divergence and curl in
  (ii) agree almost everywhere with the pointwise formulas $\Kmat^a_i\partial_aF^i$ and
  $\Kmat^a_i\partial_aF^j-\Kmat^a_j\partial_aF^i$.
\item[(iv)] The conjugated derivatives commute on their common
distributional domain:
  \[
   E_iE_j=E_jE_i.
  \]
\end{enumerate}
\end{lemma}

\begin{proof} We prove each item separately.

\emph{(i)} The change of variables $x=X(y)$ makes
$\mathcal C_X:L^2(dx)\to L^2(\mu)$ unitary, and
$E_i=\mathcal C_X\partial_{x^i}\mathcal C_X^{-1}$ by
\eqref{eq:component-derivative}.  The hypothesis $f,E_if\in L^2(\mu)$ says
exactly that $\mathcal C_X^{-1}f$ lies in the domain of $\partial_{x^i}$, and
likewise for $\phi$.  Since $\partial_{x^i}$ is skew-adjoint on $L^2(dx)$ and
unitary conjugation preserves skew-adjointness together with its domain, $E_i$
is skew-adjoint on $L^2(\mu)$, and the displayed identity is that
skew-adjointness.

\emph{(ii)} Put $u:=F\circ X^{-1}$. Because $u$ and the right-hand sides are
locally integrable, these first-order divergence and curl identities extend
from smooth to compactly supported Lipschitz tests by mollification and
dominated convergence. We may therefore test after composition with $X^{-1}$.
The weak chain rule and change of
variables show that, after precomposition with $X$, the Euclidean divergence equation
$\partial_{x^i}u^i=q\circ X^{-1}$ becomes
\[
 \int q\varphi d\mu
 =-\int F^i(E_i\varphi)d\mu
 =-\int J\Kmat_i^aF^i\partial_a\varphi dy
\]
for every $\varphi\in C_c^\infty(\Real^n_y)$. This is equivalent to the first
conservative identity. Applying the same argument to
$\partial_{x^i}u^j-\partial_{x^j}u^i=c^{ij}\circ X^{-1}$ gives the second weak
pairing and conservative identity, completing the proof of \eqref{eq:conservativeform}. In this form no derivative falls on
$\Kmat$, and we never multiply a distribution by $J^{-1}$.
Applying the divergence calculation to each row $T^{j \cdot}$ proves
\eqref{eq:conservative-tensor-divergence}.

\emph{(iii)} Put $u:=\mathcal C_X^{-1}F$. The Sobolev chain rule gives
\[
 \partial_aF^i
 =A_a^j(\partial_{x^j}u^i)\circ X
 =A_a^jE_jF^i.
\]
Multiplying by $\Kmat$ gives
\[
 \Kmat_i^a\partial_aF^j=E_iF^j.
\]
Taking the trace and the antisymmetric part gives the pointwise formulas
\[
 \Kmat_i^a\partial_aF^i,
 \qquad
 \Kmat_i^a\partial_aF^j-\Kmat_j^a\partial_aF^i.
\]
Since the distributional derivatives of $u$ have these almost-everywhere
representatives after precomposition with $X$, they agree with the conjugated
distributional divergence and curl.

\emph{(iv)} This follows directly from
$E_i=\mathcal C_X\partial_{x^i}\mathcal C_X^{-1}$ and the commutation of the
Cartesian derivatives.
\end{proof}

We now turn to the reduction of smooth ideal MHD to the wave--Hodge system.

\begin{lemma}[Wave--Hodge reduction]
\label{lem:wave-hodge-reduction}
Let $X=X(t,y)$ be smooth in spacetime, let $w:=X-\mathrm{id}$ and let
$\Kmat=(dX)^{-1}$.  Then the following statements hold.
 
\emph{(i) Algebraic pressure identity.}
Without using an evolution equation or the constraint $J=1$,
\begin{equation}\label{eq:datum-geometric}
 \mathcal Q[dw]
 =-\tr \left(
   \Kmat\mathcal A^+(dw)\Kmat\mathcal A^-(dw)\right)
 =-\Kmat^a_i\Kmat^b_j
   Q_0(\partial_aw^j,\partial_bw^i).
\end{equation}
 
\emph{(ii) Smooth MHD reduction.}
Assume in addition that $X$ is the Lagrangian flow of a smooth solution of
\eqref{eq:mhd}, with $X(0)=\mathrm{id}$ and initial magnetic field $\Binf$.
If
$\mathcal Q[dw]\in L^{r_{\mathrm{HLS}}}(\mu)$ and the pressure is normalised so
that $(\nabla p)\circ X\in L^4(\mu)$ then
\[
 (\nabla p)\circ X=F[dw]=\mathcal{P}_{\gmetric[dw]}\mathcal Q[dw],
\]
and $(w,F[dw])$ satisfies \eqref{eq:auxiliary-system}. In particular, this holds when
$v,b\in C([0,T];H^\infty)$ decay together with all their derivatives and $p$ is
the Riesz pressure \eqref{eq:return-pressure}.
\end{lemma}
 
\begin{proof}
For the algebraic identity let $\Ufield=dw$.  The two terms in the definition
of $Q_0$ give the same contraction after exchanging $(a,i)$ with $(b,j)$.
Therefore
\[
 \Kmat^a_i\Kmat^b_jQ_0(\Ufield_a^j,\Ufield_b^i)
 =\tr \left(
   \Kmat\mathcal A^+\Ufield\Kmat\mathcal A^-\Ufield\right),
\]
and \eqref{eq:datum-geometric} follows directly from the definition \eqref{eq:datum}.
 
It remains to derive the system from smooth ideal MHD.  Let $(v,B,p)$ be the
smooth solution, so that $\partial_tX=v\circ X$ and $X(0,\cdot)=\mathrm{id}$.
 
\emph{The frozen-in law.}  Since $\Div_xv=\Div_xB=0$, the induction equation is
$\partial_tB+[v,B]=0$ with $[v,B]=\nabla_vB-\nabla_Bv$, that is, $B$ is Lie
transported by $v$.  In labels this reads
\begin{equation}\label{eq:frozen-in}
 B\circ X=dX\Binf=\mathcal BX ,
\end{equation}
the second equality being the definition $\mathcal B=\Binf\cdot\nabla^{y}$,
since $(\mathcal BX)^i=\Binf^a\partial_aX^i=(dX\Binf)^i$. Once the induction
equation has been put in this Lie-transport form, the label identity does not
use the volume constraint $\det dX=1$.

The corresponding Elsasser fields satisfy
\[
 (v\pm B)\circ X=\partial_tX\pm\mathcal BX=\mathcal A^\pm X.
\]
Since $\Div_x(v\pm B)=0$, the weak pullback calculus gives directly
\[
 \Div_{\gmetric}(\mathcal A^\pm X)
 =\bigl(\Div_x(v\pm B)\bigr)\circ X=0.
\]
 
\emph{The two material derivative terms.}  For any smooth Eulerian field $g$ the chain
rule gives $\partial_t(g\circ X)=(\partial_tg+\nabla_vg)\circ X$.  Taking $g=v$
and using $v\circ X=\partial_tX$,
\[
 (\partial_tv+\nabla_vv)\circ X=\partial_t\big(v\circ X\big)=\partial_t^2X .
\]
For the magnetic term, $\nabla^{y}(B\circ X)=\big((\nabla_xB)\circ X\big)dX$;
contracting with $\Binf$ and using \eqref{eq:frozen-in} twice,
\[
 \mathcal B(B\circ X)
 =\big((\nabla_xB)\circ X\big)dX\Binf
 =\big((\nabla_xB)\circ X\big)(B\circ X)
 =(\nabla_BB)\circ X ,
\]
while the left-hand side is $\mathcal B(\mathcal BX)=\mathcal B^2X$ by \eqref{eq:frozen-in}.  Hence
\[
 (\nabla_BB)\circ X=\mathcal B^2X .
\]
 
\emph{The wave equation.}  Composing the momentum equation
$\partial_tv+\nabla_vv-\nabla_BB+\nabla p=0$ with $X$ and inserting the two
identities gives
\[
 \partial_t^2X-\mathcal B^2X=-(\nabla p)\circ X ,
 \qquad\text{that is}\qquad
 \Box X=-(\nabla p)\circ X ,
\]
with $\Box=\partial_t^2-\mathcal B^2=\mathcal A^+\mathcal A^-$.  Because
$\Binf$ is constant, $\mathcal B\mathrm{id}=\Binf$ and hence
$\Box\mathrm{id}=0$, so with $F:=(\nabla p)\circ X$,
\[
 \Box w=-F ,
 \qquad
 w(0,\cdot)=0 ,
 \qquad
 \partial_tw(0,\cdot)=v_0 .
\]
The vector fields along $X$ are $\vLag=v\circ X=\partial_tw$ and, by
\eqref{eq:frozen-in}, $\BLag=B\circ X=\Binf+\mathcal Bw$, so that
$\bLag=\mathcal Bw$.
 
\emph{The Hodge equations.}  By \eqref{eq:component-derivative} the precomposition operator
$\mathcal C_X$ is multiplicative and intertwines $\partial_{x^i}$ with $E_i$:
\[
 \mathcal C_X(g_1g_2)=(\mathcal C_Xg_1)(\mathcal C_Xg_2),
 \qquad
 \mathcal C_X\partial_{x^i}=E_i\mathcal C_X .
\]
Applied to $F^j=\mathcal C_X(\partial_{x^j}p)$, the second identity gives
\[
    E_iF^j=\mathcal C_X(\partial_{x^i}\partial_{x^j}p).
\]
We deduce immediately that $\curl_{\gmetric}F=0$ by
the symmetry of the Euclidean Hessian, while
$\Div_{\gmetric}F=E_iF^i=\mathcal C_X(\Delta_xp)$.  Taking $\Div_x$ of the
momentum equation and using $\Div_xv=\Div_xB=0$ gives
\[
 \Delta_xp=(\partial_{x^i}B^j)(\partial_{x^j}B^i)
           -(\partial_{x^i}v^j)(\partial_{x^j}v^i) ,
\]
Precomposing this identity with $X$, using the multiplicativity above, and
recalling that $\vLag=\partial_tw$ and $\BLag=\Binf+\mathcal Bw$, we have
\[
 E_i\vLag^j=\Kmat^a_i\partial_t\partial_aw^j ,
 \qquad
 E_i\BLag^j=\Kmat^a_i\mathcal B\partial_aw^j ,
\]
and therefore
\begin{align*}
 \Div_{\gmetric}F
 &=(E_i\BLag^j)(E_j\BLag^i)-(E_i\vLag^j)(E_j\vLag^i)\\
 &=\Kmat^a_i\Kmat^b_j\Big[
     \big(\mathcal B\partial_aw^j\big)\big(\mathcal B\partial_bw^i\big)
    -\big(\partial_t\partial_aw^j\big)\big(\partial_t\partial_bw^i\big)\Big]\\
 &=-\Kmat^a_i\Kmat^b_jQ_0\big(\partial_aw^j,\partial_bw^i\big)
  =\mathcal Q[dw]
\end{align*}
by part~(i).  The $L^4(\mu)$ uniqueness of Lemma~\ref{lem:hodge-solvability}
then gives $F=F[dw]$, and $X(0)=\mathrm{id}$, $\partial_tX(0)=v_0$ give the
Cauchy data of \eqref{eq:auxiliary-system}.
\end{proof}
 


\section{Wave-Adapted Anisotropic Spaces}
\label{sec:anisotropic-spaces}

Because $\mathcal B$ has constant coefficients, the Alfv\'en, longitudinal and flat calculi are
simultaneous Fourier multipliers. This section proves the three fixed-time estimates used later: the
anisotropic embedding, the trace theorem and the product estimates. Without further mention, the estimates below are first proved for Schwartz functions and eventually extended by continuity. The arguments below adapt the standard proofs for Sobolev spaces and classical
Bourgain spaces; for the latter, see for instance
\cite{Tao2006Dispersive}. Their first anisotropic versions in the present MHD
setting were established by Zhang in \cite{Zhang2024}. The formulations and
proofs given here extend that analysis to the level of generality required for
regime~\textup{(R1)} of Theorem~\ref{thm:main}.

\subsection{Embeddings and traces}

\begin{lemma}[anisotropic Sobolev embedding]\label{lem:Linftyembed}
Let
\begin{equation}\label{eq:Y-embed-range}
 \sreg>\frac{n-1}{2},\qquad \preg\ge0,\qquad \sreg+\preg>\frac n2 .
\end{equation}
Then every $f\in Y^{\sreg,\preg}(\Rn)$ satisfies
\begin{equation}\label{eq:embed-frac}
 \|f\|_{L^\infty(\Rn)}
 \le
 \frac{C(n,\sreg,\preg,v_{\rm A})}
 {(\sreg+\preg-\frac n2)^{1/2}}
 \|f\|_{Y^{\sreg,\preg}} .
\end{equation}
Both strict inequalities in \eqref{eq:Y-embed-range} are necessary.
\end{lemma}

\begin{proof}
By Fourier inversion and Cauchy--Schwarz,
\[
 \|f\|_{L^\infty}
 \le C_n\mathcal J(\sreg,\preg)^{1/2}\|f\|_{Y^{\sreg,\preg}},
 \qquad
 \mathcal J(\sreg,\preg)
 :=\int_{\Rn}
 \langle\xi\rangle^{-2\sreg}
 \langle\lambda(\xi)\rangle^{-2\preg} d\xi .
\]
Write $\xi=\mathbf e_1\xi_1+\xi'$, after rotating $\Binf$ to $v_{\rm A}\mathbf e_1$ and using
$\langle v_{\rm A}\xi_1\rangle\simeq_{v_{\rm A}}\langle\xi_1\rangle$, the change
$\xi'=\langle\xi_1\rangle\zeta$ gives
\[
 \mathcal J(\sreg,\preg)
 \simeq_{v_{\rm A},\preg}
 c_\perp(n,\sreg)
 \int_{\Real}\langle\xi_1\rangle^{n-1-2\sreg-2\preg} d\xi_1,
 \qquad
 c_\perp(n,\sreg)
 :=\int_{\Real^{n-1}}\langle\zeta\rangle^{-2\sreg} d\zeta .
\]
The two factors are finite precisely under \eqref{eq:Y-embed-range}; the last integral is
$O((\sreg+\preg-n/2)^{-1})$ near its endpoint.  For
$0\le\preg\le1$ the remaining comparison constant is uniform in $\preg$.
This proves \eqref{eq:embed-frac}.  Conversely, evaluation at a point is bounded only if the
reciprocal weight belongs to $L^2_\xi$; truncating that weight proves necessity of both strict
inequalities.
\end{proof}
\begin{corollary}[traces in the $Y$ scale]\label{cor:trace-Y}
For $\breg>\frac12$, $\sreg\ge0$, any $\preg$, and every
$f\in Y^{\sreg,\preg}_{\breg}$,
$$\sup_t\big\|\langle D_{\parallel}\rangle^{\preg}\Lambda^{\sreg}f(t)\big\|_{L^{2}}\ \le\
C_{\breg} \|f\|_{Y^{\sreg,\preg}_{\breg}},$$
and $t\mapsto f(t)$ is continuous into $Y^{\sreg,\preg}$. Consequently
\begin{equation}\label{eq:trace-class}
H^{\sreg,\preg}_{\breg}\ \hookrightarrow\ C(\Real;Y^{\sreg,\preg})\cap C^1(\Real;Y^{\sreg,\preg-1}),
\qquad
\sup_t\big(\|f(t)\|_{Y^{\sreg,\preg}}+\|\partial_tf(t)\|_{Y^{\sreg,\preg-1}}\big)
\le C_{\breg}\|f\|_{H^{\sreg,\preg}_{\breg}}.
\end{equation}
\end{corollary}

\begin{proof}
Recall the definitions \eqref{eq:taupm},\eqref{eq:weightshm} and that the hat denotes the pure-space transform while the tilde denotes the time-space Fourier transform. At fixed $\xi$, Fourier inversion in $\tau$ and Cauchy--Schwarz give
\[
 |\widehat f(t,\xi)|
 \le \Big(\int_{\Real} m^{-2\breg}(\tau,\lambda) d\tau\Big)^{1/2}
      \Big(\int_{\Real} |m^{\breg}\widetilde f(\tau,\xi)|^2 d\tau\Big)^{1/2},
 \qquad
 \sup_\lambda\int_{\Real} m^{-2\breg}(\tau,\lambda) d\tau
 \le2\int_{\Real}\langle z\rangle^{-2\breg} dz.
\]
Multiply by $\langle\lambda\rangle^{\preg}\langle\xi\rangle^{\sreg}$ and integrate in
$\xi$.  Dominated convergence yields continuity.  Applying the same argument to $\partial_tf$
gives the bound and the continuity of $t\mapsto(\partial_tf)(t)$ into
$Y^{\sreg,\preg-1}$.  Since $\langle\lambda\rangle^{\preg-1}\le
\langle\lambda\rangle^{\preg}$ gives $Y^{\sreg,\preg}\hookrightarrow
Y^{\sreg,\preg-1}$, and $\partial_tf$ is the distributional time derivative of
$f$, the fundamental theorem of calculus in the Bochner sense yields
$f(t)=f(t_0)+\int_{t_0}^{t}(\partial_tf)(r) dr$ in $Y^{\sreg,\preg-1}$; hence
$f\in C^1(\Real;Y^{\sreg,\preg-1})$ and \eqref{eq:trace-class} follows.
\end{proof}

\subsection{Products}

\begin{proposition}[products in the $Y$ scale]\label{prop:tame-product}
Let $\preg\ge0$ and
\begin{equation}\label{eq:Y-product-range}
 \sreg>\frac{n-1}{2},\qquad \sreg+\preg>\frac n2.
\end{equation}
Then $Y^{\sreg,\preg}$ is a Banach algebra and is a multiplier on $H^{\sreg}$:
\begin{equation}\label{eq:Y-product}
 \|fg\|_{Y^{\sreg,\preg}}\le C\|f\|_{Y^{\sreg,\preg}}\|g\|_{Y^{\sreg,\preg}},
 \qquad
 \|fg\|_{H^{\sreg}}\le C\|f\|_{H^{\sreg}}\|g\|_{Y^{\sreg,\preg}}.
\end{equation}
By interpolation, the second estimate also holds with $H^\varsigma$ in place of $H^{\sreg}$ for
$0\le\varsigma\le\sreg$.  In particular, at $(\sreg,\preg)=(s-1,1)$,
\begin{equation}\label{eq:tame-sp}
 \|fg\|_{H^\varsigma}\le C\|f\|_{H^\varsigma}\|g\|_{Y^{s-1,1}}
 \quad(0\le\varsigma\le s-1),
\end{equation}
and, after taking $L^2$ in time and using Corollary~\ref{cor:trace-Y},
\[
 \|fg\|_{L^2_tH^{s-1}}
 \le C\|f\|_{L^2_tH^{s-1}}\|g\|_{H^{s-1,1}_{\breg}}.
\]
The hypotheses \eqref{eq:Y-product-range} coincide with the embedding range
\eqref{eq:Y-embed-range}.
\end{proposition}

\begin{proof}
Write
\[
 \mathfrak w_{\sreg,\preg}(\xi):=\langle\xi\rangle^{-2\sreg}\langle\lambda(\xi)\rangle^{-2\preg},
 \qquad \mathfrak w_{\sreg,0}(\xi):=\langle\xi\rangle^{-2\sreg}.
\]
We first prove the two weighted convolution inequalities, which hold in the range \eqref{eq:Y-product-range},
\begin{equation}\label{eq:weight-convolution}
 \mathfrak w_{\sreg,\preg}*\mathfrak w_{\sreg,\preg}\le C\mathfrak w_{\sreg,\preg},
 \qquad \mathfrak w_{\sreg,0}*\mathfrak w_{\sreg,\preg}\le C\mathfrak w_{\sreg,0}.
\end{equation}
Write $\xi=e_1\xi_1+\xi'$ and rotate coordinates so that $\lambda(\xi)=v_{\rm A}\xi_1$ and use
$\langle v_{\rm A}t\rangle\simeq_{v_{\rm A}}\langle t\rangle$.  Set
\[
 d:=n-1,\qquad \zeta:=\xi_1,\qquad
 A:=\langle\eta_1\rangle,\qquad
 B:=\langle\zeta-\eta_1\rangle,
 \qquad \alpha:=2\sreg-d>0.
\]
We need the following preliminary estimate for the transverse convolution. On the region $A\le B$,
\begin{equation*}\label{eq:convineq}
    \int_{\Real^{d}}(A^2+|z|^2)^{-\sreg}
 (B^2+|\xi'-z|^2)^{-\sreg} dz
 \le C A^{-\alpha}(B^2+|\xi'|^2)^{-\sreg},
\tag{*}
\end{equation*}
and the symmetric estimate holds on $B<A$.  To see this, if
$|\xi'|\le2B$, bound the second factor by $B^{-2\sreg}$ and use
$\int(A^2+|z|^2)^{-\sreg}dz=C A^{-\alpha}$.  If
$|\xi'|>2B$, split at $|z|=|\xi'|/2$.  On the first part the second
factor is $O(|\xi'|^{-2\sreg})$ and the first factor has integral
$CA^{-\alpha}$.  On the second part the first factor is
$O(|\xi'|^{-2\sreg})$, while the second has integral $CB^{-\alpha}$;
this is bounded by $CA^{-\alpha}|\xi'|^{-2\sreg}$ because $A\le B$.
This proves~\eqref{eq:convineq}.

Consider first the convolution of two copies of
$\mathfrak w_{\sreg,\preg}$.  On $A\le B$, estimate~\eqref{eq:convineq} leaves
\[
 C A^{-\alpha-2\preg}B^{-2\preg}
   (B^2+|\xi'|^2)^{-\sreg}.
\]
Since $\langle\zeta\rangle\le A+B\le2B$ there,
\[
 B^{-2\preg}(B^2+|\xi'|^2)^{-\sreg}
 \lesssim
 \langle\zeta\rangle^{-2\preg}
 (\langle\zeta\rangle^2+|\xi'|^2)^{-\sreg}.
\]
The remaining integral in $\eta_1$ is finite because
$\alpha+2\preg>1$.  The region $B<A$ is identical after interchanging
$A$ and $B$.  This proves the first inequality in
\eqref{eq:weight-convolution}.

For the mixed convolution, the factor carrying the longitudinal weight is
the one with scale $B$.  On $A\le B$, after~\eqref{eq:convineq} and after extracting
$(\langle\zeta\rangle^2+|\xi'|^2)^{-\sreg}$, it remains to verify
\[
 \sup_{\zeta\in\Real}
 \int_{A\le B}A^{-\alpha}B^{-2\preg} d\eta_1<\infty.
\]
For bounded $\zeta$ this follows directly from
$A^{-\alpha}B^{-2\preg}\lesssim
\langle\eta_1\rangle^{-\alpha-2\preg}$.  If $|\zeta|\ge2$, split at
$|\eta_1|=|\zeta|/2$.  On the central part $B\simeq\langle\zeta\rangle$,
and its contribution is bounded by
\[
 C\langle\zeta\rangle^{-2\preg}
 \int_0^{|\zeta|}\langle t\rangle^{-\alpha} dt.
\]
This is uniformly bounded when $\alpha+2\preg>1$; for $\alpha=1$ the
integral grows only logarithmically.  On the complementary part subject to
$A\le B$, one has $B\ge A\gtrsim\langle\zeta\rangle$, so the integrand is
bounded by $A^{-\alpha-2\preg}$ and is integrable.  On $B<A$, the symmetric
transverse estimate leaves $B^{-\alpha-2\preg}$ and the scale $A$ controls
$\langle\zeta\rangle$; integration after the change
$\eta_1\mapsto\zeta-\eta_1$ completes the second inequality.  Thus the two
conditions needed for \eqref{eq:weight-convolution} are exactly
$\alpha>0$ and $\alpha+2\preg>1$, namely \eqref{eq:Y-product-range}.

We now apply these inequalities to the product.  For Schwartz functions put
\[
 F(\xi):=\langle\xi\rangle^{\sreg}
          \langle\lambda(\xi)\rangle^{\preg}\widehat f(\xi),
 \qquad
 G(\xi):=\langle\xi\rangle^{\sreg}
          \langle\lambda(\xi)\rangle^{\preg}\widehat g(\xi).
\]
Writing each Fourier factor as a weighted function times
$\mathfrak w_{\sreg,\preg}^{1/2}$ and using Cauchy--Schwarz in the
convolution variable gives
\[
 |\widehat{fg}(\xi)|^2
 \le
 (\mathfrak w_{\sreg,\preg}*\mathfrak w_{\sreg,\preg})(\xi)
 (|F|^2*|G|^2)(\xi).
\]
The first inequality in \eqref{eq:weight-convolution} cancels the target
weight, and Tonelli's theorem gives
\[
 \|fg\|_{Y^{\sreg,\preg}}^2
 \le C\int (|F|^2*|G|^2)(\xi) d\xi
 =C\|F\|_2^2\|G\|_2^2.
\]
For the multiplier estimate, instead put
\[
 F_0(\xi):=\langle\xi\rangle^{\sreg}\widehat f(\xi),
 \qquad
 G_{\preg}(\xi):=\langle\xi\rangle^{\sreg}
 \langle\lambda(\xi)\rangle^{\preg}\widehat g(\xi).
\]
Writing $\widehat f=\mathfrak w_{\sreg,0}^{1/2}F_0$ and
$\widehat g=\mathfrak w_{\sreg,\preg}^{1/2}G_{\preg}$, the second
inequality in \eqref{eq:weight-convolution} gives
$\|fg\|_{H^{\sreg}}\le C\|f\|_{H^{\sreg}}
\|g\|_{Y^{\sreg,\preg}}$ in the same way.

At the other endpoint, \eqref{eq:embed-frac} gives
$\|fg\|_2\le\|f\|_2\|g\|_\infty
\le C\|f\|_2\|g\|_{Y^{\sreg,\preg}}$.
Interpolation of the multiplication operator $f\mapsto fg$ between
$H^0$ and $H^{\sreg}$ proves the asserted $H^\varsigma$ estimates.  
\end{proof}

\begin{remark}[interpolation in the parallel index]\label{rem:Y-interpolation}
At fixed $\sreg$, the scale $Y^{\sreg,\preg}$ is, on the Fourier side, a
weighted $L^2$ scale whose weight
$\langle\xi\rangle^{2\sreg}\langle\lambda\rangle^{2\preg}$ depends
multiplicatively on $\preg$.  The Stein--Weiss interpolation theorem for
weighted $L^2$ spaces, see \cite[Theorem~5.5.3]{bergh}, therefore gives, with
equal norms,
\[
 \big[Y^{\sreg,\preg_0},Y^{\sreg,\preg_1}\big]_{\theta}
 =Y^{\sreg,(1-\theta)\preg_0+\theta\preg_1},
 \qquad 0\le\theta\le1 .
\]
In particular, a linear operator that is bounded on $Y^{\sreg,0}=H^{\sreg}$
and on $Y^{\sreg,1}$ is bounded on every $Y^{\sreg,\preg}$ with
$0\le\preg\le1$, with the interpolated operator norm. In particular, assume
that $\sigma>(n-1)/2$ and $\sigma+1>n/2$. Multiplication by an element of
$Y^{\sigma,1}$ is bounded on
$Y^{\sigma,0}=H^{\sigma}$ by \eqref{eq:tame-sp} and on $Y^{\sigma,1}$ by
\eqref{eq:Y-product}, hence
\begin{equation}\label{eq:Y-multiplier-interp}
 \|fg\|_{Y^{\sigma,\preg}}
 \le C\|f\|_{Y^{\sigma,\preg}}\|g\|_{Y^{\sigma,1}},
 \qquad 0\le\preg\le1 .
\end{equation}
This is how the parallel index is interpolated in the elliptic estimates of
Section~\ref{sec:elliptic}.
\end{remark}

\section{Linear and Null-Form Estimates}\label{sec:wave-estimates}
This section first records the constant-field Fourier calculus and then establishes the
linear and bilinear estimates used by the elliptic and fixed-point arguments. The estimates in this section adapt the standard Bourgain-space arguments,
for which we refer to \cite{Tao2006Dispersive}, to the anisotropic
Alfv\'en geometry. This adaptation was first carried out by Zhang in
\cite{Zhang2024}; here we use the slightly more general formulations required
by our function-space framework.

\begin{lemma}[the constant-field Fourier representation]\label{lem:alfven-action}
With $(\tau,\xi)$ the spacetime Fourier variables, $\lambda$ as in
\eqref{eq:Bmult}, and $\tau_\pm=\tau\pm\lambda$, one has
\begin{equation}\label{eq:joint-res}
 \widetilde{\mathcal Bf}=\mathrm i\lambda\widetilde f,
 \qquad
 \widetilde{\mathcal A^+f}=\mathrm i\tau_+\widetilde f,
 \qquad
 \widetilde{\mathcal A^-f}=\mathrm i\tau_-\widetilde f,
 \qquad
 \widetilde{\Box f}=-\tau_+\tau_-\widetilde f .
\end{equation}
Consequently $D_\pm=-\mathrm i\mathcal A^\pm$ are commuting self-adjoint
operators on $L^2(dt dy)$ and, for every Borel function $\Phi$,
\begin{equation}\label{eq:calculus}
 \Phi(D_+,D_-)f
 :=\big(\Phi(\tau_+,\tau_-)\widetilde f\big)^\vee
\end{equation}
on its maximal domain.  Moreover
\[
 \mathrm e^{r\mathcal B}f(t,y)=f(t,y+r\Binf),
 \qquad
 \mathrm e^{r\mathcal A^\pm}f(t,y)=f(t+r,y\pm r\Binf),
\]
so these groups are isometries on every $L^p$ and unitary on $L^2$.  In particular,
\begin{equation}\label{eq:comm-zero}
 [\Lambda^{\sreg},\mathcal B]=0
 \qquad(\sreg\in\Real),
\end{equation}
and all corresponding Fourier-multiplier commutators vanish.
\end{lemma}

\begin{proof}
All statements follow directly from the Fourier symbols and the translation formulas.
\end{proof}

\subsection{The linear estimate}
\label{ssec:linear}

Fix $\chi\in C^\infty_c(\Real)$, with $0\le\chi\le1$, $\chi\equiv1$ on
$[-1,1]$ and $\supp\chi\subset[-2,2]$, and recall
$\varepsilon=1-\breg$, $\breg\in(\frac12,1)$, $T\leq 1$.  

The free propagators are
\begin{equation}\label{eq:propagator}
\Cprop(t):=\cos\big(t|\mathcal B|\big)
=\tfrac12\big(\mathrm e^{t\mathcal B}+\mathrm e^{-t\mathcal B}\big),
\qquad
\Sprop(t):=\frac{\sin\big(t|\mathcal B|\big)}{|\mathcal B|}
=\int_0^{t}\tfrac12\big(\mathrm e^{r\mathcal B}+\mathrm e^{-r\mathcal B}\big) dr ,
\end{equation}
where the value of the symbol of $\Sprop(t)$ at $\lambda=0$ is $t$.  Thus $\Cprop(t)$ and
$\Sprop(t)$ are averages of translations along $\Binf$, with
$L^p$-operator norms at most $1$ and $|t|$, respectively. For
$\lambda\in\Real$, write
\[
 h_\lambda(\tau):=\max(\langle\tau+\lambda\rangle,\langle\tau-\lambda\rangle),
 \qquad
 m_\lambda(\tau):=\min(\langle\tau+\lambda\rangle,\langle\tau-\lambda\rangle),
 \qquad
 \Omega_\lambda:=h_\lambda m_\lambda^{\breg}.
\]
The entire linear estimate results from the following uniform scalar oscillator bound.

\begin{lemma}[localised scalar oscillator]\label{lem:oscillator}
Let $\breg\in(\frac12,1)$, $0<T\le1$, $a_0,a_1\in\mathbb C$, and
$g\in L^2(\Real)$.  With
\[
 C_\lambda(t):=\cos(t|\lambda|),\qquad
 S_\lambda(t):=\frac{\sin(t|\lambda|)}{|\lambda|},\quad S_0(t):=t,
\]
put
\[
 u_\lambda(t):=\chi(t)C_\lambda(t)a_0+\chi(t)S_\lambda(t)a_1
 +\chi(t/T)\int_0^tS_\lambda(t-t')g(t') dt'.
\]
Then, uniformly in $\lambda$ and $T$,
\begin{equation}\label{eq:scalar-oscillator}
 \big\|\Omega_\lambda\widetilde u_\lambda\big\|_{L^2_\tau}
 \le C(\chi,\breg)\Big(
 \langle\lambda\rangle|a_0|+|a_1|
 +T^{1-\breg}\|g\|_{L^2_t}\Big).
\end{equation}
\end{lemma}

\begin{proof}
For either sign, setting $\rho=\tau\mp|\lambda|$ gives
\[
 m_\lambda(\tau)\le\langle\rho\rangle,\qquad
 h_\lambda(\tau)\lesssim\langle\lambda\rangle\langle\rho\rangle.
\]
Since $\chi$ is Schwartz, this bounds the two exponential pieces of
$\chi C_\lambda$ by $C\langle\lambda\rangle$.  If $|\lambda|\ge1$, the same
argument applied to
\[
 S_\lambda(t)=\frac{\mathrm e^{it|\lambda|}-\mathrm e^{-it|\lambda|}}
 {2i|\lambda|}
\]
bounds $\chi S_\lambda$ uniformly.  If $|\lambda|\le1$, the family
$\chi S_\lambda$ is bounded in $H^2_t$, while
$\Omega_\lambda(\tau)\lesssim\langle\tau\rangle^{1+\breg}$.  Since
$1+\breg<2$, the $H^2_t$ bound controls this weight and proves the two data
bounds in \eqref{eq:scalar-oscillator}.

For the source set
\[
 u_D(t):=\int_0^tS_\lambda(t-t')g(t') dt',
 \qquad u_T:=\chi(t/T)u_D .
\]
Distributionally,
$(\partial_t^2+\lambda^2)u_D=g$, with zero Cauchy data at $t=0$.
For $|t|\le2T$, the bounds
$|S_\lambda(r)|\le|r|$, $|\partial_rS_\lambda(r)|\le1$, and
$|\lambda S_\lambda(r)|\le1$ give, by Cauchy--Schwarz on the interval
between $0$ and $t$,
\[
 |u_D(t)|\lesssim |t|^{3/2}\|g\|_2,
 \qquad
 |\partial_tu_D(t)|+|\lambda u_D(t)|
 \lesssim |t|^{1/2}\|g\|_2.
\]
Integrating these pointwise bounds in $t$ gives
\begin{equation}\label{eq:oscillator-local}
 \|u_D\|_{L^2(|t|\le2T)}\lesssim T^2\|g\|_2,\qquad
 \|\partial_tu_D\|_{L^2(|t|\le2T)}
 +\|\lambda u_D\|_{L^2(|t|\le2T)}
 \lesssim T\|g\|_2 .
\end{equation}
Writing $\chi_T(t)=\chi(t/T)$, the product rule and
\eqref{eq:oscillator-local} imply
\[
 \|u_T\|_2\lesssim T^2\|g\|_2,\qquad
 \|\partial_tu_T\|_2+\|\lambda u_T\|_2\lesssim T\|g\|_2.
\]
Moreover,
\[
 (\partial_t^2+\lambda^2)u_T
 =\chi_Tg+2T^{-1}\chi'(t/T)\partial_tu_D
   +T^{-2}\chi''(t/T)u_D.
\]
All three terms are supported where $|t|\le2T$, and therefore
\begin{equation}\label{eq:oscillator-cutoff-source}
 \|(\partial_t^2+\lambda^2)u_T\|_2\lesssim\|g\|_2.
\end{equation}
This is the only point at which derivatives of the time cutoff enter.

It remains to recover the modulation weight.  Split the Fourier transform
into the regions $m_\lambda\le2/T$ and $m_\lambda>2/T$.  On the first,
$\Omega_\lambda=h_\lambda m_\lambda^{\breg}\lesssim
T^{-\breg}h_\lambda$ and
$h_\lambda\lesssim1+|\tau|+|\lambda|$.  Plancherel's identity gives
\[
 \|\mathbf 1_{\{m_\lambda\le2/T\}}
       \Omega_\lambda\widetilde u_T\|_2
 \lesssim T^{-\breg}\|h_\lambda\widetilde u_T\|_2
 \lesssim T^{1-\breg}\|g\|_2,
\]
where $T\le1$ absorbs the $T^2\|g\|_2$ contribution.  On the second
region, $m_\lambda>2/T\ge2$, so both Japanese brackets are comparable
with the corresponding absolute values and hence
\[
 |(\tau+\lambda)(\tau-\lambda)|\gtrsim h_\lambda m_\lambda,\qquad
 \frac{\Omega_\lambda}{|(\tau+\lambda)(\tau-\lambda)|}
 \lesssim m_\lambda^{\breg-1}\lesssim T^{1-\breg}.
\]
Since the Fourier symbol of $\partial_t^2+\lambda^2$ is
$-\tau^2+\lambda^2=-(\tau+\lambda)(\tau-\lambda)$,
\eqref{eq:oscillator-cutoff-source} gives
\[
 \|\mathbf 1_{\{m_\lambda>2/T\}}
       \Omega_\lambda\widetilde u_T\|_2
 \lesssim T^{1-\breg}\|g\|_2.
\]
Combining the two regions proves the source bound and the lemma.
\end{proof}

\begin{theorem}[the linear estimate]\label{thm:linear}
Let $\breg\in(\frac12,1)$, $\varepsilon:=1-\breg$,
$\sreg\ge0$, $\preg\in\Real$ and $0<T\le1$. For every
$f_0\in Y^{\sreg,\preg}$,
$f_1\in Y^{\sreg,\preg-1}$ and
$G\in L^{2}_tY^{\sreg,\preg-1}(\Rst)$ there is
$u$ on $\Rst$ with
\[
 \Box u=G\quad\text{on }[0,T],\qquad
 (u,\partial_tu)|_{t=0}=(f_0,f_1),
\]
satisfying
\begin{equation}\label{eq:linear}
\big\|u\big\|_{H^{\sreg,\preg}_{\breg}}
 \le C_{\rm lin}\Big(
 \big\|f_0\big\|_{Y^{\sreg,\preg}}
 +\big\|f_1\big\|_{Y^{\sreg,\preg-1}}
 +T^{\varepsilon}\big\|G\big\|_{L^{2}_tY^{\sreg,\preg-1}}\Big),
\end{equation}
where $C_{\rm lin}=C(\chi,\breg)$ is independent of $T\le1$, of the
field, of the data and of the indices $\sreg,\preg$. Moreover,
\[
 u\in C([0,T],Y^{\sreg,\preg})
 \cap C^1([0,T],Y^{\sreg,\preg-1}).
\]
\end{theorem}

\begin{proof}
Take
\[
 u=\chi(t)\Cprop(t)f_0+\chi(t)\Sprop(t)f_1
   +\chi(t/T)\int_0^t\Sprop(t-t')G(t') dt'.
\]
It has the asserted equation and data on $[0,T]$.  At each spatial frequency
$\xi$, apply Lemma~\ref{lem:oscillator} with
$\lambda=\Binf\cdot\xi$, multiply by
$\langle\lambda\rangle^{\preg-1}\langle\xi\rangle^{\sreg}$, and
integrate in $\xi$.  The equivalent multiplier description \eqref{eq:Hsymbol}
then gives \eqref{eq:linear}.  The continuity statement follows from
Corollary~\ref{cor:trace-Y}.
\end{proof}

The following uniqueness statement for the free evolution is used twice
below: to identify auxiliary solutions with the fixed point in
Theorem~\ref{thm:existence}, and to propagate the volume constraint in
Proposition~\ref{prop:constraint}.

\begin{lemma}[uniqueness for the homogeneous wave equation]\label{lem:wave-uniqueness}
Let $0<T<\infty$ and let $u\in C^1([0,T];L^2(\Rn))$ satisfy
$\mathcal Bu\in C([0,T];L^2(\Rn))$,
\[
 \Box u=0\quad\text{in }\mathcal D'((0,T)\times\Rn),
 \qquad
 u(0)=\partial_tu(0)=0 .
\]
Then $u\equiv0$ on $[0,T]$. The same conclusion holds componentwise for
vector- and matrix-valued $u$.
\end{lemma}

\begin{proof}
Let
\[
 |\mathcal B|:=(-\mathcal B^2)^{1/2},
 \qquad P_N:=\mathbf 1_{[0,N]}(|\mathcal B|),
 \qquad
 u_N:=P_Nu
\]
The Fourier multiplier $P_N$ commutes with $\partial_t$ and $\mathcal B$.
Since $\mathcal B^2P_N$ is bounded on $L^2$, the equation
$\partial_t^2u_N=\mathcal B^2u_N$ and the assumed time continuity give
$u_N\in C^2([0,T];L^2)$. Hence
\[
 \mathcal E_N(t):=\frac12\left(
   \|\partial_tu_N(t)\|_2^2
   +\|\mathcal Bu_N(t)\|_2^2
 \right).
\]
By the skew-adjointness of $\mathcal B$,
\begin{align*}
 \mathcal E_N'(t)
 &=\langle\partial_tu_N,\mathcal B^2u_N\rangle
   +\langle\mathcal Bu_N,\mathcal B\partial_tu_N\rangle\\
 &=-\langle\mathcal B\partial_tu_N,\mathcal Bu_N\rangle
   +\langle\mathcal Bu_N,\mathcal B\partial_tu_N\rangle
 =0.
\end{align*}
Thus $\mathcal E_N(t)=\mathcal E_N(0)=0$. Since $P_N\to\mathrm{Id}$ strongly
on $L^2$, letting $N\to\infty$ gives $\partial_tu=0$, and $u(0)=0$ gives
$u\equiv0$.
\end{proof}

\subsection{The null-form estimate}
\label{ssec:null-form}

The nonlinearity in \eqref{eq:auxiliary-system} is the bilinear Alfv\'en null form $Q_0$ from \eqref{eq:Q0}. Without further mention, the estimates below are first proved for Schwartz functions and eventually extended by continuity.

\begin{theorem}[fractional null-form estimate]\label{thm:null-form}
Let $n\ge2$, $\breg\in(\frac12,1]$, $s>\frac{n+1}2$ and
$\rpar\in[0,1]$. Then
\begin{equation}\label{eq:null-form}
 \|Q_0(f,g)\|_{L^2_tY^{s-1,\rpar}}
 \le C\Big(
   \|f\|_{H^{s-1,1+\rpar}_{\breg}}
   \|g\|_{H^{s-1,1}_{\breg}}
  +\|f\|_{H^{s-1,1}_{\breg}}
   \|g\|_{H^{s-1,1+\rpar}_{\breg}}
 \Big),
\end{equation}
where $C=C(n,s,\breg,v_{\rm A})$ is uniform for
$\rpar\in[0,1]$. In particular,
\begin{equation}\label{eq:diagonalNF}
     \|Q_0(f,f)\|_{L^2_tY^{s-1,\rpar}}
 \le 2C\|f\|_{H^{s-1,1}_{\breg}}
          \|f\|_{H^{s-1,1+\rpar}_{\breg}}.
\end{equation}
\end{theorem}

Recall the notation in \eqref{eq:taupm} and \eqref{eq:self-adjoint-derivatives}. Put $\mathbf e=\Binf/v_{\rm A}$, write
$\xi=\xi_1\mathbf e+\xi'$, and set
$\mathsf z=(\tau_+,\tau_-)$.  The change
$(\tau,\xi_1)\mapsto(\tau_+,\tau_-)$ has Jacobian
$2v_{\rm A}$, so
\begin{equation}\label{eq:null-coordinate-plancherel}
 d\nu(\mathsf z):=\frac{d\tau_+ d\tau_-}{2v_{\rm A}},
 \qquad
 \|f\|_{L^2(\Rst)}^2
 =\iint|\widetilde f(\mathsf z,\xi')|^2 d\nu(\mathsf z) d\xi'.
\end{equation}
For a total frequency $(\mathsf z,\xi')$, let
$(\mathsf z_2,\eta')$ denote the frequency of the $g$-factor and
$(\mathsf z_1,\xi'-\eta')$ that of the $f$-factor. Then
$\mathsf z=\mathsf z_1+\mathsf z_2$, and the full spatial frequency associated
with $\eta=(\mathsf z_2,\eta')$ is
\[
 \eta=(\xi_1(\mathsf z_2),\eta')=\left(\frac{\tau_+(\mathsf z_2)-\tau_-(\mathsf z_2)}
 {2v_{\rm A}}, \eta'\right).
\]
For use below, retain the multiplier notation
\begin{equation}\label{eq:null-weight-operator}
 \Omega_{\sreg,\preg}
 :=\langle D_\parallel\rangle^{\preg-1}
   \Omega(D_+,D_-)\Lambda^{\sreg},
 \qquad \Omega=hm^{\breg},
\end{equation}
so that, by \eqref{eq:Hsymbol},
$\|\Omega_{\sreg,\preg}f\|_2\simeq
\|f\|_{H^{\sreg,\preg}_{\breg}}$.

\begin{lemma}[the resonance identity]\label{lem:resonance}
At factor frequencies $\mathsf z_1,\mathsf z_2$, the symbol of $Q_0$ is
\[
 -q_0(\mathsf z_1,\mathsf z_2),\qquad
 q_0:=\tfrac12\big(
 \tau_+(\mathsf z_1)\tau_-(\mathsf z_2)
 +\tau_-(\mathsf z_1)\tau_+(\mathsf z_2)\big).
\]
\end{lemma}

\begin{proof}
This is immediate from \eqref{eq:joint-res} and the definition \eqref{eq:Q0}.
\end{proof}

\begin{lemma}[the pointwise null-weight bound]\label{lem:pointwise}
For $0\le\breg\le1$,
\begin{equation}\label{eq:null-weight-pointwise}
 \frac{|\tau_+|}{\Omega}\le\langle\tau_-\rangle^{-\breg},
 \qquad
 \frac{|\tau_-|}{\Omega}\le\langle\tau_+\rangle^{-\breg}.
\end{equation}
\end{lemma}

\begin{proof}
Since $\Omega=\langle\tau_+\rangle\langle\tau_-\rangle/m^{1-\breg}$,
\[
 \frac{|\tau_+|}{\Omega}
 \le\frac{m^{1-\breg}}{\langle\tau_-\rangle}
 \le\langle\tau_-\rangle^{-\breg},
\]
and the other inequality is symmetric.
\end{proof}

Set
\[
 \mathfrak k(\mathsf z_1,\mathsf z_2;\eta)
 :=\frac{|q_0(\mathsf z_1,\mathsf z_2)|}
 {\Omega(\mathsf z_1)\Omega(\mathsf z_2)\langle\eta\rangle^{s-1}}.
\]
The following estimate contains the whole bilinear machinery.

\begin{proposition}[the core estimate]\label{prop:core}
For $\breg\in(\frac12,1]$ and $s>\frac{n+1}2$,
\begin{equation}\label{eq:core}
 \left\|
 \iint \mathfrak k(\mathsf z-\mathsf z_2,\mathsf z_2;\eta)
 F(\mathsf z-\mathsf z_2,\xi'-\eta')G(\mathsf z_2,\eta')
  d\eta' d\nu(\mathsf z_2)
 \right\|_{L^2_{\mathsf z,\xi'}}
 \le C\|F\|_2\|G\|_2 .
\end{equation}
One may take
\begin{equation}\label{eq:core-const}
 C^2\le
 \frac{c_\perp(n,s-1)c_\parallel(\breg)^2}{2v_{\rm A}},
 \quad
 c_\perp(n,s-1):=\int_{\Real^{n-1}}\langle\eta'\rangle^{-2(s-1)} d\eta',
 \quad
 c_\parallel(\breg):=\int_{\Real}\langle r\rangle^{-2\breg} dr .
\end{equation}
\end{proposition}

\begin{proof}
By the triangle inequality applied to the two terms of $q_0$,
$\mathfrak k\le\mathfrak k_1+\mathfrak k_2$, where
\[
 \mathfrak k_1:=\frac{\frac12|\tau_+(\mathsf z_1)\tau_-(\mathsf z_2)|}
 {\Omega(\mathsf z_1)\Omega(\mathsf z_2)\langle\eta\rangle^{s-1}},
 \qquad
 \mathfrak k_2:=\frac{\frac12|\tau_-(\mathsf z_1)\tau_+(\mathsf z_2)|}
 {\Omega(\mathsf z_1)\Omega(\mathsf z_2)\langle\eta\rangle^{s-1}} .
\]
It therefore suffices to bound each of the two corresponding convolutions and to
add the results.  We treat $\mathfrak k_1$; the argument for $\mathfrak k_2$ is
identical after interchanging the signs $+$ and $-$.
Lemma~\ref{lem:pointwise} and
$\langle\eta\rangle\ge\langle\eta'\rangle$ give
\[
 \mathfrak k_1
 \le\tfrac12
 \langle\tau_-(\mathsf z_1)\rangle^{-\breg}
 \langle\tau_+(\mathsf z_2)\rangle^{-\breg}
 \langle\eta'\rangle^{-(s-1)}.
\]
Denote the right-hand side by
$K_{\mathsf z}(\mathsf z_2,\eta')$, and let $H$ be the convolution
on the left of \eqref{eq:core} with this single cross-term in place of
$\mathfrak k$.  At every total frequency, Cauchy--Schwarz in
$(\mathsf z_2,\eta')$ gives
\begin{align*}
 |H(\mathsf z,\xi')|^2
 &\le
 \left(\iint K_{\mathsf z}(\mathsf z_2,\eta')^2
        d\eta' d\nu(\mathsf z_2)\right)\\
 &\quad\times
 \left(\iint
 |F(\mathsf z-\mathsf z_2,\xi'-\eta')|^2
 |G(\mathsf z_2,\eta')|^2
  d\eta' d\nu(\mathsf z_2)\right).
\end{align*}
The first factor is independent of $\xi'$ and is bounded uniformly in
$\mathsf z$.  Indeed, writing
$r_\pm=\tau_\pm(\mathsf z_2)$, its value is
\[
 \frac{c_\perp(n,s-1)}{8v_{\rm A}}
 \int_{\Real^2}
 \langle\tau_-(\mathsf z)-r_-\rangle^{-2\breg}
 \langle r_+\rangle^{-2\breg} dr_+ dr_-
 =
 \frac{c_\perp(n,s-1)c_\parallel(\breg)^2}
 {8v_{\rm A}}.
\]
Here $1/(2v_{\rm A})$ comes from $d\nu(\mathsf z_2)$, and the additional
factor $1/4$ is the square of the factor $1/2$ in the resonance
identity.

It remains only to integrate the second Cauchy--Schwarz factor.  Tonelli's
theorem and the translation
\[
 (\mathsf z_1,\tilde\xi')
 :=(\mathsf z-\mathsf z_2,\xi'-\eta')
\]
give
\begin{align*}
 &\iiiint
 |F(\mathsf z-\mathsf z_2,\xi'-\eta')|^2
 |G(\mathsf z_2,\eta')|^2
  d\eta' d\nu(\mathsf z_2) d\xi' d\nu(\mathsf z)\\
 &\qquad=
 \left(\iint|F(\mathsf z_1,\tilde \xi')|^2
        d\tilde \xi' d\nu(\mathsf z_1)\right)
 \left(\iint|G(\mathsf z_2,\eta')|^2
        d\eta' d\nu(\mathsf z_2)\right)\\
 &\qquad=\|F\|_2^2\|G\|_2^2.
\end{align*}
This proves \eqref{eq:core} for the cross-term $\mathfrak k_1$, with constant
$\big(c_\perp c_\parallel^2/(8v_{\rm A})\big)^{1/2}$.  Interchanging the
signs $+$ and $-$ proves the same estimate for $\mathfrak k_2$; adding the two
gives \eqref{eq:core-const} and completes the proof.
\end{proof}

\begin{proof}[Proof of Theorem~\ref{thm:null-form}]
Put $a=s-1$.  At the two factor frequencies write
$\zeta_1=\xi-\eta$, $\zeta_2=\eta$ and
$\lambda_j=\Binf\cdot\zeta_j$, so that
$\xi=\zeta_1+\zeta_2$ and $\lambda=\lambda_1+\lambda_2$.
We use
\[
 \langle\xi\rangle^a
 \lesssim_a\langle\zeta_1\rangle^a+\langle\zeta_2\rangle^a,
 \qquad
 \langle\lambda_1+\lambda_2\rangle^{\rpar}
 \le\langle\lambda_1\rangle^{\rpar}
    +\langle\lambda_2\rangle^{\rpar}
\]
where the second inequality is uniform for $0\le\rpar\le1$. By
Lemma~\ref{lem:resonance}, it is enough to estimate the four Fourier integrals
obtained from
$\langle\zeta_i\rangle^a\langle\lambda_j\rangle^{\rpar}$,
$1\le i,j\le2$. For $\theta=0,\rpar$, set
\begin{align*}
 F_\theta(\mathsf z_1,\zeta_1')
 &:=\Omega(\mathsf z_1)\langle\zeta_1\rangle^a
    \langle\lambda_1\rangle^\theta
    |\widetilde f(\mathsf z_1,\zeta_1')|,
 \\
 G_\theta(\mathsf z_2,\zeta_2')
 &:=\Omega(\mathsf z_2)\langle\zeta_2\rangle^a
    \langle\lambda_2\rangle^\theta
    |\widetilde g(\mathsf z_2,\zeta_2')|.
\end{align*}
Since $\zeta_2=\eta$, Proposition~\ref{prop:core} bounds the two terms with
isotropic weight $\langle\zeta_1\rangle^a$ by
\[
 C\|F_{\rpar}\|_2\|G_0\|_2,
 \qquad
 C\|F_0\|_2\|G_{\rpar}\|_2.
\]
The same bounds hold for the two terms with isotropic weight
$\langle\zeta_2\rangle^a$. Indeed,
$q_0(\mathsf z_1,\mathsf z_2)=q_0(\mathsf z_2,\mathsf z_1)$, and at fixed
$(\mathsf z,\xi')$ the
substitution $(\mathsf z_2,\eta')\mapsto(\mathsf z-\mathsf z_2,\xi'-\eta')$
preserves $d\nu d\eta'$ and exchanges the factor frequencies, giving the
version of Proposition~\ref{prop:core} with
$\langle\zeta_1\rangle^a$ in the denominator. Finally,
\[
 \|F_\theta\|_2\simeq\|f\|_{H^{a,1+\theta}_{\breg}},
 \qquad
 \|G_\theta\|_2\simeq\|g\|_{H^{a,1+\theta}_{\breg}}
\]
by \eqref{eq:Hsymbol}. Summing the four estimates and applying Plancherel proves
\eqref{eq:null-form}. The weight inequalities are uniform for
$\rpar\in[0,1]$, and the equivalence constants in \eqref{eq:Hsymbol} are
independent of the parallel index because the factor
$\langle\lambda\rangle^{2\preg-2}$ is extracted exactly. This gives the stated
uniformity.
\end{proof}

The conditions on the parameters become transparent by looking at \eqref{eq:core-const}:
$s>\frac{n+1}2$ is precisely $c_\perp(n,s-1)<\infty$, while
$\breg>\frac12$ is precisely $c_\parallel(\breg)<\infty$.
The dependence on the background field enters only through the $(2v_{\rm A})^{-1}$ density acquired from the change of variables performed after applying the Plancherel identity.  Throughout, the joint functional calculus in $(D_+,D_-)$ used above is legitimate by Lemma~\ref{lem:alfven-action}.

\section{Elliptic Estimates}\label{sec:elliptic}

With the pullback geometry and canonical pressure force along $X$ constructed in
Section~\ref{sec:wave-hodge}, this section develops the first-order derivative
estimates for the Hodge system, that is, the bounds for the Lagrangian pressure
Hessian.  Two preliminaries come first: stability of Sobolev regularity under
the flow, and a coefficient calculus with base space $Y^{s-1,1}$.  The
fixed-time, product, and null-form estimates used here have already been
established in Section~\ref{sec:anisotropic-spaces} and
Subsection~\ref{ssec:null-form}.

\subsection{Sobolev stability under the flow}

\begin{lemma}[properties of precomposition operators]\label{lem:composition}
Let $X=\mathrm{id}+w:\Rn_y\to\Rn_x$ be an orientation-preserving
bi-Lipschitz homeomorphism such that
\[
 \|dX\|_{L^\infty}+\|(dX)^{-1}\|_{L^\infty}\le\mathfrak m,
 \qquad 0<J_-\le\det dX\le J_+.
\]
Put $\Ufield:=dw$ and $\Kmat:=(dX)^{-1}$, and assume
$\Ufield,\Kmat-\mathrm{Id}\in Y^{s-1,1}$.  Recall $\mathcal{C}_{X}^\pm f=f\circ X^\pm$. Then, for every $0\le r\le s$,
\begin{equation}\label{eq:sobolev-composition}
 \|\mathcal C_Xf\|_{H^r_y}\le C_{\mathrm{comp}}(X,r)\|f\|_{H^r_x},
 \qquad
 \|\mathcal C_X^{-1}u\|_{H^r_x}\le C_{\mathrm{comp}}(X,r)\|u\|_{H^r_y},
\end{equation}
where, with $k(r):=\max\{\lceil r\rceil-1,0\}$,
\[
 C_{\mathrm{comp}}(X,r)\le C(n,s,r,v_{\rm A},\mathfrak m,J_\pm)
 \big(1+\|\Ufield\|_{Y^{s-1,1}}
        +\|\Kmat-\mathrm{Id}\|_{Y^{s-1,1}}\big)^{k(r)}.
\]
If these bounds hold uniformly for $X=X(t)$ on a compact interval $I$ and, in addition,
\[
 w\in C(I;H^s),\qquad
 \Ufield,\Kmat-\mathrm{Id}\in C(I;Y^{s-1,1}),
\]
then $u\in C(I;H^r_y)$ implies
$\mathcal C_{X(t)}^{-1}u(t)\in C(I;H^r_x)$. Independently of the Sobolev
time-regularity hypothesis above, suppose that $X(t)$ and $X(t)^{-1}$ vary
locally uniformly on $I$ and that the Jacobian and bi-Lipschitz bounds are
uniform. Then, for every $1\le q<\infty$, every $t_0\in I$, and every
$g\in L^q$,
\begin{equation}\label{eq:composition-Lr-continuity}
 \sup_{t\in I}\|\mathcal C_{X(t)}^{\pm1}\|_{L^q\to L^q}<\infty,\qquad
 \lim_{t\to t_0}\|\mathcal C_{X(t)}^{\pm1}g
 -\mathcal C_{X(t_0)}^{\pm1}g\|_{L^q}=0.
\end{equation}
\end{lemma}

\begin{proof}
For $r=0$, \eqref{eq:sobolev-composition} is the change-of-variables formula;
for $r=1$ it follows from the weak chain rules
\begin{equation}\label{eq:pullback-chain-rule}
 \nabla^y(\mathcal C_Xf)=\mathcal C_X(\nabla f)dX,
 \qquad
 \nabla(\mathcal C_X^{-1}u)=\mathcal C_X^{-1}\big((\nabla^yu)\Kmat\big),
\end{equation}
together with the bounds on $\|dX\|_{L^\infty}$, $\|\Kmat\|_{L^\infty}$ and
$\det dX$; and interpolation gives $0<r<1$.  For $1<r\le s$, use the elementary
characterisation
\[
 \|\phi\|_{H^r}\simeq\|\phi\|_{H^{r-1}}+\|\nabla \phi\|_{H^{r-1}}
\]
and induct on $k(r)$.  In the first chain rule, apply the product estimate
\eqref{eq:tame-sp} at $\varsigma=r-1\le s-1$ to $dX=\mathrm{Id}+\Ufield$ after composing; in the second,
apply it to $\Kmat=\mathrm{Id}+(\Kmat-\mathrm{Id})$ before composing.
Each step lowers the Sobolev index by one and introduces one factor in the
displayed bound for $C_{\mathrm{comp}}(X,r)$; the derivatives of $dX$ and
$\Kmat$ arising in the Sobolev product are controlled by \eqref{eq:tame-sp}
through $\Ufield,\Kmat-\mathrm{Id}\in Y^{s-1,1}$.

For the time-dependent statement, the uniform operator bounds reduce the claim by
density to a fixed smooth compactly supported function.  The hypotheses imply
$X(t)^{\pm1}\to X(t_0)^{\pm1}$ locally uniformly and
$\Kmat(t)-\Kmat(t_0)\to0$ in $L^\infty\cap Y^{s-1,1}$.  Dominated
convergence proves the claim for $r=0$, the chain rules give $r=1$, and interpolation
gives $0<r<1$; by the product estimate \eqref{eq:tame-sp} and induction on $k(r)$ the case $1<r\le s$ follows.

For \eqref{eq:composition-Lr-continuity}, the operators
$\mathcal C_{X(t)}^{\pm1}$ are uniformly bounded on every $L^q$,
$1\le q<\infty$, because $J$ is bounded above and below. For fixed
$g\in C_c(\Rn)$, strong continuity follows from dominated convergence and the
assumed locally uniform convergence of $X(t)^{\pm1}$; density of $C_c(\Rn)$ in
$L^q$ and the uniform operator bounds do the rest.
\end{proof}


\subsection{Coefficient calculus}\label{ssec:coeffcalculus}

By Proposition~\ref{prop:tame-product}, $Y^{s-1,1}$ is a Banach algebra. We are now going to study its behaviour under analytic functions, which in the applications will be restrictions to a small ball in the space of matrices $M$ of $\det(\mathrm{Id}+M)$, its cofactors,
$\log\det(\mathrm{Id}+M)$, $(\mathrm{Id}+M)^{-1}$ and their products. Thus the following statement and \eqref{eq:coeff-lipschitz} in particular control $J-1$, $\Kmat-\mathrm{Id}$, the two
inverse factors in $\mathcal Q$, and $\varrho$ from one statement.

\begin{lemma}[coefficient calculus]\label{lem:coeff-calculus}
Let $\Psi$ be real-analytic near the origin in a finite-dimensional matrix space. There is
$\delta_\Psi=\delta_\Psi(\Psi,n,s,v_{\rm A})>0$ such that, if
$\|M_1\|_{Y^{s-1,1}}+\|M_2\|_{Y^{s-1,1}}\le\delta_\Psi$, then
\begin{equation*}\tag{P}\label{eq:coeff-lipschitz}
\|\Psi(M_1)-\Psi(M_2)\|_{Y^{s-1,1}}
\ \le\ C_\Psi\|M_1-M_2\|_{Y^{s-1,1}}.
\end{equation*}
In particular
$\|\Psi(M_1)-\Psi(0)\|_{Y^{s-1,1}}\le C_\Psi\|M_1\|_{Y^{s-1,1}}$.

In addition, suppose that the pointwise ranges of $M_1,M_2$ lie in a fixed
compact convex set $\mathscr K$ on a neighbourhood of which $\Psi$ is defined and
$C^1$.  Then, for every $1\le p\le\infty$,
\[
\|\Psi(M_1)-\Psi(M_2)\|_{L^p}\le C_{\Psi,\mathscr K}\|M_1-M_2\|_{L^p},
\]
and, separately, for every $\partial\in\{\partial_t,\partial_a\}$, $1\le a\le n$,
for which the indicated weak derivative of $M_1$ belongs to $L^1_{\rm loc}$,
\[
\partial\Psi(M_1)=\Psi'(M_1)\partial M_1\quad\text{in }\mathcal D'.
\]
Consequently, the same difference estimates hold with $C_tL^p$ in place of
$L^p$. If $M\in C_tY^{s-1,1}$, then
$\Psi(M)-\Psi(0)\in C_tY^{s-1,1}$; more generally, differences
$\Psi(M_1)-\Psi(M_2)$ are continuous there by
\eqref{eq:coeff-lipschitz}. This formulation is necessary because
$Y^{s-1,1}$ is nonunital when $\Psi(0)\ne0$.
\end{lemma}

\begin{proof}
Expand $\Psi$ at the origin.  Carefully rewriting each polynomial to highlight differences $M_1^j-M_2^j$ in
$\Psi(M_1)-\Psi(M_2)$ and using the $Y^{s-1,1}$ algebra estimate \eqref{eq:Y-product} gives a summable geometric series majorant as soon as
$\|M_1\|_{Y^{s-1,1}}+\|M_2\|_{Y^{s-1,1}}\le\delta_\Psi$, proving \eqref{eq:coeff-lipschitz}; take $M_2=0$ for the one-function estimate.

For the $L^p$ clause, $\Psi'$ is bounded on $\mathscr K$ and the segment joining
$M_2(y)$ to $M_1(y)$ stays in $\mathscr K$ by convexity, so the mean-value
theorem applies pointwise; no derivative hypothesis is used here.  The
chain rule is the Lipschitz chain rule obtained by mollification and a.e. convergence (rather than by density, which would fail for $p=\infty$).
\end{proof}

\subsection{The pressure Hessian}
\label{ssec:pressure-hessian}

Fix an orientation-preserving global bi-Lipschitz map $X$ and the associated
precomposition operator $\mathcal C_X$, metric $\gmetric$ and measure $\mu$ from
Section~\ref{sec:wave-hodge}.  The canonical vector field
$\mathcal P_{\gmetric}q$ along $X$ and its Hardy--Littlewood--Sobolev bound are given by
Lemma~\ref{lem:hodge-solvability}.  We now estimate its componentwise label
derivative, namely the Lagrangian pressure Hessian; this is the only (derivative) elliptic estimate the paper needs. 

The idea is that $\mathcal P_{\gmetric}$ has order $-1$, so $d\mathcal P_{\gmetric}$
should behave like an operator of order zero.  To make this rigorous one needs
only observe that, in label coordinates, the metric Hodge operator is a
principal but small-coefficient perturbation of the flat Hodge operator in the
labels; the smallness will be guaranteed by the parameter $\delta$ of the
fixed-point class $\mathcal U_{R,\delta}$ of \S\ref{sec:fixedpoint}.

\begin{lemma}[Hodge a priori estimates: general rhs]\label{lem:label-hodge}
Let $\Ufield=dw$ be exact, let $X=\mathrm{id}+w$, recall $\sigma=s-1>(n-1)/2$,
and put
\begin{equation}\label{eq:KMdef}
    \Kmat=\Kmat[\Ufield]:=(\mathrm{Id}+\Ufield)^{-1},\qquad
    \Mmat=\Mmat[\Ufield]:=\Kmat[\Ufield]-\mathrm{Id}.
\end{equation}
There is $\delta_{\rm H}>0$, depending only on the standing parameters, such that if
\begin{equation}\label{eq:hodge-smallness}
 \|\Ufield\|_{Y^{\sigma,1}}\le\delta_{\rm H},
\end{equation}
then $X$ is an orientation-preserving global bi-Lipschitz map, and there is a
bounded linear operator $d\mathcal P_{\gmetric}$ on $Y^{\sigma,\preg}$,
$0\le\preg\le1$, with
\begin{equation}\label{eq:elliptic-Y}
 \|d\mathcal P_{\gmetric}q\|_{Y^{\sigma,\preg}}
 \le C_{\rm H}\|q\|_{Y^{\sigma,\preg}},
 \qquad 0\le\preg\le1.
\end{equation}
In particular, since $Y^{\sigma,0}=H^\sigma$,
\begin{equation}\label{eq:elliptic-Hsigma}
 \|d\mathcal P_{\gmetric}q\|_{H^\sigma}
 \le C_{\rm H}\|q\|_{H^\sigma}.
\end{equation}
The constant is uniform on the ball \eqref{eq:hodge-smallness}.  If moreover
$q$ belongs to $L^p(\mu)$ for some $1<p<n$, which is automatic with $p=2$ when
$n\ge3$, then
\begin{equation}\label{eq:elliptic-identification}
 (d\mathcal P_{\gmetric}q)_a^i
 =\partial_a(\mathcal P_{\gmetric}q)^i
 \qquad\text{in }\mathcal D',
\end{equation}
that is, on this class $d\mathcal P_{\gmetric}$ is the distributional label
derivative of the canonical vector field along $X$ defined by \eqref{eq:pressure-gradient}, as the
notation indicates.  For $n\ge3$ this identifies the operator completely; for
$n=2$ the canonical vector field along $X$ need not be defined for a general
$q\in Y^{\sigma,\preg}$, and $d\mathcal P_{\gmetric}$ is its bounded extension.
The operator depends only on $\Ufield$, not on the spatially constant part of
$w$.
\end{lemma}

\begin{proof}
We construct a candidate for $d\mathcal P_{\gmetric}q$ at the level of the
differentiated Hodge system, by means of Riesz-transform type operators, and
then identify it with the distributional derivative of the canonical vector field
$\mathcal P_{\gmetric}q$ along $X$ from Lemma~\ref{lem:hodge-solvability}.  The threshold
$\delta_{\rm H}$ is shrunk finitely many times below.

\emph{Geometry and coefficients.}
By the embedding \eqref{eq:embed-frac} at $(\sreg,\preg)=(\sigma,1)$ we may
assume $\|\Ufield\|_{L^\infty}\le\tfrac12$.  Then
$w$ has a Lipschitz representative with constant $\tfrac12$, so
\[
 \tfrac12|y-z|\le|X(y)-X(z)|\le\tfrac32|y-z|,
\]
and for every $x$ the contraction $y\mapsto x-w(y)$ of the complete space
$\Rn$ produces the unique solution of $X(y)=x$; hence $X$ is a global
bi-Lipschitz homeomorphism.  Since $\det(\mathrm{Id}+\theta\Ufield)$ vanishes
nowhere for $0\le\theta\le1$, the Jacobian retains the sign of the identity,
so $X$ is orientation-preserving with $2^{-n}\le J\le(3/2)^n$.  Finally, Lemma~\ref{lem:coeff-calculus}, applied to the real-analytic map
$\Psi(M):=(\mathrm{Id}+M)^{-1}-\mathrm{Id}$, gives the first of
\begin{equation}
    \|\Mmat\|_{Y^{\sigma,1}}\le C_\Psi\|\Ufield\|_{Y^{\sigma,1}},
 \qquad
 \|\Kmat\|_{L^\infty;\mathrm{op}}\le2 ,
\end{equation}
the second being the Neumann series bound
$\|(\mathrm{Id}+\Ufield)^{-1}\|_{L^\infty;\mathrm{op}}
\le(1-\|\Ufield\|_{L^\infty;\mathrm{op}})^{-1}$.

\emph{The differentiated resolvent.}
From the Cartesian calculus for vector fields along $X$ in \S\ref{sec:wave-hodge}, we deduce that the metric Hodge system \eqref{eq:invdiv} is exactly
\begin{equation}\label{eq:label-hodge-perturbation}
 \begin{cases}
  \Div_{y}F
    &=q-\Mmat^a_i\partial_aF^i,\\
  (\curl_{y}F)^{ij}
    &=-\Mmat^a_i\partial_aF^j
      +\Mmat^a_j\partial_aF^i.
 \end{cases}
\end{equation}
Here $\Div_y$ and $\curl_y$ denote $\Div_{\gmetric}$ and $\curl_{\gmetric}$ with
$\Kmat^a_i$ replaced by $\delta^a_i$, that is
$\Div_yF=\partial_iF^i$ and $(\curl_yF)^{ij}=\partial_iF^j-\partial_jF^i$, label and Cartesian indices being identified in these operators and
throughout the flat construction that follows, so that a label derivative and
the Fourier variable may carry a Cartesian index; this is used nowhere else.
Note that $\curl_y$ antisymmetrises the Cartesian index of $F$ against the
derivative index, whereas the rowwise label curl $\Curl_y$ of
\eqref{eq:Curl} antisymmetrises the two label indices of an
$\Rn$-valued $1$-form and leaves its Cartesian index untouched.  This makes
precise the perturbation structure announced at the start of the subsection,
\[
 (\Div_{\gmetric},\curl_{\gmetric})
 = (\Div_{y},\curl_{y})+(\Kmat-\mathrm{Id})d .
\]
The remainder is of principal order, rather than a lower-order Christoffel
term, but its coefficient is small.

Motivated by
\begin{equation}\label{eq:motivation}
    \Delta_yF^j=\partial_i^2F^j=\partial_i(\partial_iF^j-\partial_jF^i)+\partial_i\partial_jF^i=\partial_i(\curl_y F)^{ij}+\partial_j\Div_y F
\end{equation}
giving
\[
\partial_aF^j =\partial_a\Delta^{-1}_y\left[\partial_i(\curl_y F)^{ij}+\partial_j\Div_y F\right],
\]
we define, for a scalar $r$ and an antisymmetric $2$-tensor $c=(c^{ij})$, the flat
Hodge reconstruction operator
\begin{equation}\label{eq:Hodgereconstruction}
    \big(\mathscr H(r,c)\big)_a^j
 :=\partial_a\Delta_y^{-1}\big(\partial_jr+\partial_ic^{ij}\big).
\end{equation}
Indeed, if $r=\Div_{y}F$ and
$c^{ij}=\partial_iF^j-\partial_jF^i$, then \eqref{eq:motivation} shows that
$dF=\mathscr H(r,c)$.  Recalling $\Hess=dF$, the system
\eqref{eq:label-hodge-perturbation} becomes the resolvent identity
\begin{equation}\label{eq:pressure-Hessian}
 \Hess=\mathscr T_0q+\mathscr T_{\Mmat}\Hess,
\end{equation}
where
\[
 \mathscr T_{\Mmat}\Hess
 :=\mathscr H\big(-\Mmat^a_i\Hess_a^i,
       -\Mmat^a_i\Hess_a^j+\Mmat^a_j\Hess_a^i\big), \qquad \mathscr T_0q:=\mathscr H(q,0).
\]
The entries of $\mathscr H$ are Fourier multipliers with the
homogeneous symbols $\xi_a\xi_j/|\xi|^2$, with the value at $\xi=0$ fixed
to be zero. Equivalently, they are compositions of Riesz transforms and are
bounded by one at the symbol level. Since every
$Y^{\sigma,\preg}$ is a weighted $L^2$ space
on the Fourier side, with weight depending only on $\xi$, each entry is
bounded on every $Y^{\sigma,\preg}$ with norm at most one.  By \eqref{eq:Y-multiplier-interp} with $g=\Mmat$, the operator norm of
$\mathscr T_{\Mmat}$ on $Y^{\sigma,\preg}$
is at most $C_0\|\Mmat\|_{Y^{\sigma,1}}$, with $C_0=C_0(n,s,v_{\rm A})$
independent of $\preg$.  Shrinking $\delta_{\rm H}$ so that
$C_0C_\Psi\delta_{\rm H}\le\tfrac12$, the operator
$\mathrm{Id}-\mathscr T_{\Mmat}$ is inverted by its Neumann series, uniformly
for $0\le\preg\le1$; in particular \eqref{eq:pressure-Hessian} has at most one
solution in $Y^{\sigma,\preg}$.  Define
\begin{equation}\label{eq:neumann-hodge-derivative}
 d\mathcal P_{\gmetric}q
 :=(\mathrm{Id}-\mathscr T_{\Mmat})^{-1}\mathscr T_0q .
\end{equation}
The preceding bounds prove \eqref{eq:elliptic-Y}--\eqref{eq:elliptic-Hsigma},
and \eqref{eq:pressure-Hessian} contains only $\Ufield$ and $\Kmat$, so the
operator does not see the spatially constant part of $w$.

\emph{Identification.}
It remains to prove \eqref{eq:elliptic-identification} when, in addition,
$q\in L^p(\mu)$ for some $1<p<n$. Choose a standard nonnegative unit-mass
mollifier $\rho_\ell$ and set
$w_\ell:=w*\rho_\ell$, so that each $\Ufield_\ell:=dw_\ell=\Ufield*\rho_\ell$
is smooth and exact, satisfies \eqref{eq:hodge-smallness} and
$\|\Ufield_\ell\|_{L^\infty}\le\tfrac12$, and $\Ufield_\ell\to\Ufield$ in
$Y^{\sigma,1}$; choose also smooth, compactly supported $q_\ell\to q$ in
$H^\sigma\cap L^p$.  Lemma~\ref{lem:coeff-calculus} gives
$\Kmat_\ell\to\Kmat$, $\Mmat_\ell\to\Mmat$ and $J_\ell\to J$ in $L^\infty$,
with $\Mmat_\ell\to\Mmat$ also in $Y^{\sigma,1}$; the maps
$X_\ell=\mathrm{id}+w_\ell$ are orientation-preserving and uniformly
bi-Lipschitz.

For these smooth pairs the canonical vector field
$F_\ell:=\mathcal P_{\gmetric_\ell}q_\ell$ along $X_\ell$ from
Lemma~\ref{lem:hodge-solvability} is smooth.  Put
$u_\ell:=\nabla_x\Delta_x^{-1}\mathcal C_{X_\ell}^{-1}q_\ell$, so that
$F_\ell=\mathcal C_{X_\ell}u_\ell$ by \eqref{eq:pressure-gradient}.  Applying
the first identity of \eqref{eq:pullback-chain-rule} componentwise gives
$dF_\ell=\mathcal C_{X_\ell}(\nabla_xu_\ell) dX_\ell$, that is
\[
 \Hess_\ell:=dF_\ell
 =\mathcal C_{X_\ell}\big(\nabla_x^2\Delta_x^{-1}\mathcal C_{X_\ell}^{-1}q_\ell\big) dX_\ell ,
\]
which belongs to $H^\sigma$ by Lemma~\ref{lem:composition} and
\eqref{eq:tame-sp}.
In this class the computation \eqref{eq:motivation} is justified, so
$\Hess_\ell$ solves \eqref{eq:pressure-Hessian} with coefficient $\Mmat_\ell$
and datum $q_\ell$; by the uniqueness just established,
$\Hess_\ell=d\mathcal P_{\gmetric_\ell}q_\ell$.  The uniform resolvent bound
and \eqref{eq:canonical-hls} give, after extraction,
\[
 \Hess_\ell\rightharpoonup \overline{\Hess}\quad\hbox{in }H^\sigma,
 \qquad
 F_\ell\rightharpoonup F\quad\hbox{in }L^{p^*} .
\]
Passing to the limit in \eqref{eq:pressure-Hessian}, with the coefficient
converging strongly in $Y^{\sigma,1}$ and the solution weakly in $H^\sigma$,
shows that $\overline{\Hess}=d\mathcal P_{\gmetric}q$, while
$\Hess_\ell=dF_\ell$ gives $\overline{\Hess}=dF$ in distributions.

It remains to identify $F$. For each $\ell$, the conservative weak identities
in Lemma~\ref{lem:weak-pullback}(ii) hold for $F_\ell$, with divergence datum
$q_\ell$ and zero curl. Since $F_\ell\rightharpoonup F$ in $L^{p^*}$,
$q_\ell\to q$ in $L^p$, and $J_\ell\Kmat_\ell\to J\Kmat$ and
$J_\ell\to J$ in $L^\infty$, these identities pass directly to the limit
against every compactly supported smooth test function.
Lemma~\ref{lem:weak-pullback}(ii) therefore shows that $F$ solves the metric
Hodge system in the conjugated distributional sense of
Lemma~\ref{lem:hodge-solvability}. The uniform Jacobian bounds give
$F\in L^{p^*}(\mu)$, and uniqueness identifies $F=\mathcal P_{\gmetric}q$.
\end{proof}

In the remainder of this subsection $I=[0,T]$, and $\Ufield,\Vfield$ denote
spatially exact $\Rn$-valued $1$-forms in $H^{\sigma,1}_{\breg}(\Rst)$
satisfying
\begin{equation}\label{eq:admissible}
 \|\Ufield\|_{H^{\sigma,1}_{\breg}}\le R,
 \qquad
 \sup_{t\in I}\|\Ufield(t)\|_{Y^{\sigma,1}}\le\delta ,
\end{equation}
and likewise for $\Vfield$.  Here $\delta>0$ is small enough, depending only
on the standing parameters, that the following conditions are met. First, every $\Ufield(t)$ satisfies
\eqref{eq:hodge-smallness}, hence $\|\Ufield(t)\|_{L^\infty}\le\frac12$ by
\eqref{eq:embed-frac}. Second, Lemma~\ref{lem:coeff-calculus} applies to pairs of such coefficients; the
explicit choice \eqref{eq:delta-choice} is made in \S\ref{sec:fixedpoint}.

All time norms are taken over $I$, and $C_\delta$, $C_{\delta,\rpar}$ denote
constants depending only on $\delta$, $\rpar$ and the standing parameters,
never on $R$, on $T$ or on the fields.  We abbreviate, repeated indices being
summed,
\[
 N_{\Ufield,ab}^{ij}:=Q_0(\Ufield_a^j,\Ufield_b^i),
 \qquad
 \Ccoef_{\Ufield,ij}^{ab}:=\Kmat[\Ufield]^a_i\Kmat[\Ufield]^b_j,
 \qquad
 \widetilde{\Ccoef}_{\Ufield,ij}^{ab}
 :=\Ccoef_{\Ufield,ij}^{ab}-\delta_i^a\delta_j^b,
 \qquad
 \mathcal Q[\Ufield]=-\Ccoef_{\Ufield}:N_{\Ufield},
\]
and write $D:=\|\Ufield-\Vfield\|_{H^{\sigma,1}_{\breg}}$,
$\Hess_{\Ufield}:=dF[\Ufield]$.

\begin{lemma}[Hodge a priori estimates: null-form rhs]\label{lem:pressure}
Let $\Ufield,\Vfield$ be as above, recall $\sigma=s-1>(n-1)/2$, and let
$F[\Ufield]$ be the canonical Hodge representative \eqref{eq:FU}.  Then $F[\Ufield]\in C(I;L^4)$, with
\begin{equation}\label{eq:timecontL4}
 \|F[\Ufield]\|_{L^\infty(I;L^4)}\le C_\delta R^2 ,
\end{equation}
and, for every $\rpar\in[0,1]$ such that
$\Ufield\in H^{\sigma,1+\rpar}_{\breg}$,
\begin{equation}\label{eq:pressure-anisotropic}
 \|dF[\Ufield]\|_{L^2(I;Y^{\sigma,\rpar})}
 \le C_{\delta,\rpar} 
 \|\Ufield\|_{H^{\sigma,1}_{\breg}} 
 \|\Ufield\|_{H^{\sigma,1+\rpar}_{\breg}} .
\end{equation}
In particular at $\rpar=0$,
\begin{equation}\label{eq:pressure-isotropic}
 \|dF[\Ufield]\|_{L^2(I;H^\sigma)}\le C_\delta R^2 .
\end{equation}
Finally,
\begin{equation}\label{eq:pressure-lipschitz}
 \|dF[\Ufield]-dF[\Vfield]\|_{L^2(I;H^\sigma)}
 \le C_\delta R(1+R) D .
\end{equation}
\end{lemma}

\begin{proof}
Everything below depends only on the values of the fields on $I$.  We first
record three facts.

\emph{(a) Time regularity.}  Corollary~\ref{cor:trace-Y} gives
$\Ufield\in C(I;Y^{\sigma,1})$ and
$\partial_t\Ufield,\mathcal B\Ufield\in C(I;H^\sigma)$, the latter two with
norm at most $C_{\breg}R$.  Since
$\Psi(Z):=(\mathrm{Id}+Z)^{-1}-\mathrm{Id}$ is Lipschitz on the $\delta$-ball
of $Y^{\sigma,1}$ by \eqref{eq:coeff-lipschitz},
$\Mmat[\Ufield]\in C(I;Y^{\sigma,1})$ with norm at most $C_\delta$, and
likewise $\widetilde{\Ccoef}_{\Ufield}\in C(I;Y^{\sigma,1})$; moreover
\begin{equation}\label{eq:coeff-difference}
 \|\Mmat_{\Ufield}-\Mmat_{\Vfield}\|_{L^\infty_tY^{\sigma,1}}
 +\|\Ccoef_{\Ufield}-\Ccoef_{\Vfield}\|_{L^\infty_tY^{\sigma,1}}
 \le C_\delta D .
\end{equation}
In particular $t\mapsto(\mathrm{Id}-\mathscr T_{\Mmat[\Ufield](t)})^{-1}$ is
norm-continuous, so every time-dependent quantity below is (strongly)
measurable.

\emph{(b) The coefficient step.} Write
$\Ccoef_{\Ufield}:N=(\delta_i^a\delta_j^b):N
+\widetilde{\Ccoef}_{\Ufield}:N$. The constant contraction is bounded directly,
whereas the remainder is controlled by \eqref{eq:Y-multiplier-interp}. Together
with (a) and \eqref{eq:coeff-difference}, this gives, uniformly in
$\rpar\in[0,1]$ and for every tensor $N$,
\begin{equation}\label{eq:coefficient-step}
 \|\Ccoef_{\Ufield}:N\|_{L^2_tY^{\sigma,\rpar}}
 \le C_\delta\|N\|_{L^2_tY^{\sigma,\rpar}},
 \qquad
 \|(\Ccoef_{\Ufield}-\Ccoef_{\Vfield}):N\|_{L^2_tY^{\sigma,\rpar}}
 \le C_\delta D \|N\|_{L^2_tY^{\sigma,\rpar}} .
\end{equation}

\emph{(c) The null-form step.}  Applying the bilinear estimate
\eqref{eq:null-form} to each pair of components,
\begin{equation}\label{eq:nullform-step}
 \|N_{\Ufield}\|_{L^2_tY^{\sigma,\rpar}}
 \le C\|\Ufield\|_{H^{\sigma,1}_{\breg}}
      \|\Ufield\|_{H^{\sigma,1+\rpar}_{\breg}},
 \qquad
 \|N_{\Ufield}-N_{\Vfield}\|_{L^2_tH^\sigma}\le CRD ,
\end{equation}
the second because bilinearity gives
$N_{\Ufield}-N_{\Vfield}
=Q_0(\Ufield-\Vfield,\Ufield)+Q_0(\Vfield,\Ufield-\Vfield)$ componentwise.

\emph{Fixed-time integrability and \eqref{eq:timecontL4}.}  Since
\[
 \sigma>\frac{n-1}2\ \ge\ \frac{3n-4}{8}=\frac n2-\frac n{2r_{\mathrm{HLS}}}
 \qquad(n\ge2),
\]
Sobolev embedding gives $H^\sigma\hookrightarrow L^{2r_{\mathrm{HLS}}}$, with
$r_{\mathrm{HLS}}$ as in \eqref{eq:rHLS}.  By (a), each factor of
$N_{\Ufield}=\partial_t\Ufield \partial_t\Ufield
-\mathcal B\Ufield \mathcal B\Ufield$ lies in
$C(I;L^{2r_{\mathrm{HLS}}})$ with norm at most $CR$, so H\"older gives
$N_{\Ufield}\in C(I;L^{r_{\mathrm{HLS}}})$ with norm at most $CR^2$; since
$\Ccoef_{\Ufield}\in C(I;L^\infty)$ with norm at most $C_\delta$, by (a) and
\eqref{eq:embed-frac},
\begin{equation}\label{eq:datum-fixed-time}
 \mathcal Q[\Ufield]\in C(I;L^{r_{\mathrm{HLS}}}),
 \qquad
 \sup_{t\in I}\|\mathcal Q[\Ufield](t)\|_{L^{r_{\mathrm{HLS}}}}
 \le C_\delta R^2 .
\end{equation}
In particular, the canonical vector field along $X(t)$ given by
$F[\Ufield](t)=\mathcal C_{X(t)}\nabla\Delta_x^{-1}
\mathcal C_{X(t)}^{-1}\mathcal Q[\Ufield](t)$ is defined at every time, since
$1<r_{\mathrm{HLS}}<n$ and $r_{\mathrm{HLS}}^*=4$. Normalize a primitive of
$\Ufield(t)$ by zero mean on the unit ball. The homotopy formula used in
\S \ref{sec:fixedpoint}, together with
$\Ufield\in C(I;Y^{\sigma,1})\hookrightarrow C(I;C_b)$, shows that $X(t)$
varies locally uniformly; the uniform bi-Lipschitz bounds give the same for
$X(t)^{-1}$. By \eqref{eq:composition-Lr-continuity}, the operators
$\mathcal C_{X(t)}^{\pm1}$ are uniformly bounded on $L^{r_{\mathrm{HLS}}}$ and
on $L^4$, and strongly continuous in $t$.  Since $\nabla\Delta_x^{-1}$ is
bounded from $L^{r_{\mathrm{HLS}}}$ to $L^4$ by \eqref{eq:flat-hls} at
$p=r_{\mathrm{HLS}}$, the composition is uniformly bounded and strongly
continuous in $t$, and \eqref{eq:datum-fixed-time} gives
$F[\Ufield]\in C(I;L^4(dy))$ together with \eqref{eq:timecontL4}. The
fixed-time HLS estimate also gives a uniform $L^4(\mu[\Ufield](t))$ bound, and
the two bounds are equivalent uniformly because the Jacobians are bounded
above and below.

\emph{Proof of \eqref{eq:pressure-anisotropic} and \eqref{eq:pressure-isotropic}.}
The continuous representative from \eqref{eq:datum-fixed-time} belongs to
$L^{r_{\mathrm{HLS}}}$ at every time, while (b)--(c) give
$\mathcal Q[\Ufield](t)\in Y^{\sigma,\rpar}$ for almost every time. Thus
\eqref{eq:elliptic-identification} applies for almost every $t$ and gives
$dF[\Ufield](t)=d\mathcal P_{\gmetric(t)}\mathcal Q[\Ufield](t)$. Combining
(c) with the first estimate of \eqref{eq:coefficient-step} bounds
$\|\mathcal Q[\Ufield]\|_{L^2_tY^{\sigma,\rpar}}$ by the right-hand side of
\eqref{eq:pressure-anisotropic}; squaring the fixed-time estimate
\eqref{eq:elliptic-Y} at $\preg=\rpar$ and integrating in $t$ converts this
into \eqref{eq:pressure-anisotropic}.  Taking $\rpar=0$ and
$\|\Ufield\|_{H^{\sigma,1}_{\breg}}\le R$ gives \eqref{eq:pressure-isotropic}.

\emph{Proof of \eqref{eq:pressure-lipschitz}.}  The datum difference
decomposes as
\begin{equation}\label{eq:datum-difference-decomposition}
 \mathcal Q[\Ufield]-\mathcal Q[\Vfield]
 =-(\Ccoef_{\Ufield}-\Ccoef_{\Vfield}):N_{\Vfield}
   -\Ccoef_{\Ufield}:(N_{\Ufield}-N_{\Vfield}) .
\end{equation}
By \eqref{eq:coefficient-step} and \eqref{eq:nullform-step} at $\rpar=0$, the
first term costs $C_\delta R^2D$, because the geometry changes, and the
second $C_\delta RD$, because one factor of the quadratic null form is
replaced by the difference; hence
\begin{equation}\label{eq:datum-lipschitz}
 \|\mathcal Q[\Ufield]-\mathcal Q[\Vfield]\|_{L^2_tH^\sigma}
 \le C_\delta R(1+R)D .
\end{equation}
Subtracting the two resolvent identities \eqref{eq:pressure-Hessian},
\[
 (\mathrm{Id}-\mathscr T_{\Mmat_{\Ufield}})
 \big(\Hess_{\Ufield}-\Hess_{\Vfield}\big)
 =\mathscr T_0\big(\mathcal Q[\Ufield]-\mathcal Q[\Vfield]\big)
  +\mathscr T_{\Mmat_{\Ufield}-\Mmat_{\Vfield}}\Hess_{\Vfield} .
\]
The inverse on the left has norm at most $2$, by the Neumann argument of
Lemma~\ref{lem:label-hodge}.  The first term on the right is bounded by
$C_\delta R(1+R)D$ through \eqref{eq:datum-lipschitz}, and the second, by the
linearity of $\mathscr T_{\Mmat}$ in its coefficient together with
\eqref{eq:coeff-difference} and \eqref{eq:pressure-isotropic} applied to
$\Vfield$, by
\[
 C\|\Mmat_{\Ufield}-\Mmat_{\Vfield}\|_{L^\infty_tY^{\sigma,1}}
   \|\Hess_{\Vfield}\|_{L^2_tH^\sigma}
 \le C_\delta R^2D .
\]
Combining the two contributions proves \eqref{eq:pressure-lipschitz}.
\end{proof}

\section{The Differentiated Fixed Point}
\label{sec:fixedpoint}

We combine the linear, null-form and elliptic estimates to produce a
spatially exact spacetime field $\Ufield\in H^{s-1,1}_{\breg}$ solving the differentiated wave-Hodge system. The fixed-time smallness $\delta$ built into the class
supplies the coefficient and Hodge estimates, while short time and the datum
modulus $\omega_{v_0}$ yield this smallness for arbitrary-size data $R$.  The displacement $w$ is
reconstructed by the Duhamel formula only after the differentiated fixed point $\Ufield$ is found, and it will be such that $dw=\Ufield$ on the interval $[0,T_*]$ of definition.

\subsection{Preliminaries} We first introduce all the needed objects.\\

\textbf{The class.} Fix $0<T\le1$ and write $I:=[0,T]$. Let $\mathcal{Y}$ be the Hilbert space
\[
 \mathcal Y:=H^{s-1,1}_{\breg}(\Rst)
\]
and for $v_0\in H^{s}(\Rn)$ divergence-free, of arbitrary size, we let
$C_{\rm lin}$ be the constant of \eqref{eq:linear} at $(\sreg,\preg)=(s-1,1)$, and put
\begin{equation}\label{eq:Rdefn}
    R:=2C_{\rm lin}\|v_0\|_{H^s}.
\end{equation}
For an $\Rn$-valued $1$-form $\Ufield$, write
\begin{equation}\label{eq:Curl}
    (\Curl_{y}\Ufield)^j_{ab}
 :=\partial_a\Ufield_b^j-\partial_b\Ufield_a^j, 
\end{equation}

we then let the fixed-point class be:
\begin{equation}\label{eq:fixedpoint-class}
  \mathcal U_{R,\delta} := \left\{ \Ufield \in \mathcal Y : 
  \begin{aligned}
    &\Curl_{y}\Ufield = 0, \quad (\Ufield,\partial_t\Ufield)(0) = (0,\nabla v_0), \\
    &\|\Ufield\|_{\mathcal Y} \le R, \quad \sup_{t\in I}\|\Ufield(t)\|_{Y^{s-1,1}} \le \delta
  \end{aligned}
  \right\}.
\end{equation}
Note that the exactness and the trace conditions are closed. An element of $\mathcal U_{R,\delta}$ should be thought of as a `trial field' for the differentiated wave Hodge system before a fixed point is found.  Moreover, at $t=0$ the identity $X_\Ufield(0,y)=y$ for any trial flow map (see \eqref{eq:trial-bilipschitz} below) identifies Eulerian initial data with label fields, namely $x=y$, and we can write $(\nabla v_0)_a^i=\partial_av_0^i$ unequivocally.\\

\textbf{Reconstruction of the displacement $\bar w_\Ufield$ and map $X_\Ufield$ for a trial field $\Ufield$.} We now explain how from the exactness condition of a trial field $\Ufield$ we can go back to a (primitive) displacement $\bar w_\Ufield$ and map $X_\Ufield$; this is important because these objects will enter the definition of the fixed-point map $\mathcal{M}$, which will be defined next.  By the distributional Poincar\'e
lemma, see for example
\cite[Theorem~2.1]{Mardare2008}, an exact $\Ufield$ admits a primitive
$\bar w_{\Ufield}$ with $d\bar w_{\Ufield}=\Ufield$.  Here the primitive can
be written down explicitly: by Corollary~\ref{cor:trace-Y},
$\Ufield\in C(\Real;Y^{s-1,1})$. For every compactly supported smooth spatial
test tensor $\Phi$, the map
\[
 t\longmapsto\langle\Curl_y\Ufield(t),\Phi\rangle
\]
is continuous. Its pairing against every compactly supported smooth time test
vanishes because $\Curl_y\Ufield=0$ in spacetime distributions. It therefore
vanishes identically. Thus $\Curl_y\Ufield(t)=0$ in $\mathcal D'(\Rn)$ for
every $t$, so the fixed-time Poincar\'e lemma and the homotopy formula below
apply at every time. The proof of
Lemma~\ref{lem:Linftyembed} shows that elements of $Y^{s-1,1}$ have integrable
Fourier transforms, so that $Y^{s-1,1}\hookrightarrow C_b(\Rn)$; hence
$\Ufield$ is jointly continuous and bounded on $I\times\Rn$, and the homotopy
formula
\[
 \bar w_{\Ufield}^{ j}(t,y):=\int_0^1\Ufield_a^j(t,\theta y) y^a d\theta
\]
defines a jointly continuous function.  It is a primitive: mollifying
$\Ufield(t)$ in $y$ yields smooth closed forms $\Ufield_\varepsilon$, for which
the formula gives a primitive $\bar w_\varepsilon$ classically,
$d\bar w_\varepsilon=\Ufield_\varepsilon$.  As $\varepsilon\to0$ both
$\Ufield_\varepsilon\to\Ufield(t)$ and $\bar w_\varepsilon\to\bar
w_{\Ufield}(t)$ locally uniformly, so this identity passes to the limit in
$\mathcal D'$ and $d\bar w_{\Ufield}=\Ufield$.  Since $\Ufield(t)\in L^\infty_y$,
$\bar w_{\Ufield}(t)$ is Lipschitz in $y$ for every $t$.
We then fix the remaining spatial translation gauge by subtracting the mean over the
unit ball, so that
\[
 |B_1|^{-1}\int_{B_1}\bar w_{\Ufield}(t,y) dy=0.
\]
For $t\in I$ the anisotropic embedding \eqref{eq:embed-frac} and the definition
of the class give
\[
 L_t:=\|d\bar w_{\Ufield}(t)\|_{L^\infty}
 \le C_{\rm e}\|\Ufield(t)\|_{Y^{s-1,1}}\le C_{\rm e}\delta\le\tfrac12,
\]
by the definition of $\delta$ which for convenience will be given later in \eqref{eq:delta-choice}.  Off $I$ the class imposes no
fixed-time smallness, and the discussion below is not claimed there.
Consequently, for all $t\in I$ and all $y,z\in\Rn$,
\begin{equation}\label{eq:trial-bilipschitz}
 (1-L_t)|y-z|
 \le|X_{\Ufield}(t,y)-X_{\Ufield}(t,z)|
 \le(1+L_t)|y-z|,
 \qquad X_{\Ufield}:=\mathrm{id}+\bar w_{\Ufield}.
\end{equation}
This proves injectivity.  Surjectivity follows separately: for each $x\in\Rn$, the map
$y\mapsto x-\bar w_{\Ufield}(t,y)$ is a contraction of the complete space
$\Rn$ and therefore has a unique fixed point, which is precisely the unique
solution of $X_{\Ufield}(t,y)=x$.  Thus $X_{\Ufield}(t,\cdot)$ is a global
bi-Lipschitz homeomorphism; its inverse has Lipschitz constant at most
$(1-L_t)^{-1}$, and the singular values of $dX_{\Ufield}$ lie between
$1-L_t$ and $1+L_t$.  The path
$\mathrm{Id}+\theta\Ufield$, $0\le\theta\le1$, remains invertible, so its
determinant has the positive sign of the identity.  In particular,
\[
 0<(1-L_t)^n\le J[\Ufield](t,y)\le(1+L_t)^n
\]
almost everywhere and the map is orientation preserving. Any other primitive
differs by a spatial translation. Such a translation leaves $\Ufield$,
$\Kmat$, $\gmetric$, $\mu$, $\mathcal Q[\Ufield]$ and the operators $E_i$
unchanged. The $L^4(\mu)$ uniqueness in Lemma~\ref{lem:hodge-solvability}
therefore shows that the canonical field $F[\Ufield]$ is unchanged as well.
Hence the construction below is independent of the chosen normalisation,
denoted by the overbar in $\bar w_\Ufield$.\\

\textbf{The map.} For $\Ufield\in\mathcal U_{R,\delta}$, let $F[\Ufield]$ be the canonical
$L^4$ representative on $I$ characterised by the Hodge system
\eqref{eq:force-hodge} for the metric $\gmetric[\Ufield]$, namely \eqref{eq:FU}. By
\eqref{eq:pressure-gradient}, it is given by
\[
 F[\Ufield]
 =\mathcal C_{X_{\Ufield}}\nabla\Delta_{x}^{-1}\mathcal C_{X_{\Ufield}}^{-1}\mathcal Q[\Ufield],
\]
here $X_\Ufield$ was constructed above in \eqref{eq:trial-bilipschitz} and Lemma~\ref{lem:pressure} provides the needed bounds on this forcing.  With this at hand, we define the time-localised map
\begin{equation}\label{eq:differential-map}
 (\mathcal M\Ufield)(t)
 :=\chi(t)\Sprop(t)\nabla v_0
 -\chi(t/T)\int_0^t\Sprop(t-t')
       \mathbf 1_I(t')dF[\Ufield](t') dt'.
\end{equation}
Since $\chi$ fixed in \S \ref{ssec:linear} satisfies $\chi\equiv1$ on $[-1,1]$ and $T\le1$, both cutoffs equal one on $I$, and
the field defined by $\Vfield(t):=(\mathcal M\Ufield)(t)$ solves
\[
    \Box\Vfield=-dF[\Ufield] \quad \text{on} \quad I \quad \text{with data} \quad (0,\nabla v_0).
\]
In particular, a fixed point of $\mathcal{M}$ solves the differentiated wave
equation associated with \eqref{eq:auxiliary-system}; reconstruction of $w$ and
the identity $dw=\Ufield$ are proved below.

We now show some preliminary bounds and make some observations. Firstly, since both terms in \eqref{eq:differential-map} are spatial differentials, we obtain
$\Curl_y(\mathcal M\Ufield)=0$ and thus $\mathcal{M}$ preserves the exactness condition of $\mathcal U_{R,\delta}$. Moreover, for $G\in L^2(I;H^{s-1})$, the time-localised Duhamel term satisfies
\begin{equation}\label{eq:cutoff-duhamel}
\begin{aligned}
 \left\|\chi(t/T)\int_0^t\Sprop(t-t')
       \mathbf 1_I(t')G(t') dt'\right\|_{\mathcal Y}
 &\le C_{\rm lin}T^\varepsilon\|G\|_{L^2_tH^{s-1}},\\
 \sup_{t\in I}\left\|\int_0^t\Sprop(t-t')G(t') dt'\right\|_{Y^{s-1,1}}
 &\le CT^{1/2}\|G\|_{L^2_tH^{s-1}}.
\end{aligned}
\end{equation}
To see this, we first note that the free propagator \eqref{eq:propagator} satisfies the elementary scalar bound
\begin{equation}\label{eq:prop-gain}
\langle\lambda\rangle \Big|\frac{\sin(t\lambda)}{\lambda}\Big|\ \le\
\sqrt2 \min\big(1,\ |t| \langle\lambda\rangle\big),\qquad |t|\le1 ,
\end{equation}
obtained from $|\sin(t\lambda)|\le\min(1,|t\lambda|)$ in the regions
$|\lambda|\ge1$ and $|\lambda|\le1$. Then, for the first estimate, extend $G$ by zero outside $I$ and apply Theorem~\ref{thm:linear} with zero Cauchy data and
$(\sreg,\preg)=(s-1,1)$.  This is where time localisation gives the factor
$T^\varepsilon=T^{1-\breg}$.  For the second, we use the scalar multiplier estimate
\[
 \|\Sprop(r)f\|_{Y^{s-1,1}}
 \le C\|f\|_{H^{s-1}},\qquad |r|\le1,
\]
which follows directly from \eqref{eq:prop-gain} and the definition \eqref{eq:Y}, together with  Minkowski and Cauchy--Schwarz,
uniformly for $0\le t\le T$, to bound
\[
 \left\|\int_0^t\Sprop(t-t')G(t') dt'\right\|_{Y^{s-1,1}}
 \le C\int_0^t\|G(t')\|_{H^{s-1}} dt'
 \le CT^{1/2}\|G\|_{L^2(I;H^{s-1})}.
\]
\\

\textbf{The smallness condition $\delta$.} We now define $\delta$ and explain why this condition comes for free in our setup. This was used already to ensure that $X_\Ufield=\mathrm{id}+\bar w_\Ufield$ is globally bi-Lipschitz with two-sided Jacobian
bounds at every $t\in I$. 

Let $\Psi_1,\dots,\Psi_N$ be the finitely many real-analytic functions used
below, namely $Z\mapsto\det(\mathrm{Id}+Z)$, its cofactors,
$\log\det(\mathrm{Id}+Z)$, $(\mathrm{Id}+Z)^{-1}-\mathrm{Id}$ and the pairwise
products of the entries of the last, and let $C_{\rm e}$ be the constant of
\eqref{eq:embed-frac} at $(\sreg,\preg)=(s-1,1)$.  Fix
\begin{equation}\label{eq:delta-choice}
 \delta:=\min\Big\{\tfrac1{2C_{\rm e}},\
 \tfrac12\min_{1\le k\le N}\delta_{\Psi_k},\
 \delta_{\rm H}\Big\}>0,
\end{equation}
which depends only on the standing parameters.  The first entry makes
\eqref{eq:embed-frac} give $\|\Ufield\|_{L^\infty(I\times\Rn)}\le\frac12$ whenever
$\sup_{t\in I}\|\Ufield(t)\|_{Y^{s-1,1}}\le\delta$; the factor $\frac12$ in the
second entry is what allows Lemma~\ref{lem:coeff-calculus} to be applied to
pairs $\Ufield(t),\Vfield(t)$ of trial fields, as is done in the
difference estimate \eqref{eq:pressure-lipschitz}; the third entry places
every $\Ufield(t)$ in the ball \eqref{eq:hodge-smallness} of
Lemma~\ref{lem:label-hodge}, whose conclusions then apply at each time.\\
We emphasise that $\delta$ can be fixed once and for all,
independently of $R$ and $T$. This is because the constants $C_\delta$ of
Lemma~\ref{lem:pressure} depend on neither, and we can ensure that the norm of the free evolution of the initial data in the definition \eqref{eq:differential-map} of $\mathcal{M}$, given $R$, does not exceed $\delta$ simply by controlling $T$.  Indeed, Plancherel, the definition of the space \eqref{eq:Y} and
estimate \eqref{eq:prop-gain} give, for $0\le t\le T$,
\begin{align}
 \|\Sprop(t)\nabla v_0\|_{Y^{s-1,1}}^2
 &\le C\int_{\Rn}
   \min(1,T\langle\lambda(\xi)\rangle)^2
   \langle\xi\rangle^{2(s-1)}
   |\widehat{\nabla v_0}(\xi)|^2 d\xi \notag\\
 &=C\omega_{v_0}(T)^2.
 \label{eq:free-fixed-time-modulus}
\end{align}
Thus the free wave part in the map $\mathcal{M}$ is small in the fixed-time geometric norm even though its
spacetime norm is of size $\|v_0\|_{H^s}\simeq R$.  This is the only point at which
regime~\textup{(R1)} needs the datum-dependent modulus.  We use the two
properties recorded after \eqref{eq:datum-modulus}:
$\omega_{v_0}(T)\downarrow0$, and
$\omega_{v_0}(T)\le T^{\rpar}
\|\nabla v_0\|_{Y^{s-1,\rpar}}$ when that norm is finite.\\



We are now ready to state the existence result.
\begin{lemma}[the map is a self-map of the class]\label{lem:selfmap}
Let $v_0\in H^{s}(\Rn)$ be divergence-free,
of arbitrary size. Then, for every $T\le T_*$, $\mathcal M$ is a well-defined self-map of the
nonempty class $\mathcal U_{R,\delta}$, with $T_*>0$ depending on
$n,s,\breg,v_{\rm A}$, on $\|v_0\|_{H^{s}}$ and on the modulus $\omega_{v_0}$.
\end{lemma}

\begin{proof}
For $\Ufield\in\mathcal U_{R,\delta}$, Lemma~\ref{lem:pressure}, the definition of $R$ in \eqref{eq:Rdefn} and the bounds \eqref{eq:free-fixed-time-modulus},
\eqref{eq:cutoff-duhamel} give
\[
 \|\mathcal M\Ufield\|_{\mathcal Y}
 \le \frac R2+C_\delta T^\varepsilon R^2,
 \qquad
 \sup_{t\in I}\|\mathcal M\Ufield(t)\|_{Y^{s-1,1}}
 \le C\big(\omega_{v_0}(T)+T^{1/2}C_\delta R^2\big).
\]
The map has the prescribed time traces because both cutoffs are constant near
$t=0$, the Duhamel integral and its first time derivative vanish there, and
$\partial_t\Sprop(0)=\mathrm{Id}$.  It is exact because the free term is
$d(\chi\Sprop(t)v_0)$ and the source in the second term is $dF[\Ufield]$.

We now fix the lifespan once and for all. After enlarging the common constants in the preceding estimates, choose
$T_*\in(0,1]$ so that
\begin{equation}\label{eq:Tstar}
\begin{aligned}
 C_\delta T_*^\varepsilon R^2&\le\tfrac R2,
 &&\text{spacetime radius of the self-map},\\
 C\omega_{v_0}(T_*)&\le\tfrac\delta2,
 &&\text{fixed-time size of the free wave},\\
 CC_\delta T_*^{1/2}R^2&\le\tfrac\delta2,
 &&\text{fixed-time size of the Duhamel term},\\
 C_\delta T_*^\varepsilon R(1+R)&\le\tfrac12,
 &&\text{contraction constant}.
\end{aligned}
\end{equation}
Note that the fourth line is recorded already here so that one choice of $T_*$ works for
both the self-map and contraction lemmas.  Such a positive time exists in
regime~\textup{(R1)} because $\omega_{v_0}(T)\to0$; the first, third and fourth
left-hand sides are ordinary positive powers of $T$.  For every $T\le T_*$,
the first three lines give respectively
\[
 \|\mathcal M\Ufield\|_{\mathcal Y}\le R,
 \qquad
 \sup_{t\in I}\|\mathcal M\Ufield(t)\|_{Y^{s-1,1}}\le\delta.
\]
Together with exactness and the traces, these are precisely the defining
conditions of $\mathcal U_{R,\delta}$.

For non-emptiness, take the free evolution
$\Ufield^{(0)}(t):=\chi(t)\Sprop(t)\nabla v_0$.  The linear estimate and the
definition of $R$ give
$\|\Ufield^{(0)}\|_{\mathcal Y}\le R/2$, while
\eqref{eq:free-fixed-time-modulus} and the second line of
\eqref{eq:Tstar} give its fixed-time bound by $\delta/2$.  It is exact and has
the prescribed data, hence belongs to the class.
\end{proof}

\begin{lemma}[the map is a contraction]\label{lem:contraction}
With $T_*$ chosen by \eqref{eq:Tstar}, the map $\mathcal M$ satisfies
\begin{equation}\label{eq:map-contraction}
 \|\mathcal M\Ufield-\mathcal M\Vfield\|_{\mathcal Y}
 \le C_\delta T^\varepsilon R(1+R)
       \|\Ufield-\Vfield\|_{\mathcal Y}
 \le\tfrac12\|\Ufield-\Vfield\|_{\mathcal Y}
\end{equation}
for all $\Ufield,\Vfield\in\mathcal U_{R,\delta}$ and every $T\le T_*$.
\end{lemma}

\begin{proof}
The free terms in \eqref{eq:differential-map} cancel. Applying the first
estimate in \eqref{eq:cutoff-duhamel} to the difference and then
\eqref{eq:pressure-lipschitz} gives
\[
 \|\mathcal M\Ufield-\mathcal M\Vfield\|_{\mathcal Y}
 \le C_{\rm lin}T^\varepsilon
 \|dF[\Ufield]-dF[\Vfield]\|_{L^2(I;H^{s-1})}
 \le C_\delta T^\varepsilon R(1+R)
       \|\Ufield-\Vfield\|_{\mathcal Y}.
\]
The last line of \eqref{eq:Tstar} makes the coefficient at most $1/2$.
\end{proof}

Let us make the distinction between the two existence regimes quantitative.
In regime~\textup{(R1)}, the second condition in \eqref{eq:Tstar} depends on the
actual high-frequency tail of the datum.  Indeed, for every $\lambda_0\ge1$,
\[
 \omega_{v_0}(T)
 \le T\lambda_0 \|v_0\|_{H^{s}}
 +\big\|\mathbf 1_{\{\langle D_{\parallel}\rangle>\lambda_0\}}
       \Lambda^{s-1}\nabla v_0\big\|_{L^2}.
\]
For a fixed datum, first choose $\lambda_0$ so that the second term is small and
then choose $T$ so that the first is small.  The tail is not uniformly
controlled on bounded subsets of $H^s$, which is why the regime~\textup{(R1)}
lifespan is not a function of $\|v_0\|_{H^s}$ alone.

In regime~\textup{(R2)}, put
$D_{\rpar}:=\|\nabla v_0\|_{Y^{s-1,\rpar}}$.  The estimate
$\omega_{v_0}(T)\le T^{\rpar}D_{\rpar}$ from \eqref{eq:omegav0bound} shows that the second line of
\eqref{eq:Tstar} is ensured by the explicit norm condition
\begin{equation}\label{eq:Tstar-R2}
 CT_*^{\rpar}D_{\rpar}\le\tfrac\delta2.
\end{equation}
All quantities in \eqref{eq:Tstar} and \eqref{eq:Tstar-R2} are then controlled
by $n,s,\breg,\rpar,v_{\rm A}$ and
$\|v_0\|_{H^s}+D_{\rpar}$.  This is the asserted norm dependence of the
regime~\textup{(R2)} lifespan.  To propagate the stronger spacetime norm we
will separately impose
\begin{equation}\label{eq:Tstar-enhanced}
 C_{\rm lin}C_{\delta,\rpar}T_*^\varepsilon R\le\tfrac12.
\end{equation}
Condition \eqref{eq:Tstar-enhanced} is not needed to construct a fixed point by Picard iteration; it only ensures the improved regularity estimate is propagated along the Picard sequence.

\subsection{Lagrangian existence} With the contractive self-map $\mathcal{M}$ at hand,  existence in regime \textup{(R1)} follows quickly from Banach's fixed-point Theorem once $w$ is reconstructed from the fixed point $\Ufield$. Propagation of any positive amount of parallel regularity under Picard iterations addresses regime \textup{(R2)}.

\begin{theorem}[existence for the auxiliary system]\label{thm:existence}
Let $n\ge2$, $\frac{n+1}2<s\le\frac n2+1$, $\breg\in(\frac12,1)$, $\varepsilon=1-\breg$,
let $\Binf\in\Rn\setminus\{0\}$ be fixed, and let
$v_0\in H^{s}(\Rn)$ be divergence-free, of arbitrary size, and with no relation to
$\Binf$. Put
$R:=2C_{\rm lin}\|v_0\|_{H^{s}}$ and let $\delta$ be as
fixed above. There is $T_*>0$, depending on $n,s,\breg,v_{\rm A}$, on
$\|v_0\|_{H^{s}}$ and on the modulus $\omega_{v_0}$, such that $\mathcal M$ carries $\mathcal U_{R,\delta}$
into itself and has a unique fixed point $\Ufield$ in it.  The potential reconstructed from this
fixed point solves \eqref{eq:auxiliary-system} on $[0,T_*]$ and satisfies
$\Ufield=dw$ on $[0,T_*]\times\Rn$. Moreover,
\begin{equation}\label{eq:existence-traces}
 dw\in C\big([0,T_*],Y^{s-1,1}\big)\cap C^1\big([0,T_*],H^{s-1}\big),
 \qquad
 E_i\partial_tX^j\in C\big([0,T_*],H^{s-1}\big).
\end{equation}
If $\|\nabla v_0\|_{Y^{s-1,\rpar}}<\infty$ for some
$\rpar\in(0,1]$, then $T_*$ may be chosen to satisfy
\eqref{eq:Tstar}, \eqref{eq:Tstar-R2} and
\eqref{eq:Tstar-enhanced}; in particular, it depends only on the listed
parameters and on
$\|v_0\|_{H^s}+\|\nabla v_0\|_{Y^{s-1,\rpar}}$.  For this choice,
\begin{equation}\label{eq:longitudinal-propagation}
 \|\Ufield\|_{H^{s-1,1+\rpar}_{\breg}}
 \le 2C_{\rm lin}\|\nabla v_0\|_{Y^{s-1,\rpar}}.
\end{equation}
This is additional Lagrangian spacetime regularity of the differentiated
displacement; no Eulerian anisotropic norm is asserted.
\end{theorem}

\begin{proof}
Take $T=T_*$ and $I=[0,T_*]$. By Corollary~\ref{cor:trace-Y}, convergence in
$\mathcal Y$ implies convergence in $C_tY^{s-1,1}$ and of the first time
derivatives in $C_tH^{s-1}$. Thus the traces and fixed-time bound in
\eqref{eq:fixedpoint-class} are closed. The norm bound is closed as well, and
the continuity of $\Curl_y$ from $\mathcal Y$ to distributions makes the
exactness condition closed. Hence $\mathcal U_{R,\delta}$ is complete.
Lemmas~\ref{lem:selfmap} and~\ref{lem:contraction} and Banach's theorem give a
unique fixed point $\Ufield\in\mathcal U_{R,\delta}$.

\textit{Back to the undifferentiated wave--Hodge system.} Reconstruct the displacement by
\begin{equation}\label{eq:potential-reconstruction}
 w(t):=\Sprop(t)v_0-\int_0^t\Sprop(t-t')F[\Ufield](t') dt'.
\end{equation}
Lemma~\ref{lem:pressure} gives $F[\Ufield]\in C(I;L^4)$.  Since
$v_0\in L^4$, $\partial_t\Sprop(r)=\Cprop(r)$, and
\[
 \|\Sprop(r)\|_{L^4\to L^4}\le |r|,
 \qquad
 \|\Cprop(r)\|_{L^4\to L^4}\le 1,
\]
the Bochner integral in \eqref{eq:potential-reconstruction} is continuously
differentiable in $L^4$, with derivative
\[
 \partial_tw(t)=\Cprop(t)v_0
 -\int_0^t\Cprop(t-t')F[\Ufield](t') dt'.
\]
Thus $w\in C^1(I;L^4)$. Then
$(w,\partial_tw)(0)=(0,v_0)$ and, by \eqref{eq:differential-map},
\[
 dw(t)=\Sprop(t)\nabla v_0
       -\int_0^t\Sprop(t-t')dF[\Ufield](t') dt'
       =(\mathcal M\Ufield)(t)=\Ufield(t)
\]
on $I$ in distributions, where the cutoffs in \eqref{eq:differential-map}
equal one.

This identity is important: it says
that the primitive used to define the geometry of the trial field and the
dynamical displacement reconstructed from the wave equation have the same
spatial differential.  They can differ only by a spatial translation, which
does not change $\gmetric$, $\mathcal Q$ or the Hodge system.  Consequently, $F[\Ufield]=F[dw]$, and this identity recovers the full
system \eqref{eq:auxiliary-system}.
Corollary~\ref{cor:trace-Y} applied to the fixed point gives
$dw\in C_tY^{s-1,1}$ and $\partial_tdw\in C_tH^{s-1}$, which is the first
statement in \eqref{eq:existence-traces}.  By \eqref{eq:component-derivative} we have
\[
 E_i\partial_tX^j
 =\Kmat^a_i\partial_t\Ufield_a^j
 =\partial_t\Ufield_i^j
  +(\Kmat^a_i-\delta_i^a)\partial_t\Ufield_a^j,
\]
then Lemma~\ref{lem:coeff-calculus} gives
$\Kmat-\mathrm{Id}\in C_tY^{s-1,1}$, and the product estimate \eqref{eq:tame-sp} gives the
second statement in \eqref{eq:existence-traces}.

\textit{Uniqueness in the adapted fixed-point class.} Let $\widetilde w$ be another solution of
\eqref{eq:auxiliary-system} on $I$ with the same Cauchy data, in the sense that
$\widetilde w\in C^1(I;L^1_{\rm loc})$,
$\Box\widetilde w=-F[d\widetilde w]$ in $\mathcal D'((0,T_*)\times\Rn)$ and
$d\widetilde w\in C(I;Y^{s-1,1})\cap C^1(I;H^{s-1})$, and suppose that
$d\widetilde w$ agrees on $I$ with the restriction of some
$\widetilde\Ufield\in\mathcal U_{R,\delta}$. Set
$\Vfield:=\mathcal M\widetilde\Ufield$. On $I$, both $\Vfield$ and
$d\widetilde w$ solve
\[
 \Box Z=-dF[\widetilde\Ufield],
 \qquad
 (Z,\partial_tZ)(0)=(0,\nabla v_0).
\]
Lemma~\ref{lem:wave-uniqueness} gives $\Vfield=d\widetilde w$ on $I$. Since
$\Vfield=\mathcal M\widetilde\Ufield$, the self-map property puts $\Vfield$ in
$\mathcal U_{R,\delta}$. Since $\mathcal M$ depends only on the restriction to
$I$, $\mathcal M\Vfield=\mathcal M\widetilde\Ufield=\Vfield$, so fixed-point uniqueness gives
$\Vfield=\Ufield$. Thus $d\widetilde w=dw$ and
$\widetilde w-w=c(t)$. Subtracting the two wave equations gives $c''=0$, and
the common Cauchy data give $c=0$. Hence $\widetilde w=w$.

\textit{Additional parallel regularity \textup{(R2)}.} Suppose now that
$D_{\rpar}:=\|\nabla v_0\|_{Y^{s-1,\rpar}}<\infty$ for some
$0<\rpar\le1$. Start the Picard iteration from the free field
\[
 \Ufield^{(0)}(t):=\chi(t)\Sprop(t)\nabla v_0,
 \qquad
 \Ufield^{(j+1)}:=\mathcal M\Ufield^{(j)}.
\]
The proof of Lemma~\ref{lem:selfmap} puts the free field in
$\mathcal U_{R,\delta}$, and the self-map property keeps every iterate there.
The contraction estimate makes the sequence converge strongly in $\mathcal Y$
to the fixed point.
To obtain a bound in the stronger space, put
\[
 C_j:=\|\Ufield^{(j)}\|_{H^{s-1,1+\rpar}_{\breg}}.
\]
The linear estimate with $(\sreg,\preg)=(s-1,1+\rpar)$ first gives
$C_0\le C_{\rm lin}D_{\rpar}$.  For the next iterate, the same estimate
places the source in $Y^{s-1,\rpar}$, and
\eqref{eq:pressure-anisotropic}, together with the uniform bound
$\|\Ufield^{(j)}\|_{\mathcal Y}\le R$, guaranteed by the class, gives
\[
 C_{j+1}
 \le C_{\rm lin}D_{\rpar}
   +C_{\rm lin}C_{\delta,\rpar}T^\varepsilon R C_j.
\]
Condition \eqref{eq:Tstar-enhanced} makes the coefficient of $C_j$ at most
$1/2$. Starting from $C_0\le C_{\rm lin}D_{\rpar}$, induction gives
$C_j\le2C_{\rm lin}D_{\rpar}$ for every $j$. A subsequence therefore converges
weakly in the Hilbert space $H^{s-1,1+\rpar}_{\breg}$. Its continuous embedding
into $\mathcal Y$ and the strong convergence there identify its weak limit with
$\Ufield$. Weak lower semicontinuity proves
\eqref{eq:longitudinal-propagation}.
\end{proof}

\section{Determinant Constraint and Return to MHD}\label{sec:return-mhd}
\subsection{The volume constraint propagates}
\label{ssec:constraint}

The system \eqref{eq:auxiliary-system} does not contain $\det dX=1$; however, any solution
propagates it dynamically, at the regularity of the class.

\begin{proposition}[weak Jacobi's identity]\label{prop:constraint}
Let $w$ solve \eqref{eq:auxiliary-system} on $[0,T]$ with $v_0$ divergence-free and
\begin{equation}\label{eq:constraint-hyp}
\Ufield=dw\ \text{on }[0,T]\times\Rn
\ \text{for some }\Ufield\in H^{s-1,1}_{\breg},
\qquad \big\|dw\big\|_{L^\infty}\le\tfrac12,\qquad
d(\Box X)\in L^1_tL^2_{y} .
\end{equation}
Then $\det dX\equiv1$ on $[0,T]$.
\end{proposition}

\begin{proof}
Recall the notation
\[
 \sigma=s-1
 ,
 \qquad
 A=dX=\mathrm{Id}+dw,
 \qquad
 \Kmat=A^{-1},
 \qquad
 \varrho:=\log\det A.
\]
The trace estimate of Corollary~\ref{cor:trace-Y} applied to $\Ufield$, together with
$\Ufield=dw$ on $[0,T]$, gives
\[
 dw\in C_tY^{\sigma,1},
 \qquad
 \partial_tdw,\ \mathcal Bdw\in C_tH^\sigma.
\]
Consequently,
\[
 \mathcal A^\pm dw
 =\partial_tdw\pm\mathcal Bdw
 \in L^2_tH^\sigma
 \hookrightarrow L^2_tL^4_y,
\]
because $\sigma>(n-1)/2\ge n/4$.  The pointwise assumption
$\|dw\|_{L^\infty}\le 1/2$ keeps $A$ in a fixed compact subset of
the invertible matrices and therefore gives
\[
 \Kmat\in L^\infty.
\]

For smooth $X$ the following computation is immediate from Jacobi's formula.
At the present regularity, differentiating the first-order identity once more
and identifying the resulting quadratic product are not automatic, so we
justify the full second-order identity by mollification. Fix
$I'\Subset(0,T)$, mollify $dw$ in spacetime using only values inside
$(0,T)$, and set
\[
 A_\ell:=\mathrm{Id}+(dw)_\ell,
 \qquad
 \Kmat_\ell:=A_\ell^{-1},
 \qquad
 \varrho_\ell:=\log\det A_\ell.
\]
The mollifiers preserve
$\|A_\ell-\mathrm{Id}\|_{L^\infty}\le1/2$, so the matrices
$\Kmat_\ell$ are uniformly bounded. All the differential operators involved
have constant coefficients and commute with mollification. Jacobi's formula
and differentiation of $\Kmat_\ell A_\ell=\mathrm{Id}$ give
\[
 \mathcal A^\pm\varrho_\ell
 =\tr\bigl(\Kmat_\ell\mathcal A^\pm A_\ell\bigr),
 \qquad
 \mathcal A^+\Kmat_\ell
 =-\Kmat_\ell(\mathcal A^+A_\ell)\Kmat_\ell.
\]
Therefore
\begin{align*}
 \Box\varrho_\ell
 &=\mathcal A^+\tr\bigl(\Kmat_\ell\mathcal A^-A_\ell\bigr)\\
 &=\tr\bigl(\Kmat_\ell d(\Box X)_\ell\bigr)
   -\tr\bigl(
       \Kmat_\ell(\mathcal A^+A_\ell)
       \Kmat_\ell(\mathcal A^-A_\ell)
     \bigr).
\end{align*}
Here $\mathcal A^\pm A_\ell=\mathcal A^\pm(dw)_\ell$ and
$\Box A_\ell=d(\Box X)_\ell$. We use the cyclicity of the trace explicitly:
\[
 \tr\bigl(
   \Kmat_\ell\mathcal A^+(dw)_\ell
   \Kmat_\ell\mathcal A^-(dw)_\ell
 \bigr)
 =\tr\bigl(
   \Kmat_\ell\mathcal A^-(dw)_\ell
   \Kmat_\ell\mathcal A^+(dw)_\ell
 \bigr).
\]
After averaging these equal expressions, the definition of $Q_0$ gives
\begin{align*}
 \tr(\Kmat_\ell \mathcal A^+(dw)_\ell \Kmat_\ell \mathcal A^-(dw)_\ell)
 &=\frac12\tr\bigl(
      \Kmat_\ell \mathcal A^+(dw)_\ell \Kmat_\ell \mathcal A^-(dw)_\ell
      +\Kmat_\ell \mathcal A^-(dw)_\ell \Kmat_\ell \mathcal A^+(dw)_\ell
    \bigr)\\
 &=\frac12\Kmat_{\ell,i}^a\Kmat_{\ell,j}^b
   \Big[
      (\mathcal A^+(dw)_\ell)_a^j
      (\mathcal A^-(dw)_\ell)_b^i
     +(\mathcal A^-(dw)_\ell)_a^j
      (\mathcal A^+(dw)_\ell)_b^i
   \Big]\\
 &=\Kmat_{\ell,i}^a\Kmat_{\ell,j}^b
   Q_0\bigl((dw)_{\ell,a}^{ j},(dw)_{\ell,b}^{ i}\bigr)\\
 &=-\mathcal Q[(dw)_\ell].
\end{align*}
Hence
\begin{equation*}
    \Box\varrho_\ell
 =\mathcal Q[(dw)_\ell]
  +\tr\bigl(\Kmat_\ell d(\Box X)_\ell\bigr).
\end{equation*}
Put
\[
 P_\ell^\pm:=\Kmat_\ell\mathcal A^\pm(dw)_\ell,
 \qquad
 P^\pm:=\Kmat\mathcal A^\pm dw,
 \qquad
 G_\ell:=d(\Box X)_\ell,
 \qquad
 G:=d(\Box X).
\]
After passing to a subsequence, $\Kmat_\ell\to\Kmat$ almost everywhere,
and the matrices $\Kmat_\ell$ are uniformly bounded. Standard mollifier
convergence gives
$\mathcal A^\pm(dw)_\ell\to\mathcal A^\pm dw$ in
$L^2(I';L^4)$ and $G_\ell\to G$ in $L^1(I';L^2)$. The decompositions
\begin{align*}
 P_\ell^\pm-P^\pm
 &=\Kmat_\ell\bigl(\mathcal A^\pm(dw)_\ell
                    -\mathcal A^\pm dw\bigr)
   +(\Kmat_\ell-\Kmat)\mathcal A^\pm dw,\\
 \Kmat_\ell G_\ell-\Kmat G
 &=\Kmat_\ell(G_\ell-G)+(\Kmat_\ell-\Kmat)G
\end{align*}
and dominated convergence therefore give
\[
 \begin{gathered}
  \varrho_\ell\longrightarrow\varrho
  \quad\text{in }L^2_{\mathrm{loc}}(I'\times\Rn),\\
  P_\ell^\pm\longrightarrow P^\pm
  \quad\text{in }L^2(I';L^4_y),\\
  \tr(\Kmat_\ell G_\ell)\longrightarrow\tr(\Kmat G)
  \quad\text{in }L^1(I';L^2_y).
 \end{gathered}
\]
Here the first convergence follows as well from the Lipschitz continuity of
$M\mapsto\log\det(\mathrm{Id}+M)$ on the fixed compact matrix range. Since
$\mathcal Q$ is a finite sum of products of $P_\ell^+$ and $P_\ell^-$,
the $L^2_tL^4_y$ convergence gives
$\mathcal Q[(dw)_\ell]\to\mathcal Q[dw]$ in $L^1(I';L^2_y)$. Since $I'$ was
arbitrary, passing to the limit gives
\begin{equation}\label{eq:weak-jacobi}
 \Box\varrho
 =\mathcal Q[dw]+\tr\bigl(\Kmat d(\Box X)\bigr)
 \qquad\text{in }\mathcal D'((0,T)\times\Rn).
\end{equation}

Passing to the limit in the first-order Jacobi formulas also gives
\begin{equation}\label{eq:constraint-first-order-jacobi}
 \partial_t\varrho
 =\tr(\Kmat\partial_tdw),
 \qquad
 \mathcal B\varrho
 =\tr(\Kmat\mathcal Bdw).
\end{equation}
Since $dw\in C_tY^{\sigma,1}\hookrightarrow C_tL^\infty$ and the maps
$A\mapsto A^{-1}$ and $A\mapsto\log\det A$ are Lipschitz on the fixed set of
matrices under consideration, we have
\[
 \Kmat\in C([0,T];L^\infty),
 \qquad
 \varrho\in C([0,T];L^2).
\]
Together with $\partial_tdw,\mathcal Bdw\in C_tL^2$, the first-order
identities imply
\[
 \varrho\in C^1([0,T];L^2),
 \qquad
 \mathcal B\varrho\in C([0,T];L^2).
\]
In particular, they hold also at $t=0$ and $t=T$.

The trace bounds give
$\mathcal Q[dw]\in L^\infty_tL^{r_{\mathrm{HLS}}}_y$, so the
Hardy--Littlewood--Sobolev estimate and the uniform Jacobian bounds give
$F\in L^\infty_tL^4_y$. Together with
$dF=-d(\Box X)\in L^1_tL^2_y$, this gives
$F\in L^1_tW^{1,2}_{\rm loc}$. Bi-Lipschitz composition gives the same local
regularity for $\mathcal C_X^{-1}F$, so Lemma~\ref{lem:weak-pullback}
identifies $\tr(\Kmat dF)$ with $\Div_{\gmetric}F$. Using the Hodge equation
in \eqref{eq:weak-jacobi}, we conclude that
\[
 \Box\varrho
 =\mathcal Q[dw]-\Div_{\gmetric}F
 =0
 \qquad\text{in }\mathcal D'((0,T)\times\Rn).
\]
The Cauchy data for $X$ give
\[
 A(0)=\mathrm{Id},
 \qquad
 \varrho(0)=0,
\]
and \eqref{eq:constraint-first-order-jacobi} gives
\[
 \partial_t\varrho(0)
 =\tr(dv_0)
 =\Div_yv_0
 =0.
\]

Since $\varrho\in C^1([0,T];L^2)$ with $\mathcal B\varrho\in C([0,T];L^2)$,
$\Box\varrho=0$ in $\mathcal D'$ and $\varrho(0)=\partial_t\varrho(0)=0$,
Lemma~\ref{lem:wave-uniqueness} gives $\varrho\equiv0$.
Therefore $\det dX=\exp(\varrho)\equiv1$, as claimed.
\end{proof}


\begin{remark}[generalized Jacobi identity]\label{rem:generalcaseJacobi}
More generally, \eqref{eq:weak-jacobi} gives
$\Box\varrho=\mathcal Q[dw]-q$ whenever
$\Box X=-F$ and $\Div_{\gmetric}F=q$. Thus the inverse-deformation factors
$\Kmat=A^{-1}$ in \eqref{eq:datum}, rather than the cofactor factors
$J\Kmat$, are exactly what makes the unconstrained extension propagate volume
in contrast with the constrained formulation of \cite{Zhang2024}.
\end{remark}

\begin{remark}[weak limits and the relaxation of the volume constraint]
\label{rem:weak-limit-heuristic}
The following heuristic, stated in $n=3$ for concreteness, indicates that
dropping the volume constraint on the trial fields is the natural relaxation.
Let $(v_m,B_m)$ be smooth solutions of \eqref{eq:mhd} and let
$E_m:=B_m\times v_m$ be the associated electric fields, as prescribed by
the ideal Ohm law. Assume $v_m$, $B_m$ and $E_m$ are bounded in
$L^2_{\rm loc}(\Real^{1+3})$ and converge weakly there:
\[
 v_m\rightharpoonup v,\qquad B_m\rightharpoonup B,\qquad
 E_m\rightharpoonup E,
\]
with $B\neq0$ almost everywhere.  Write $\mathrm{vol}_x=dx^1\wedge dx^2\wedge
dx^3$ and let
\[
 F_m:=\iota_{B_m}\mathrm{vol}_x+E_m^\flat\wedge dt
\]
be the Faraday $2$-form. Then $dF_m=0$: the $dx$-part is $\Div_xB_m=0$ and the
$dt$-part is Faraday's law $\partial_tB_m+\curl E_m=0$. Moreover
$F_m\wedge F_m=-2(B_m\cdot E_m)dt\wedge\mathrm{vol}_x=0$. The div--curl lemma of Murat and
Tartar~\cite{Murat1978,Tartar1979}, in the form of the weak continuity of wedge
products of closed forms, therefore passes this to the limit:
\[
 B\cdot E=0 .
\]
 
What compensated compactness does \emph{not} give is the Ohm law itself: there
is no reason for $B_m\times v_m$ to converge weakly to $B\times v$. All that
survives is that $E$ lies pointwise in the range of $B\times(\cdot)$,
which is the plane $B^\perp$. Pointwise almost everywhere, choose a measurable
scalar $\alpha$ and set
\[
 V_{\rm eff}:=\frac{E\times B}{|B|^2}+\alpha B,
 \qquad
 u:=V_{\rm eff}-v.
\]
Then $E=B\times V_{\rm eff}=B\times(v+u)$ and letting
$\beta:=\iota_B\mathrm{vol}_x$ be the
$2$-form dual to $B$, after passing by linearity to the limit in the Faraday--Ohm systems one formally obtains
\[
 \partial_t\beta+\mathcal L_{V_{\rm eff}}\beta=0.
\]
There is no reason, however, to have $\Div V_{\rm eff}=0$ and this could be read as the Eulerian counterpart of running the fixed point on trial
deformations with $\det(\mathrm{Id}+\Ufield)\neq1$.
\end{remark}

\subsection{Eulerian reconstruction and existence}
\label{ssec:eulerian}

We now show that an auxiliary solution determines a distributional solution of
\eqref{eq:mhd}.
Here $\nabla v$ and $\nabla B$ lie only in $C_tH^{s-1}$ and $\nabla^2p$ only in
$C_tH^{s-2}$, so the Eulerian derivatives entering the formal smooth change-of-variables
calculation that leads to \eqref{eq:auxiliary-system} have no pointwise meaning,
and that calculation cannot simply be reversed through $X^{-1}$.  The four Eulerian statements, that is, the two
divergence conditions, the induction equation and the momentum equation, are therefore
verified below against test functions in a weak sense. 

Combining the work from \S \ref{sec:fixedpoint} with the volume constraint of
Proposition~\ref{prop:constraint}, it remains to pass to Eulerian variables
and prove the conservation laws. We use the trace of
Corollary~\ref{cor:trace-Y} and the multiplier bound of
Proposition~\ref{prop:tame-product} from Section~\ref{sec:anisotropic-spaces}.

We also recall that the component derivations $E_i$ are skew-adjoint on $L^2(\mu)$ by the
pullback calculus of Lemma~\ref{lem:weak-pullback}.
Constancy gives $\nabla\Binf=0$ and makes $\mathcal B$ skew-adjoint. We finally remark that the
estimates depend on the field only through $v_{\rm A}$.


\begin{proof}[Proof of clauses~\textup{(1)}--\textup{(3)} and the existence
statements in Theorem~\ref{thm:main}]
Set $T:=T_*$. Let $\Ufield$ be the fixed point of Theorem~\ref{thm:existence}, let $w$ be its
reconstructed potential, and put
\[
 X=\mathrm{id}+w,\qquad F=F[dw],\qquad
 \vLag=\partial_tX,\qquad \BLag=\mathcal BX,\qquad
 v=\vLag\circ X^{-1},\qquad B=\BLag\circ X^{-1}.
\]
Then $(w,\partial_tw)(0)=(0,v_0)$, so
$(v,B)|_{t=0}=(v_0,\Binf)$.  Moreover
$\Ufield=dw$ on $[0,T_*]\times\Rn$, with
$\Ufield\in\mathcal U_{R,\delta}$, where
$R=2C_{\rm lin}\|v_0\|_{H^s}$; this is \eqref{eq:main-adapted-bound}.

\emph{The constraint and preliminary bounds.}
The fixed-point bounds and \eqref{eq:embed-frac} give
$\|dw\|_{L^\infty}\le\frac12$, while Lemma~\ref{lem:pressure} gives
$dF\in L^2_tH^{s-1}\subset L^1_tL^2$.  Proposition~\ref{prop:constraint}
therefore yields
\[
 \det dX=1.
\]
The bound $\|dw\|_{L^\infty}\le\frac12$ and arguing as in \eqref{eq:trial-bilipschitz} gives that
$X(t,\cdot)$ is a volume-preserving bi-Lipschitz homeomorphism, which proves
the geometric assertions in clause~(2).

Recall $\sigma=s-1$ and
$\bLag:=b\circ X=\BLag-\Binf=\mathcal Bw=(dw)\Binf$.
The trace estimates \eqref{eq:existence-traces} give
\begin{equation}\label{eq:preliminaryeulerian}
    \partial_tdw,\ \mathcal Bdw\in C_tH^\sigma,\qquad
 \bLag\in C_tH^\sigma,
\end{equation}
and $H^\sigma\hookrightarrow L^4$ because
$\sigma>(n-1)/2\ge n/4$; in particular $v_0\in L^4$.  Also
$\mathcal B^2w=(\mathcal Bdw)\Binf\in C_tH^\sigma$, while
$F\in C_tL^4$ by Lemma~\ref{lem:pressure}.  Integrating
$\partial_t\vLag=\mathcal B^2w-F$ from $\vLag(0)=v_0$ gives $\vLag\in C_tL^4$;
together with $\bLag\in C_tH^\sigma$ this yields
\begin{equation}\label{eq:return-preliminary}
 \vLag,\bLag\in C_tL^4,\qquad
 d\vLag=\partial_tdw,\quad d\bLag=\mathcal Bdw
 \quad\text{in }C_tH^\sigma.
\end{equation}

\emph{Return to Eulerian variables.}
Set $g:=F\circ X^{-1}$, the candidate Eulerian pressure gradient. For
$\varphi\in C_c^\infty([0,T)\times\Rn;\Rn)$, divergence-free or not, put
$\Phi:=\varphi\circ X$. The Sobolev chain rule gives
\begin{equation}\label{eq:sobolevchainrule}
 \partial_t\Phi=(\partial_t\varphi)\circ X
 +((\nabla\varphi)\circ X)\vLag,
 \qquad
 \mathcal B\Phi=((\nabla\varphi)\circ X)\BLag
 \quad\text{in }\mathcal D'\text{ and a.e.}
\end{equation}
These compositions have supports contained in a fixed compact set in label
space. Indeed, \eqref{eq:return-preliminary} and $w(0)=0$ give
\[
 w(t)=\int_0^t\vLag(r)dr\in L^\infty_tL^4,
\]
and the Gagliardo--Nirenberg inequality gives
\[
 \|w\|_{L^\infty_{t,y}}
 \lesssim
 \|w\|_{L^\infty_tL^4}^{4/(4+n)}
 \|dw\|_{L^\infty_{t,y}}^{n/(4+n)}<\infty.
\]
Thus $\supp(\varphi\circ X)(t)$ lies in a fixed compact set. Since
$\Phi\in W^{1,2}$ has fixed compact support and
$\vLag,\BLag,F\in L^2_{\rm loc}$, the distributional identity extends to
$\Phi$ by density; the initial term is justified by the $W^{1,2}$ trace
theorem. The same argument applies to $\Psi=\psi\circ X$ below.

Using volume preservation, $\Phi(0)=\varphi(0)$ and
$\partial_t\vLag-\mathcal B\BLag=-F$, we obtain
\begin{equation}\label{eq:return-momentum}
\begin{aligned}
 &\iint v\cdot\partial_t\varphi  dx dt
   +\iint (v\otimes v-B\otimes B):\nabla\varphi  dx dt
   +\int v_0\cdot\varphi(0) dx\\
 &\qquad=\iint \vLag\cdot\partial_t\Phi  dy dt
   -\iint \BLag\cdot\mathcal B\Phi  dy dt
   +\int \vLag(0)\cdot\Phi(0) dy\\
 &\qquad=-\iint(\partial_t\vLag-\mathcal B\BLag)\cdot\Phi  dy dt\\
 &\qquad=\iint F\cdot\Phi  dy dt
 =\iint g\cdot\varphi  dx dt.
\end{aligned}
\end{equation}
For tests supported in $(0,T)\times\Rn$, this says
\[
 \partial_tv+\Div_x(v\otimes v-B\otimes B)=-g
 \quad\text{in }\mathcal D'.
\]

For $\psi\in C_c^\infty([0,T)\times\Rn;\Rn)$, put
$\Psi:=\psi\circ X$. Since $\BLag(0)=\Binf$ and
$\partial_t\BLag=\mathcal B\vLag$, the same calculation gives
\begin{equation}\label{eq:return-induction}
\begin{aligned}
 &\iint B\cdot\partial_t\psi  dx dt
   +\iint (B\otimes v-v\otimes B):\nabla\psi  dx dt
   +\int\Binf\cdot\psi(0) dx\\
 &\qquad=\iint \BLag\cdot\partial_t\Psi  dy dt
   -\iint \vLag\cdot\mathcal B\Psi  dy dt
   +\int\BLag(0)\cdot\Psi(0) dy\\
 &\qquad=-\iint(\partial_t\BLag-\mathcal B\vLag)\cdot\Psi  dy dt=0.
\end{aligned}
\end{equation}
This is the weak induction equation of Section~\ref{sec:introduction}.

The first-order identities in \eqref{eq:constraint-first-order-jacobi} from the proof of Proposition~\ref{prop:constraint} together with the definition \eqref{eq:divcurlg} now read
\[
 \Div_{\gmetric}\vLag=\partial_t\log\det dX=0,\qquad
 \Div_{\gmetric}\BLag=\mathcal B\log\det dX=0
\]
and we can conclude the corresponding Eulerian ones. Indeed, for every scalar
test function $\zeta\in C_c^\infty(\Rn)$, put $Z:=\zeta\circ X$. The
conservative weak formulation in Lemma~\ref{lem:weak-pullback}(ii), extended
from smooth tests to compactly supported Lipschitz tests by density, gives
\begin{equation}\label{eq:diveulerian}
    \langle\Div_xv,\zeta\rangle
 =-\int v^i\partial_i\zeta dx
 =-\int \vLag^i\Kmat_i^a\partial_aZ dy=0.
\end{equation}
The identical calculation with $\BLag$ proves $\Div_xB=0$.

\emph{Pressure recovery.}
Define
\begin{equation}\label{eq:return-pressure}
 p_{\mathrm R}:=\Riesz_i\Riesz_j
 \big(v^iv^j-\Binf^ib^j-b^i\Binf^j-b^ib^j\big),
 \qquad b:=B-\Binf.
\end{equation}
The preliminary bounds, the component-change identity and volume preservation
give
\[
 (\partial_{x^i}v^j)\circ X
 =E_i\vLag^j=\Kmat_i^a\partial_a\vLag^j,
 \qquad
 (\partial_{x^i}b^j)\circ X
 =E_i\bLag^j=\Kmat_i^a\partial_a\bLag^j.
\]
Consequently,
\[
 \|\nabla p_{\mathrm R}\|_{L^2}
 \lesssim\|v\|_{L^4}\|\nabla v\|_{L^4}
 +\|b\|_{L^4}\|\nabla b\|_{L^4}
 +v_{\rm A}\|\nabla b\|_{L^2},
\]
uniformly in time. By the conjugated component calculus, the Hodge equation
$\curl_{\gmetric}F=0$ gives $\curl_xg(t)=0$ for almost every $t$. Taking
$\Div_x$ in the interior distributional form of \eqref{eq:return-momentum}
and using $\Div_xv=0$ gives
\[
 \Div_xg
 =-\partial_i\partial_j(v^iv^j-B^iB^j)
 =\Delta_xp_{\mathrm R}
 \quad\text{in }\mathcal D'((0,T)\times\Rn).
\]
Consequently, for almost every $t$, the field
$g_0:=g-\nabla p_{\mathrm R}$ belongs to $L^4+L^2$ and satisfies
$\curl_xg_0=\Div_xg_0=0$. Hence $\Delta_xg_0=0$ in distributions. Since
$g_0$ is tempered, its Fourier transform is supported at the origin and
$g_0$ is a polynomial. No nonzero polynomial belongs to $L^4+L^2$, so
$g_0=0$. Therefore
\begin{equation}\label{eq:return-force}
 g=\nabla p_{\mathrm R},\qquad
 F=(\nabla p_{\mathrm R})\circ X\in L^\infty_tL^2.
\end{equation}
Taking $p=p_{\mathrm R}$, equations \eqref{eq:return-momentum} and
\eqref{eq:return-induction}, together with the two divergence conditions,
show that $(v,B,p)$ solves \eqref{eq:mhd} in distributions with data
$(v_0,\Binf)$ and normalised pressure $p$.

\textit{Final regularity statements.} We now upgrade the regularity of the objects to the claimed $C_t H^s$.  Since by \eqref{eq:preliminaryeulerian} we have $\mathcal B^2w=(\mathcal Bdw)\Binf\in C_tH^\sigma \subset  C_tL^2$, the newly established
\eqref{eq:return-force} gives
\[
 \partial_t\vLag=\mathcal B^2w-F\in L^\infty_tL^2.
\]
Since $\vLag(0)=v_0\in L^2$, it follows by time integration that
$\vLag\in W^{1,\infty}_tL^2\subset C_tL^2$. Moreover,
$\bLag\in C_tH^\sigma\subset C_tL^2$ by
\eqref{eq:preliminaryeulerian}. Combining these facts with
$d\vLag,d\bLag\in C_tH^\sigma$ and the elementary characterisation
$\|f\|_{H^s}\simeq\|f\|_{L^2}+\|df\|_{H^{s-1}}$, this yields
\[
 \vLag,\bLag\in C_tH^s.
\]
Finally, $\partial_tw=\vLag$ and $w(0)=0$ give $w\in C_tH^s$, and
Lemma~\ref{lem:coeff-calculus} gives $\Kmat-\mathrm{Id}\in C_tY^{s-1,1}$; with
$dw\in C_tY^{s-1,1}$ these are the hypotheses of Lemma~\ref{lem:composition},
which applied to $v=\mathcal C_X^{-1}\vLag$ and $b=\mathcal C_X^{-1}\bLag$
gives $v,b\in C_tH^s$.  Since $H^s$ is an algebra and the Riesz transforms
are bounded on it, \eqref{eq:return-pressure} gives $p\in C_tH^s$.  This proves clause~(1).

\emph{Conservation.}
Since $J=1$, the quantities in \eqref{eq:energy} are
\[
 E=\tfrac12\int\big(|\vLag|^2+|\bLag|^2\big) dy,\qquad
 H_{\rm c}=\int\vLag\cdot\bLag dy.
\]
They are absolutely continuous, and
\[
 \partial_t\vLag=\mathcal B\bLag-F,\qquad
 \partial_t\bLag=\mathcal B\vLag.
\]
Indeed, the right-hand sides belong to $L^\infty_tL^2$, so
$\vLag,\bLag\in W^{1,\infty}_tL^2$ and the Hilbert-space product rule
applies.  Skew-adjointness of $\mathcal B$ gives
\[
\begin{aligned}
 \frac{dE}{dt}
 &=\langle\mathcal B\bLag,\vLag\rangle
   +\langle\mathcal B\vLag,\bLag\rangle
   -\langle F,\vLag\rangle,\\
 \frac{dH_{\rm c}}{dt}
 &=\langle\mathcal B\bLag,\bLag\rangle
   +\langle\vLag,\mathcal B\vLag\rangle
   -\langle F,\bLag\rangle.
\end{aligned}
\]
The two Alfv\'en terms in the first line cancel, while the first two terms in
the second line vanish separately.  It remains only to use the orthogonality
of the exact force $F$ to the $\gmetric$-divergence-free vector fields along $X$,
$\vLag$ and $\bLag$.  More explicitly,
\[
 \int F\cdot\vLag dy
 =\int\nabla p_{\mathrm R}\cdot v dx=0,
 \qquad
 \int F\cdot\bLag dy
 =\int\nabla p_{\mathrm R}\cdot b dx=0.
\]
Therefore
\[
 \frac{dE}{dt}=0,\qquad \frac{dH_{\rm c}}{dt}=0.
\]
This proves clause~(3); subtracting $\Binf$ is essential because the
unrenormalised magnetic energy is infinite.

\emph{The existence regimes.}
Regime~\textup{(R1)} has the lifespan furnished by
Theorem~\ref{thm:existence} and Lemma~\ref{lem:selfmap}.  If the datum has a positive
longitudinal fraction, then
$\omega_{v_0}(T)\le T^{\rpar}
\|\nabla v_0\|_{Y^{s-1,\rpar}}$; the Picard argument in
Theorem~\ref{thm:existence} also gives \eqref{eq:longitudinal-propagation}.
This proves regime~\textup{(R2)}.
\end{proof}

\subsection{Eulerian uniqueness}
\label{ssec:unconditional-uniqueness}

Banach's theorem gives uniqueness of the differentiated solution inside the
fixed-point class. We now prove the comparison statement needed for uniqueness
in the full Eulerian $C_tH^s$ class. The preceding construction has already
shown that the map $X$ is volume preserving and that its induced Eulerian
fields solve MHD. No flow is constructed for the comparison solution.

We use italic letters for Eulerian fields and sans serif letters for fields
composed with $X$. We first recover the pressure directly from the weak
formulation.

\begin{lemma}[the pressure of an $H^s$ distributional solution]
\label{lem:pressure-recovery}
Let $s>\frac n2$ and let $(v,B)$ be a distributional solution of
\eqref{eq:mhd} on $[0,T]\times\Rn$ with data $(v_0,\Binf)$, in the sense of
Section~\ref{sec:introduction}, such that
$v,b=B-\Binf\in C([0,T],H^s)$. Define
\[
 p_{\mathrm R}:=\Riesz_i\Riesz_j
 \big(v^iv^j-\Binf^ib^j-b^i\Binf^j-b^ib^j\big).
\]
Then $p_{\mathrm R}\in C([0,T],H^s)$,
$\nabla p_{\mathrm R}\in C([0,T],H^{s-1})$, and
\[
 \partial_tv+\Div_x(v\otimes v-B\otimes B)+\nabla p_{\mathrm R}=0
 \quad\text{in }\mathcal D'((0,T)\times\Rn).
\]
\end{lemma}

\begin{proof}
Since $H^s$ is an algebra and the Riesz transforms are bounded on $H^s$, the
tensor
\[
 \mathbb T:=v\otimes v-B\otimes B+\Binf\otimes\Binf
 =v\otimes v-\Binf\otimes b-b\otimes\Binf-b\otimes b
\]
and $p_{\mathrm R}=\Riesz_i\Riesz_j\mathbb T^{ij}$ belong to
$C([0,T],H^s)$. Put $G:=\Div_x\mathbb T\in C([0,T],H^{s-1})$. The weak
momentum equation says that $\partial_tv+G$ annihilates every compactly
supported smooth divergence-free test field. A standard solenoidal cutoff
approximates every divergence-free Schwartz field by such tests. Hence, with
$\mathbb P$ the Euclidean Leray projector,
\[
 \mathbb P(\partial_tv+G)=0
 \quad\text{in }\mathcal S'((0,T)\times\Rn).
\]
Since $\Div_xv=0$, one has $\mathbb P\partial_tv=\partial_tv$ in
distributions, and therefore
\[
 \partial_tv=-\mathbb P G.
\]
Moreover,
\[
 G-\mathbb PG=\nabla\Delta_x^{-1}\Div_xG=-\nabla p_{\mathrm R}.
\]
Thus $\partial_tv+G+\nabla p_{\mathrm R}=0$. The constant tensor
$\Binf\otimes\Binf$ has zero divergence, so this is the required momentum
equation.
\end{proof}

\begin{proposition}[weak-strong uniqueness]
\label{prop:unconditional-uniqueness}
Let $(v,B,p)$ and $X$ be the solution and the volume-preserving map
constructed in \S\ref{ssec:eulerian} on $[0,T_*]$.
If $(\widetilde v,\widetilde B)$ is any other distributional solution with
the same initial data and
\[
 \widetilde v,\ \widetilde B-\Binf\in C([0,T_*],H^s),
\]
then
\[
 \widetilde v=v,\qquad \widetilde B=B
 \quad\text{on }[0,T_*].
\]
The two normalised Riesz pressures also agree. In particular, clause~\textup{(4)} of Theorem~\ref{thm:main} holds.
\end{proposition}

\begin{proof}
\emph{Step 1: the difference equation in the reference geometry.}
Apply Lemma~\ref{lem:pressure-recovery} to both solutions and call the two
normalised pressures $p$ and $\widetilde p$.  Let
\[
 Z^\pm=v\pm B,
 \qquad
 \widetilde Z^\pm=\widetilde v\pm\widetilde B,
 \qquad
 \delta Z^\pm=\widetilde Z^\pm-Z^\pm,
\]
put $\delta p:=\widetilde p-p$, and, after composition with $X$, set
\[
 \mathsf Z^\pm:=Z^\pm\circ X=\mathcal A^\pm X,
 \qquad
 \delta\mathsf Z^\pm:=\delta Z^\pm\circ X,
 \qquad
 \delta\mathsf p:=\delta p\circ X.
\]
Subtracting the two weak Elsasser systems gives in $\mathcal D'$
\[
 \partial_t\delta Z^\pm
 +\Div_x\bigl(\delta Z^\pm\otimes Z^\mp
 +(Z^\pm+\delta Z^\pm)\otimes\delta Z^\mp\bigr)
 +\nabla\delta p=0.
\]
The recovered momentum equation and the weak induction equation imply
$\partial_t\delta Z^\pm\in C_tH^{s-1}$, and hence
$\delta Z^\pm\in W^{1,2}_{\mathrm{loc}}$ in spacetime.  Moreover, the map
$(t,y)\mapsto(t,X(t,y))$ is bi-Lipschitz because $X(t,\cdot)$ is uniformly
bi-Lipschitz and $\partial_tX=v\circ X\in L^\infty$.  The Sobolev chain
rule therefore gives, componentwise,
\[
 \partial_t\delta\mathsf Z^\pm
 =(\partial_t\delta Z^\pm)\circ X
   +(\partial_tX)^jE_j\delta\mathsf Z^\pm,
 \qquad
 \mathcal B\delta\mathsf Z^\pm
 =(\mathcal BX)^jE_j\delta\mathsf Z^\pm
 \quad\text{a.e. and in }\mathcal D'.
\]
Since $\mathcal A^\mp X=\mathsf Z^\mp=Z^\mp\circ X$ and
$\Div_xZ^\mp=0$, these identities identify the pullback of
$\partial_t\delta Z^\pm+
\Div_x(\delta Z^\pm\otimes Z^\mp)$ with
$\mathcal A^\mp\delta\mathsf Z^\pm$.  The volume-preserving weak change of
variables pulls back the remaining tensor divergence and pressure gradient,
with $(\nabla\delta p)\circ X=\grad_{\gmetric}\delta\mathsf p$, and gives
\begin{equation}\label{eq:reference-difference}
 \mathcal A^\mp\delta\mathsf Z^\pm
 +\Div_{\gmetric}\bigl(
   (\mathsf Z^\pm+\delta\mathsf Z^\pm)
   \otimes\delta\mathsf Z^\mp\bigr)
 +\grad_{\gmetric}\delta\mathsf p=0.
\end{equation}
Here $\Div_{\gmetric}\delta\mathsf Z^\pm=0$ and
$\delta\mathsf Z^\pm(0)=0$.
Lemma~\ref{lem:composition} and \eqref{eq:reference-difference} yield
\[
 \delta\mathsf Z^\pm
 \in C_tH^s\cap C_t^1H^{s-1}\subset C_t^1L^2,
\]
which justifies the distributional identities and Hilbert-space product
rules used below.

\emph{Step 2: transport of the energy density and the interaction flux.}
After an orthogonal change of variables, assume $\Binf=v_{\rm A}e_1$, and put
\[
    e_\pm=\frac12 |\delta\mathsf Z^\pm|^2.
\]
Recall $\Kmat=(dX)^{-1}$.  Taking the scalar product of
\eqref{eq:reference-difference} with
$\delta\mathsf Z^\pm$, using
$\Div_{\gmetric}\delta\mathsf Z^\pm=0$ and writing the metric
divergences in conservative form, gives
\[
 (\partial_t\mp v_{\rm A}\partial_1)e_\pm
 +\partial_a\left(
   \Kmat_i^a\delta\mathsf Z^{\mp,i}e_\pm
   +\Kmat_j^a\delta\mathsf p \delta\mathsf Z^{\pm,j}\right)
 =-E_i\mathsf Z^{\pm,j}
   \delta\mathsf Z^{\mp,i}\delta\mathsf Z^{\pm,j}.
\]
Thus the constant transports send the two line energies in opposite
directions.

Let $t_0$ be a time through which the two solutions agree and put
$I=[t_0,t_1]\subset[0,T_*]$.  Writing $r=y_1$, define
\[
\begin{aligned}
 \rho_\pm(t,r)
 &:=\int_{\Real^{n-1}}e_\pm(t,r,y') dy',\\
 f_\pm(t,r)
 &:=\int_{\Real^{n-1}}
   \left(\Kmat_i^1\delta\mathsf Z^{\mp,i}e_\pm
   +\delta\mathsf p \Kmat_j^1\delta\mathsf Z^{\pm,j}\right)
   (t,r,y') dy',\\
 h_\pm(t,r)
 &:=-\int_{\Real^{n-1}}
   E_i\mathsf Z^{\pm,j}
   \delta\mathsf Z^{\mp,i}\delta\mathsf Z^{\pm,j}
   (t,r,y') dy'.
\end{aligned}
\]
Integration in $y'$ gives the balance laws:
\begin{equation}\label{eq:balancelaw}
    (\partial_t\mp v_{\rm A}\partial_r)\rho_\pm
 +\partial_rf_\pm=h_\pm.
\end{equation}
Introduce the interaction functional and the interaction flux
\[
 \mathscr I(t):=\int_{r_+>r_-}
 \rho_+(t,r_+)\rho_-(t,r_-) dr_+ dr_-,
 \qquad
 \Phi_I:=\left(
 \int_I\int_{\Real}\rho_+(t,r)\rho_-(t,r) dr dt\right)^{1/2}.
\]
The quantity $\Phi_I$ is finite by applying
\eqref{eq:bochner-embedding} below to $\delta\mathsf Z^\pm$.  These definitions were motivated geometrically in \S\ref{ssec:overview}.

For a rigorous derivation of the interaction identity, define the cumulative
density and source
\[
 R_-(t,r):=\int_{-\infty}^r\rho_-(t,q)dq,
 \qquad
 H_-(t,r):=\int_{-\infty}^rh_-(t,q)dq.
\]
Then
\[
 \mathscr I(t)=\int_{\Real}\rho_+(t,r)R_-(t,r)dr.
\]
Indeed, the Hilbert-valued one-dimensional Sobolev embedding and
$\delta\mathsf Z^\pm\in C_tH^s$ give
\[
 \rho_\pm\in C(I;L^1(\Real))\cap L^\infty(I\times\Real),
 \qquad
 \sup_{r\in\Real}\rho_\pm(t,r)
 \lesssim\|\delta\mathsf Z^\pm(t)\|_{H^s}^2.
\]
Moreover, the definitions of $f_\pm,h_\pm$, Sobolev embedding, the uniform
bound for $\Kmat$, the fact that $\delta\mathsf p\in L^\infty_tL^2$, and H\"older's
inequality give
\[
 f_\pm,h_\pm\in L^1(I\times\Real).
\]
Thus all the products below are integrable. Taking the primitive from
$-\infty$ in the minus balance law in \eqref{eq:balancelaw} gives
\[
 \partial_tR_-=-v_{\rm A}\rho_--f_-+H_-
 \quad\text{in }\mathcal D'((t_0,t_1)\times\Real),
\]
while the plus balance law reads
\[
 \partial_t\rho_+
 =v_{\rm A}\partial_r\rho_+-\partial_rf_++h_+.
\]

We next justify the product rule and the integrations by parts. Let
$\eta_\varepsilon$ be a standard nonnegative mollifier in $r$, append a
subscript $\varepsilon$ to denote convolution with $\eta_\varepsilon$, and
take $\chi_R(r)=\chi(r/R)$, where $\chi\in C_c^\infty(\Real)$ equals one on
$[-1,1]$. Set
\[
 \mathscr I_{\varepsilon,R}(t)
 :=\int_{\Real}\chi_R(r)
 \rho_{+,\varepsilon}(t,r)R_{-,\varepsilon}(t,r)dr.
\]
For fixed $\varepsilon>0$, the spatially mollified balance laws give the
time regularity needed for the usual product rule. For almost every
$t\in(t_0,t_1)$, that rule and integration by parts give
\[
\begin{aligned}
 \mathscr I_{\varepsilon,R}'(t)
 &=
 \int_{\Real}\chi_R
 \bigl(
  v_{\rm A}\partial_r\rho_{+,\varepsilon}
  -\partial_rf_{+,\varepsilon}
  +h_{+,\varepsilon}
 \bigr)R_{-,\varepsilon}dr+
 \int_{\Real}\chi_R\rho_{+,\varepsilon}
 \bigl(
  -v_{\rm A}\rho_{-,\varepsilon}
  -f_{-,\varepsilon}
  +H_{-,\varepsilon}
 \bigr)dr\\
 &=
 -2v_{\rm A}\int_{\Real}\chi_R
     \rho_{+,\varepsilon}\rho_{-,\varepsilon}dr+
 \int_{\Real}\chi_R
 \bigl(
  f_{+,\varepsilon}\rho_{-,\varepsilon}
  -\rho_{+,\varepsilon}f_{-,\varepsilon}
 \bigr)dr+
 \int_{\Real}\chi_R
 \bigl(
  h_{+,\varepsilon}R_{-,\varepsilon}
  +\rho_{+,\varepsilon}H_{-,\varepsilon}
 \bigr)dr\\
 &\quad+
 \int_{\Real}\chi_R'
 \bigl(f_{+,\varepsilon}
       -v_{\rm A}\rho_{+,\varepsilon}\bigr)
 R_{-,\varepsilon}dr.
\end{aligned}
\]
Since
\[
 \|R_{-,\varepsilon}(t)\|_{L^\infty_r}
 \le\|\rho_-(t)\|_{L^1_r},
 \qquad
 \|\chi_R'\|_\infty\lesssim R^{-1},
\]
by dominated convergence, the last line tends to zero in $L^1(I)$ as $R\to\infty$. Letting first
$R\to\infty$ and then $\varepsilon\to0$ is justified by the stated
$L^1$ and $L^\infty$ bounds together with
\[
 \|R_{-,\varepsilon}(t)-R_-(t)\|_{L^\infty_r}
 \le\|\rho_{-,\varepsilon}(t)-\rho_-(t)\|_{L^1_r},
 \qquad
 \|H_{-,\varepsilon}(t)-H_-(t)\|_{L^\infty_r}
 \le\|h_{-,\varepsilon}(t)-h_-(t)\|_{L^1_r}.
\]
The same bounds give
$\mathscr I_{\varepsilon,R}\to\mathscr I$ in $L^1(I)$. Consequently,
$\mathscr I\in W^{1,1}(I)$ and, for almost every $t\in I$,
\[
\begin{aligned}
 \mathscr I'(t)
 &=
 -2v_{\rm A}\int_{\Real}\rho_+(t,r)\rho_-(t,r)dr\\
 &\quad+
 \int_{\Real}\bigl(f_+(t,r)\rho_-(t,r)
                  -\rho_+(t,r)f_-(t,r)\bigr)dr\\
 &\quad+
 \int_{\Real}h_+(t,r)R_-(t,r)dr
 +\int_{\Real}\rho_+(t,r)H_-(t,r)dr.
\end{aligned}
\]
Expanding the two cumulative functions, this is equivalently
\begin{subequations}\label{eq:interaction-flux-identity}
\begin{align}
 \mathscr I'(t)+2v_{\rm A}\int_{\Real}\rho_+\rho_- dr
 =&\int_{\Real}\bigl(f_+\rho_--\rho_+f_-\bigr) dr \label{eq:longitudinalflux}\\
 &+\int_{\Real}h_+(r_+)
       \int_{r_-<r_+}\rho_-(r_-) dr_- dr_+\label{eq:sourceterm1}\\
 &+\int_{\Real}\rho_+(r_+)
       \int_{r_-<r_+}h_-(r_-) dr_- dr_+ \label{eq:sourceterm2}.
\end{align}
\end{subequations}
At $t=t_0$, $\mathscr I(t_0)=0$.  Integrating
\eqref{eq:interaction-flux-identity}, discarding the nonnegative value
$\mathscr I(t_1)$, and estimating the time-integrated right-hand side will
allow us to conclude $\Phi_I=0$.  In the next step, we show that
we can control the energy by this auxiliary quantity and then provide the
missing bounds in the following ones.

\emph{Step 3: control of the energy by the flux.}
Define the averaged decaying Elsasser fields
\[
 z^\pm:=\frac12\bigl(Z^\pm+\widetilde Z^\pm\bigr)\mp\Binf
\]
and, for $|I|\le1$, set
\[
 c_I:=
 \sup_{t\in I}\bigl(
  \|\delta Z^+(t)\|_{H^s}+\|\delta Z^-(t)\|_{H^s}\bigr)+\left(\int_I
  \bigl(\|z^+(t)\|_{H^s}^2+\|z^-(t)\|_{H^s}^2\bigr) dt
  \right)^{1/2}.
\]
Since the difference vanishes at $t_0$ and belongs to $C_tH^s$, one has
$c_I\to0$ as $t_1\downarrow t_0$.  The identities
\[
 Z^\pm\mp\Binf=z^\pm-\frac12\delta Z^\pm,
 \qquad
 E_i\mathsf Z^{\pm,j}=(\partial_iZ^{\pm,j})\circ X
\]
Lemma~\ref{lem:composition} and the uniform bound on $\Kmat$ imply
\begin{equation}\label{eq:uniqueness-small-coefficients}
 \sup_{t\in I}\bigl(
  \|\delta\mathsf Z^+(t)\|_\infty+
  \|\delta\mathsf Z^-(t)\|_\infty\bigr)+\sum_{\pm}
 \left\|\bigl(E_i\mathsf Z^{\pm,j}\bigr)_{i,j}\right\|_
 {L^2(I;L^2_{y_1}H^{s-1}_{y'})}
 \lesssim c_I .
\end{equation}
For the second term, we also used $|I|\le1$ to bound the contribution of
$\delta Z^\pm$ in $L^2_tH^s$ by its time supremum.
All implicit constants in the remainder of the proof depend only on the
standing parameters and on the uniform bounds for $X$ and $\Kmat$ on
$[0,T_*]$; in particular, they are independent of the interval $I$.

Write the time-- localised energy
\[
 \mathcal E_I
 :=\sup_{t\in I}\int_{\Rn}(e_++e_-)(t,y) dy.
\]
Taking the $L^2$ inner products of \eqref{eq:reference-difference} with
$\delta\mathsf Z^\pm$, integrating from $t_0$ to $t$, and using the same
divergence cancellations as above, transverse Sobolev embedding and
Cauchy--Schwarz in $(t,r)$ gives
\[
\begin{aligned}
\mathcal E_I
&\leq
\sup_{t\in I}\sum_{\pm}
\left|
\int_{t_0}^{t}\int_{\Rn}
\Bigl[
\delta\mathsf Z^{\pm,j}\delta\mathsf Z^{\mp,i}
E_i\bigl(\mathsf Z^{\pm,j}+\delta\mathsf Z^{\pm,j}\bigr)
+\delta\mathsf Z^{\pm,j}E_j\delta\mathsf p
\Bigr] dy d\tau
\right|                                                        \\
&=
\sup_{t\in I}\sum_{\pm}
\left|
\int_{t_0}^{t}\int_{\Rn}
\Bigl[
\delta\mathsf Z^{\pm,j}\delta\mathsf Z^{\mp,i}
E_i\mathsf Z^{\pm,j}
+\Div_{\gmetric}\bigl(e_\pm\delta\mathsf Z^\mp\bigr)
+\Div_{\gmetric}\bigl(\delta\mathsf p \delta\mathsf Z^\pm\bigr)
\Bigr] dy d\tau
\right|                                                        \\
&=
\sup_{t\in I}\sum_{\pm}
\left|
\int_{t_0}^{t}\int_{\Rn}
\delta\mathsf Z^{\pm,j}\delta\mathsf Z^{\mp,i}
E_i\mathsf Z^{\pm,j} dy d\tau
\right|                                                        \\
&\leq
\int_{I\times\Real}
\|\delta\mathsf Z^+(t,r,\cdot)\|_{L^2_{y'}}
\|\delta\mathsf Z^-(t,r,\cdot)\|_{L^2_{y'}}
\sum_{i,j}\Bigl(
\|E_i\mathsf Z^{+,j}(t,r,\cdot)\|_{L^\infty_{y'}}
+\|E_i\mathsf Z^{-,j}(t,r,\cdot)\|_{L^\infty_{y'}}
\Bigr) dr dt                                                \\
&=
2\int_{I\times\Real}
\bigl(\rho_+(t,r)\rho_-(t,r)\bigr)^{1/2}
\sum_{i,j}\Bigl(
\|E_i\mathsf Z^{+,j}(t,r,\cdot)\|_{L^\infty_{y'}}
+\|E_i\mathsf Z^{-,j}(t,r,\cdot)\|_{L^\infty_{y'}}
\Bigr) dr dt                                                \\
&\lesssim
\Phi_I
\sum_{\pm}
\left\|
\bigl(E_i\mathsf Z^{\pm,j}\bigr)_{i,j}
\right\|_{L^2(I;L^2_{y_1}H^{s-1}_{y'})}
\lesssim c_I\Phi_I .
\end{aligned}
\]
Here we used $s-1>(n-1)/2$.  Thus
\begin{equation}\label{eq:energy-by-interaction}
 \mathcal E_I\lesssim c_I\Phi_I
\end{equation}
as desired.

\emph{Step 4: control of the source terms
\eqref{eq:sourceterm1}--\eqref{eq:sourceterm2}.}
By the definition of $h_\pm$ and the estimate above,
\[
\begin{aligned}
 \int_{I\times\Real}(|h_+|+|h_-|) dr dt
 &\le
 \int_{I\times\Rn}
 |\delta\mathsf Z^+||\delta\mathsf Z^-|
 \sum_{i,j}\bigl(
 |E_i\mathsf Z^{+,j}|+|E_i\mathsf Z^{-,j}|
 \bigr) dy dt\\
 &\lesssim c_I\Phi_I.
\end{aligned}
\]
Moreover, for every $t\in I$,
\[
 \sup_{r_+\in\Real}\int_{r_-<r_+}\rho_-(t,r_-) dr_-
 \le\mathcal E_I,
 \qquad
 \sup_{r_-\in\Real}\int_{r_+>r_-}\rho_+(t,r_+) dr_+
 \le\mathcal E_I.
\]
Consequently, after applying Fubini to the second source term
\eqref{eq:sourceterm2}, we obtain
\begin{equation}\label{eq:sourcetermbound}
\begin{split}
    &\int_I
 \left|
 \int_{\Real}h_+(t,r_+)
     \int_{r_-<r_+}\rho_-(t,r_-) dr_- dr_+
 \right| dt+
 \int_I
 \left|
 \int_{\Real}\rho_+(t,r_+)
     \int_{r_-<r_+}h_-(t,r_-) dr_- dr_+
 \right| dt\\
 &\le
 \mathcal E_I
 \int_{I\times\Real}(|h_+|+|h_-|) dr dt\\
 &\le C\mathcal E_Ic_I\Phi_I
 \lesssim c_I^2\Phi_I^2,
 \end{split}
\end{equation}
where the last inequality follows from
\eqref{eq:energy-by-interaction}.

\emph{Step 5: control of the full longitudinal flux
\eqref{eq:longitudinalflux}.}
The fluxes $f_\pm$ contain $\delta\mathsf p$ itself, so we
estimate $\delta p$ first rather than projecting the equation onto
divergence-free fields in Lagrangian coordinates.

The constant-field terms disappear under the two Riesz transforms by
incompressibility.  Pressure recovery and its divergence form give the two
representations
\[
\begin{aligned}
 \delta p
 &=\Riesz_i\Riesz_j\bigl(
   \delta Z^{+,i}z^{-,j}+z^{+,i}\delta Z^{-,j}\bigr),\\
 \nabla\delta p
 &=\nabla(-\Delta_x)^{-1}\Div_x\bigl(
   (\delta Z^+\cdot\nabla)z^-
   +(\delta Z^-\cdot\nabla)z^+\bigr).
\end{aligned}
\]
Choose
\[
 \frac12<\alpha<s-\frac n2 .
\]
Since $s-\alpha>n/2$, duality and the usual Sobolev product estimate
give
\[
 \|ag\|_{H^{\alpha-1}}
 \le C\|a\|_{H^{s-1}}\|g\|_2.
\]
The standard inhomogeneous elliptic estimate
\[
 \|q\|_{H^\alpha}
 \le C\bigl(\|q\|_2+\|\nabla q\|_{H^{\alpha-1}}\bigr)
\]
follows directly from Plancherel and
$\langle\xi\rangle^{2\alpha}\lesssim
1+|\xi|^2\langle\xi\rangle^{2(\alpha-1)}$.  The two representations above
yield separately
\[
\begin{aligned}
 \|\delta p\|_2
 &\le C\|\delta Z^+\otimes z^-+z^+\otimes\delta Z^-\|_2,\\
 \|\nabla\delta p\|_{H^{\alpha-1}}
 &\le C\|(\delta Z^+\cdot\nabla)z^-
          +(\delta Z^-\cdot\nabla)z^+\|_{H^{\alpha-1}}.
\end{aligned}
\]
Combining these two bounds with the inhomogeneous elliptic estimate,
$H^s\hookrightarrow L^\infty$, and the product estimate yields
\[
 \|\delta p(t)\|_{H_x^\alpha}
 \le C\bigl(\|z^+(t)\|_{H^s}+\|z^-(t)\|_{H^s}\bigr)
 \bigl(\|\delta Z^+(t)\|_2+\|\delta Z^-(t)\|_2\bigr).
\]

Since $\alpha>1/2$, we have, with equivalence of norms,
\begin{equation}\label{eq:bochner-embedding}
H^\alpha(\mathbb{R}^n)
=H^\alpha\bigl(\mathbb{R}_{y_1};L^2(\mathbb{R}^{n-1})\bigr)
 \cap L^2\bigl(\mathbb{R}_{y_1};H^\alpha(\mathbb{R}^{n-1})\bigr)
\hookrightarrow L^\infty_{y_1}L^2_{y'}.
\end{equation}
Now Lemma~\ref{lem:composition}, volume preservation, and
the Hilbert-valued one-dimensional Sobolev embedding imply
\[
 \sup_{r\in\Real}
 \|\delta\mathsf p(t,r,\cdot)\|_{L^2_{y'}}
 \le C\bigl(\|z^+(t)\|_{H^s}+\|z^-(t)\|_{H^s}\bigr)
 \left(\int_{\Rn}(e_++e_-)(t,y) dy\right)^{1/2}.
\]
Write the pressure part of $f_\pm$ as
\[
 P_\pm(t,r):=
 \int_{\Real^{n-1}}
 \delta\mathsf p\Kmat_j^1\delta\mathsf Z^{\pm,j}
 (t,r,y')dy',
 \qquad
 A(t):=\|z^+(t)\|_{H^s}+\|z^-(t)\|_{H^s}.
\]
The inverse deformation $\Kmat$ is uniformly bounded, and hence
\[
 |P_\pm(t,r)|
 \le CA(t)
 \mathcal E_I^{1/2}\rho_\pm(t,r)^{1/2}.
\]
Here $\|\delta\mathsf Z^\pm(t,r,\cdot)\|_{L^2_{y'}}
 =\left(\int_{\Real^{n-1}}
 |\delta\mathsf Z^\pm(t,r,y')|^2dy'\right)^{1/2}
 =\bigl(2\rho_\pm(t,r)\bigr)^{1/2}$.  Since
$\int_{\Real}\rho_\pm(t,r)dr\le\mathcal E_I$, Cauchy--Schwarz in $r$
gives, for each $t\in I$,
\[
\begin{aligned}
 \int_{\Real}\rho_-(t,r)|P_+(t,r)|dr
 &\le CA(t)\mathcal E_I^{1/2}
 \int_{\Real}\rho_-(t,r)\rho_+(t,r)^{1/2}dr\\
 &\le CA(t)\mathcal E_I^{1/2}
 \left(\int_{\Real}\rho_+(t,r)\rho_-(t,r)dr\right)^{1/2}
 \left(\int_{\Real}\rho_-(t,r)dr\right)^{1/2}\\
 &\le CA(t)\mathcal E_I
 \left(\int_{\Real}\rho_+(t,r)\rho_-(t,r)dr\right)^{1/2}.
\end{aligned}
\]
The same estimate holds with the signs interchanged.  Therefore,
Cauchy--Schwarz in $t$, the definition of $c_I$, and
\eqref{eq:energy-by-interaction} show that the time-integrated pressure
contribution to \eqref{eq:longitudinalflux} satisfies
\begin{equation}\label{eq:pressurefluxbound}
\begin{aligned}
 &\int_I\int_{\Real}
 \bigl(\rho_-(t,r)|P_+(t,r)|
       +\rho_+(t,r)|P_-(t,r)|\bigr)drdt\\
 &\quad\le C\mathcal E_I
 \|A\|_{L^2(I)}
 \left(\int_I\int_{\Real}\rho_+(t,r)\rho_-(t,r)drdt\right)^{1/2}\\
 &\quad\le Cc_I\mathcal E_I\Phi_I
 \lesssim c_I^2\Phi_I^2.
\end{aligned}
\end{equation}

It remains to control the advective parts of $f_\pm$.  The uniform bound
for $\Kmat$ and \eqref{eq:uniqueness-small-coefficients} give, pointwise
in $(t,r)$,
\[
\begin{aligned}
 \left|
 \int_{\Real^{n-1}}
 \Kmat_i^1\delta\mathsf Z^{\mp,i}e_\pm(t,r,y') dy'
 \right|
 &\le
 \|\Kmat\|_{L^\infty}
 \|\delta\mathsf Z^\mp(t)\|_{L^\infty}
 \int_{\Real^{n-1}}e_\pm(t,r,y') dy'\\
 &\le Cc_I\rho_\pm(t,r).
\end{aligned}
\]
Therefore their contribution to
$f_+\rho_--\rho_+f_-$ satisfies
\begin{equation}\label{eq:advectivefluxbound}
\begin{split}
 &\int_I\int_{\Real}
 \rho_-(t,r)
 \left|
 \int_{\Real^{n-1}}
 \Kmat_i^1\delta\mathsf Z^{-,i}e_+(t,r,y') dy'
 \right| dr dt\\
 &\quad+
 \int_I\int_{\Real}
 \rho_+(t,r)
 \left|
 \int_{\Real^{n-1}}
 \Kmat_i^1\delta\mathsf Z^{+,i}e_-(t,r,y') dy'
 \right| dr dt\\
 &\le
 Cc_I\int_I\int_{\Real}\rho_+(t,r)\rho_-(t,r) dr dt
 \lesssim c_I\Phi_I^2.   
\end{split}
\end{equation}

\emph{Step 6: conclusion.}  Integrating
\eqref{eq:interaction-flux-identity} over $I=[t_0,t_1]$ and using
$\mathscr I(t_0)=0$ gives
\[
\begin{aligned}
 \mathscr I(t_1)+2v_{\rm A}\Phi_I^2
 &=\int_I\int_{\Real}
   \bigl(f_+\rho_--\rho_+f_-\bigr)drdt\\
 &\quad+
 \int_I\int_{\Real}h_+(t,r_+)
       \int_{r_-<r_+}\rho_-(t,r_-)dr_-dr_+dt\\
 &\quad+
 \int_I\int_{\Real}\rho_+(t,r_+)
       \int_{r_-<r_+}h_-(t,r_-)dr_-dr_+dt.
\end{aligned}
\]
We now bound each part of this time-integrated right-hand side.  Splitting
$f_\pm$ into the advective and pressure terms in its definition,
\eqref{eq:advectivefluxbound} and \eqref{eq:pressurefluxbound} give
\[
 \int_I\left|
 \int_{\Real}\bigl(f_+\rho_--\rho_+f_-\bigr)dr
 \right|dt
 \le Cc_I\Phi_I^2+Cc_I^2\Phi_I^2.
\]
The two remaining integrals are precisely the source contributions bounded
in \eqref{eq:sourcetermbound}; together they are at most
$Cc_I^2\Phi_I^2$.  Since $\mathscr I(t_1)\ge0$, the integrated identity
therefore yields
\[
 2v_{\rm A}\Phi_I^2
 \le
 \bigl(Cc_I\Phi_I^2+Cc_I^2\Phi_I^2\bigr)
 +Cc_I^2\Phi_I^2
 \le C(c_I+c_I^2)\Phi_I^2.
\]
Since $c_I\to0$ as $t_1\downarrow t_0$, we may choose $I$ so that
$C(c_I+c_I^2)<2v_{\rm A}$.  The preceding inequality then forces
$\Phi_I=0$, and \eqref{eq:energy-by-interaction} gives
$\mathcal E_I=0$.  Thus the two solutions agree on $I$.

Finally, define
\[
 \mathcal S:=\left\{t\in[0,T_*]:
 \widetilde v=v\ \text{and}\ \widetilde B=B\ \text{on }[0,t]\right\}.
\]
The set $\mathcal S$ is nonempty and closed by continuity.  If
$\tau:=\sup\mathcal S<T_*$, then $\tau\in\mathcal S$ and restarting the
same short-interval argument at $t_0=\tau$ extends the agreement beyond
$\tau$, a contradiction.  Thus the solutions agree on $[0,T_*]$.  Their
normalised pressures agree by Lemma~\ref{lem:pressure-recovery}. 
\end{proof}


\bibliographystyle{amsplain}
\bibliography{mhd-constant-field}

\end{document}
